\documentclass{article}
\usepackage{eso-pic}
\usepackage{datetime}
\usepackage{lastpage}           %-to identify last page
\usepackage{xspace}             %-to force add a space
\usepackage{tabularx}
\usepackage{tikz}               %-to create Pfaudler logo
	\usetikzlibrary{shapes.geometric} %-for creating hexagon
	\usetikzlibrary{calc}
\usepackage{pgfplots} 
\pgfplotsset %-Global settings for pgfplots
  {
    compat = newest,
    xtick style = {draw=none},
    width = .95 \textwidth,
    height = .81 \textheight
  }

\usepackage[linktoc=all,        %-for clickable links and pdf metadata
            colorlinks=true,
            linkcolor=black,
            urlcolor=black,
            citecolor=black,
            pdfauthor={Pfaudler},
            pdftitle={Quality Review},
            pdfsubject={},
            pdfkeywords={},
            pdfproducer={},
            pdfcreator={}]{hyperref}

\usepackage{bookmark}

\newdateformat{monthyeardate}{%
  \monthname[\THEMONTH], \THEYEAR}
\newcommand*{\fyyear}{\Roman{year}}



%-----Font configuration-----%
\usepackage{fontspec}
\defaultfontfeatures{Ligatures={NoCommon,TeX}}
\setmainfont[
  Path          = ./,
  UprightFont   = font/OTF/Pfaudler-Book.otf,
  BoldItalicFont= font/OTF/Pfaudler-BoldItalic.otf,
  BoldFont      = font/OTF/Pfaudler-Bold.otf,
  ItalicFont    = font/OTF/Pfaudler-BookItalic.otf
]{Pfaudler}
%----------------------------%



%-----Margin formatting-----%

%\usepackage[paperwidth=8.27in, paperheight=11.69in, landscape]{geometry}
%\usepackage[paperwidth=7.5in, paperheight=13.33in, landscape]{geometry}
\input{layout}%  ^one of these written to layout.tex by the Python script
\geometry{top=0.9in, bottom=0.9in, left=0.9in, right=0.9in}
%---------------------------%




%-----Pfaudler colorscheme-----%
\definecolor{pfblue}{RGB}{52,134,199}
\definecolor{pfgrey}{RGB}{125,125,130}
%------------------------------%



%--------------Pfaudler logo (invoke with \pflogo)---------------%
%---(Must be contained in one paragraph for use in the header)---%
\newcommand\pflogo{%
\raisebox{-9pt}{%
\begin{tikzpicture}
\node [right,scale=1,inner sep=1pt,outer sep=0pt, text=pfblue]
at (0,0) {\fontsize{10.73}{5}\selectfont{%
%
%--Strapline--%
Defining the Standard}};
%
\node [right,scale=1,inner sep=0pt,outer sep=0pt, text=pfgrey]
at (0.01,0.70) {\addfontfeature{LetterSpace=0.0}\fontsize{30}{0}\selectfont{%
%
%--Manually kerned Main text--%
\textbf{%
P{\kern 0.0pt}%
f{\kern-1.0pt}%
a{\kern-1.1pt}%
u{\kern-1.1pt}%
d{\kern-1.1pt}%
l{\kern-1.1pt}%
e{\kern-1.1pt}%
r}%
}};
%
%--Pf hexagon--%
\path (4.70,1.32) node [regular polygon,
		regular polygon sides=6,
		rounded corners=1.5pt,
		inner sep=8.8pt,
		fill = pfblue](hexagon){};
%--Pf hexagon text--%
\node [text=white, scale=1.1] at (4.70,1.32) {\large{\textbf{Pf}}};
\end{tikzpicture}}}
%----------------------------------------------------------------%
\newcommand\pflogowhite{%
\raisebox{-9pt}{%
\begin{tikzpicture}
\node [right,scale=1,inner sep=1pt,outer sep=0pt, text=white]
at (0,0) {\fontsize{10.73}{5}\selectfont{%
%
%--Strapline--%
Defining the Standard}};
%
\node [right,scale=1,inner sep=0pt,outer sep=0pt, text=white]
at (0.01,0.70) {\addfontfeature{LetterSpace=0.0}\fontsize{30}{0}\selectfont{%
%
%--Manually kerned Main text--%
\textbf{%
P{\kern 0.0pt}%
f{\kern-1.0pt}%
a{\kern-1.1pt}%
u{\kern-1.1pt}%
d{\kern-1.1pt}%
l{\kern-1.1pt}%
e{\kern-1.1pt}%
r}%
}};
%
%--Pf hexagon--%
\path (4.70,1.32) node [regular polygon,
		regular polygon sides=6,
		rounded corners=1.5pt,
		inner sep=8.8pt,
		fill = white](hexagon){};
%--Pf hexagon text--%
\node [text=pfblue, scale=1.1] at (4.70,1.32) {\large{\textbf{Pf}}};
\end{tikzpicture}}}
%----------------------------------------------------------------%




%---------------------------------Header Formatting---------------------------------%
\usepackage{fancyhdr}
\pagestyle{fancy}

%--First page header style--%
\fancypagestyle{style1}{
\fancyhf{}
\fancyhead[L]{\pflogowhite}
\fancyfoot[R]{\raisebox{-2.5em}{
\definecolor{c3486c7}{RGB}{52,134,199}


\begin{tikzpicture}[y=0.80pt, x=0.80pt, yscale=-1.000000, xscale=1.000000, inner sep=0pt, outer sep=0pt]
\path[fill=c3486c7,line width=0.153pt] (127.8015,27.8018) -- (125.1667,27.8018)
  -- (125.1667,33.4979) -- (120.2802,33.4979) -- (120.2802,13.8363) --
  (128.8677,13.8363) .. controls (131.3350,13.8363) and (133.3108,14.4627) ..
  (134.7951,15.7157) .. controls (136.2794,16.9686) and (137.0215,18.6457) ..
  (137.0215,20.7468) .. controls (137.0215,22.0961) and (136.6938,23.3008) ..
  (136.0384,24.3610) .. controls (135.3830,25.4212) and (134.4578,26.2501) ..
  (133.2626,26.8477) -- (137.3918,33.4979) -- (131.5223,33.4979) --
  cycle(125.1667,23.4069) -- (128.7809,23.4069) .. controls (130.1652,23.3346)
  and (131.2985,22.0074) .. (131.3180,20.7468) .. controls (131.2025,19.1994)
  and (129.7217,18.1538) .. (128.7809,18.1445) -- (125.1667,18.1445) --
  cycle(112.1053,18.0867) -- (112.1053,21.6720) -- (117.6624,21.6720) --
  (117.6624,25.8357) -- (112.1053,25.8357) -- (112.1053,29.2475) --
  (122.0807,29.2475) -- (122.0807,33.4979) -- (107.2188,33.4979) --
  (107.2188,13.8363) -- (118.9057,13.8363) -- (118.9057,18.0867) --
  cycle(110.6994,29.1897) -- (110.6994,33.4979) -- (96.0977,33.4979) --
  (96.0977,13.8363) -- (100.9842,13.8363) -- (100.9842,29.1897) --
  cycle(76.3949,13.8363) -- (84.2884,13.8363) .. controls (87.3726,13.8363) and
  (89.9171,14.7519) .. (91.9218,16.5831) .. controls (93.9072,18.3951) and
  (94.9000,20.7371) .. (94.9000,23.6093) .. controls (94.9000,26.5200) and
  (93.9072,28.9006) .. (91.9218,30.7511) .. controls (89.9171,32.5823) and
  (87.3726,33.4979) .. (84.2884,33.4979) -- (76.3949,33.4979) --
  cycle(81.2814,29.1030) -- (84.2884,29.1030) .. controls (85.8884,29.1030) and
  (87.2281,28.5825) .. (88.3075,27.5416) .. controls (89.3870,26.4621) and
  (89.9267,25.1514) .. (89.9267,23.6093) .. controls (89.9267,22.0865) and
  (89.3870,20.8142) .. (88.3075,19.7926) .. controls (87.2473,18.7517) and
  (85.9076,18.2312) .. (84.2884,18.2312) -- (81.2814,18.2312) --
  cycle(74.8185,13.8363) -- (74.8185,24.3321) .. controls (74.8185,27.3585) and
  (73.9993,29.7198) .. (72.3608,31.4161) .. controls (70.7031,33.0931) and
  (68.4670,33.9316) .. (65.6527,33.9316) .. controls (62.7998,33.9316) and
  (60.5638,33.0931) .. (58.9446,31.4161) .. controls (57.3254,29.7391) and
  (56.5158,27.3778) .. (56.5158,24.3321) -- (56.5158,13.8363) --
  (61.4312,13.8363) -- (61.4312,24.4189) .. controls (61.4312,25.9224) and
  (61.8071,27.1272) .. (62.5589,28.0331) .. controls (63.3299,28.9198) and
  (64.3612,29.3632) .. (65.6527,29.3632) .. controls (66.9442,29.3632) and
  (67.9755,28.9198) .. (68.7465,28.0331) .. controls (69.5176,27.1272) and
  (69.9031,25.9224) .. (69.9031,24.4189) -- (69.9031,13.8363) --
  cycle(53.6176,30.1150) -- (45.2903,30.1150) -- (43.9891,33.4979) --
  (38.7846,33.4979) -- (47.1408,13.7205) -- (51.8249,13.7206) --
  (60.1522,33.4979) -- (54.9476,33.4979) -- cycle(51.9984,25.9513) --
  (49.4539,19.4167) -- (46.9095,25.9513) -- cycle(36.2805,18.0261) --
  (36.2805,22.1898) -- (40.4902,22.1898) -- (40.4902,26.4402) --
  (36.2805,26.4402) -- (36.2805,33.3796) -- (31.3940,33.3796) --
  (31.3940,13.7205) -- (47.1408,13.7205) -- (47.1408,18.0261) --
  cycle(13.6824,13.8363) -- (21.6049,13.8363) .. controls (24.1879,13.8363) and
  (26.2119,14.4820) .. (27.6769,15.7735) .. controls (29.1804,17.0650) and
  (29.9322,18.8191) .. (29.9322,21.0359) .. controls (29.9322,23.3298) and
  (29.1515,25.1514) .. (27.5901,26.5007) .. controls (26.0288,27.8307) and
  (23.9084,28.4958) .. (21.2290,28.4958) -- (18.5689,28.4958) --
  (18.5689,33.4979) -- (13.6824,33.4979) -- cycle(18.5689,24.2165) --
  (21.5182,24.2165) .. controls (22.5976,24.2165) and (23.4554,23.9273) ..
  (24.0915,23.3490) .. controls (24.7469,22.7708) and (25.0746,21.9997) ..
  (25.0746,21.0359) .. controls (25.0746,20.1492) and (24.7469,19.4456) ..
  (24.0915,18.9252) .. controls (23.4554,18.4047) and (22.5976,18.1445) ..
  (21.5182,18.1445) -- (18.5689,18.1445) -- cycle;\path[fill=black,line width=0.048pt] (20.1299,41.7333) -- (20.1299,44.8963) --
  (19.4565,44.8963) -- (19.3945,43.7334) .. controls (18.9338,44.5955) and
  (18.1925,45.0265) .. (17.1706,45.0265) .. controls (16.2787,45.0265) and
  (15.5374,44.7102) .. (14.9467,44.0776) .. controls (14.3501,43.4388) and
  (14.0518,42.6760) .. (14.0518,41.7891) .. controls (14.0518,40.8278) and
  (14.3560,40.0309) .. (14.9644,39.3983) .. controls (15.5728,38.7595) and
  (16.3230,38.4401) .. (17.2149,38.4401) .. controls (17.8883,38.4401) and
  (18.4967,38.6416) .. (19.0401,39.0448) .. controls (19.5718,39.4417) and
  (19.9232,39.9937) .. (20.0945,40.7007) -- (19.1908,40.7007) .. controls
  (19.0431,40.2355) and (18.7861,39.8851) .. (18.4199,39.6495) .. controls
  (18.0537,39.4138) and (17.6550,39.2959) .. (17.2238,39.2959) .. controls
  (16.5622,39.2959) and (16.0099,39.5285) .. (15.5669,39.9937) .. controls
  (15.1239,40.4588) and (14.9024,41.0573) .. (14.9024,41.7891) .. controls
  (14.9024,42.4465) and (15.1151,43.0078) .. (15.5403,43.4730) .. controls
  (15.9715,43.9381) and (16.5150,44.1707) .. (17.1706,44.1707) .. controls
  (17.7141,44.1707) and (18.1659,44.0063) .. (18.5262,43.6776) .. controls
  (18.8866,43.3427) and (19.1199,42.8931) .. (19.2262,42.3287) --
  (17.1086,42.3287) -- (17.1086,41.5100) -- (20.1299,41.5100) -- cycle;\path[fill=black,line width=0.048pt] (25.8949,44.0963) -- (25.8949,44.8963) --
  (21.7837,44.8963) -- (21.7837,38.5703) -- (22.6255,38.5703) --
  (22.6255,44.0963) -- cycle;\path[fill=black,line width=0.048pt] (30.9640,43.2683) -- (27.9516,43.2683) --
  (27.2782,44.8963) -- (26.3744,44.8963) -- (29.0591,38.5331) --
  (29.8565,38.5331) -- (32.5412,44.8963) -- (31.6374,44.8963) --
  cycle(30.6273,42.4496) -- (29.4578,39.6029) -- (28.2882,42.4496) -- cycle;\path[fill=black,line width=0.048pt] (35.4886,38.4308) .. controls
  (36.1147,38.4308) and (36.6522,38.5920) .. (37.1011,38.9145) .. controls
  (37.5559,39.2308) and (37.8188,39.7146) .. (37.8897,40.3658) --
  (37.0480,40.3658) .. controls (36.9180,39.5843) and (36.4159,39.1936) ..
  (35.5417,39.1936) .. controls (35.1342,39.1936) and (34.8063,39.2897) ..
  (34.5582,39.4820) .. controls (34.3102,39.6743) and (34.1861,39.9316) ..
  (34.1861,40.2541) .. controls (34.1861,40.7379) and (34.5553,41.0604) ..
  (35.2936,41.2217) -- (36.1354,41.3984) .. controls (36.7201,41.5224) and
  (37.1543,41.7209) .. (37.4378,41.9938) .. controls (37.7272,42.2605) and
  (37.8720,42.6233) .. (37.8720,43.0822) .. controls (37.8720,43.6590) and
  (37.6623,44.1304) .. (37.2429,44.4963) .. controls (36.8294,44.8498) and
  (36.2771,45.0265) .. (35.5860,45.0265) .. controls (34.8536,45.0265) and
  (34.2599,44.8343) .. (33.8051,44.4498) .. controls (33.3503,44.0652) and
  (33.0933,43.5319) .. (33.0343,42.8497) -- (33.8760,42.8497) .. controls
  (33.9410,43.2962) and (34.1241,43.6404) .. (34.4253,43.8823) .. controls
  (34.7325,44.1180) and (35.1342,44.2358) .. (35.6303,44.2358) .. controls
  (36.0556,44.2358) and (36.3893,44.1428) .. (36.6315,43.9567) .. controls
  (36.8796,43.7645) and (37.0037,43.5102) .. (37.0037,43.1939) .. controls
  (37.0037,42.7039) and (36.6670,42.3907) .. (35.9936,42.2543) --
  (34.9747,42.0496) .. controls (33.8524,41.8263) and (33.2912,41.2496) ..
  (33.2912,40.3193) .. controls (33.2912,39.7549) and (33.4921,39.2990) ..
  (33.8937,38.9517) .. controls (34.2954,38.6044) and (34.8270,38.4308) ..
  (35.4886,38.4308) -- cycle;\path[fill=black,line width=0.048pt] (41.1992,38.4308) .. controls
  (41.8254,38.4308) and (42.3629,38.5920) .. (42.8118,38.9145) .. controls
  (43.2666,39.2308) and (43.5295,39.7146) .. (43.6004,40.3658) --
  (42.7586,40.3658) .. controls (42.6287,39.5843) and (42.1266,39.1936) ..
  (41.2524,39.1936) .. controls (40.8448,39.1936) and (40.5170,39.2897) ..
  (40.2689,39.4820) .. controls (40.0208,39.6743) and (39.8968,39.9316) ..
  (39.8968,40.2541) .. controls (39.8968,40.7379) and (40.2660,41.0604) ..
  (41.0043,41.2217) -- (41.8460,41.3984) .. controls (42.4308,41.5224) and
  (42.8650,41.7209) .. (43.1485,41.9938) .. controls (43.4379,42.2605) and
  (43.5826,42.6233) .. (43.5826,43.0822) .. controls (43.5826,43.6590) and
  (43.3730,44.1304) .. (42.9536,44.4963) .. controls (42.5401,44.8498) and
  (41.9878,45.0265) .. (41.2967,45.0265) .. controls (40.5643,45.0265) and
  (39.9706,44.8343) .. (39.5158,44.4498) .. controls (39.0610,44.0652) and
  (38.8040,43.5319) .. (38.7450,42.8497) -- (39.5867,42.8497) .. controls
  (39.6517,43.2962) and (39.8348,43.6404) .. (40.1360,43.8823) .. controls
  (40.4432,44.1180) and (40.8448,44.2358) .. (41.3410,44.2358) .. controls
  (41.7663,44.2358) and (42.1000,44.1428) .. (42.3422,43.9567) .. controls
  (42.5903,43.7645) and (42.7143,43.5102) .. (42.7143,43.1939) .. controls
  (42.7143,42.7039) and (42.3777,42.3907) .. (41.7043,42.2543) --
  (40.6854,42.0496) .. controls (39.5631,41.8263) and (39.0019,41.2496) ..
  (39.0019,40.3193) .. controls (39.0019,39.7549) and (39.2028,39.2990) ..
  (39.6044,38.9517) .. controls (40.0061,38.6044) and (40.5377,38.4308) ..
  (41.1992,38.4308) -- cycle;\path[fill=black,line width=0.048pt] (47.2977,41.9845) -- (47.2977,42.8404) --
  (44.3207,42.8404) -- (44.3207,41.9845) -- cycle;\path[fill=black,line width=0.048pt] (53.1596,44.0963) -- (53.1596,44.8963) --
  (49.0484,44.8963) -- (49.0484,38.5703) -- (49.8902,38.5703) --
  (49.8902,44.0963) -- cycle;\path[fill=black,line width=0.048pt] (55.0125,38.5703) -- (55.0125,44.8963) --
  (54.1708,44.8963) -- (54.1708,38.5703) -- cycle;\path[fill=black,line width=0.048pt] (61.7602,38.5703) -- (61.7602,44.8963) --
  (61.0691,44.8963) -- (57.3833,40.2076) -- (57.3833,44.8963) --
  (56.5416,44.8963) -- (56.5416,38.5703) -- (57.2061,38.5703) --
  (60.9096,43.3334) -- (60.9096,38.5703) -- cycle;\path[fill=black,line width=0.048pt] (64.1323,39.3890) -- (64.1323,41.3519) --
  (67.2156,41.3519) -- (67.2156,42.1706) -- (64.1323,42.1706) --
  (64.1323,44.0776) -- (67.5700,44.0776) -- (67.5700,44.8963) --
  (63.2906,44.8963) -- (63.2906,38.5703) -- (67.5700,38.5703) --
  (67.5700,39.3890) -- cycle;\path[fill=black,line width=0.048pt] (68.8282,38.5703) -- (71.0078,38.5703) ..
  controls (71.9352,38.5703) and (72.6972,38.8649) .. (73.2937,39.4541) ..
  controls (73.8962,40.0433) and (74.1975,40.7906) .. (74.1975,41.6961) ..
  controls (74.1975,42.6264) and (73.8962,43.3923) .. (73.2937,43.9939) ..
  controls (72.7031,44.5955) and (71.9411,44.8963) .. (71.0078,44.8963) --
  (68.8282,44.8963) -- cycle(69.6699,44.0776) -- (71.0078,44.0776) .. controls
  (71.6812,44.0776) and (72.2335,43.8544) .. (72.6647,43.4078) .. controls
  (73.1077,42.9489) and (73.3292,42.3783) .. (73.3292,41.6961) .. controls
  (73.3292,41.0387) and (73.1077,40.4898) .. (72.6647,40.0495) .. controls
  (72.2276,39.6091) and (71.6753,39.3890) .. (71.0078,39.3890) --
  (69.6699,39.3890) -- cycle;\path[fill=black,line width=0.048pt] (18.6503,55.3044) -- (17.9503,54.6346) ..
  controls (17.4246,55.1679) and (16.7985,55.4346) .. (16.0719,55.4346) ..
  controls (15.4163,55.4346) and (14.8906,55.2579) .. (14.4948,54.9044) ..
  controls (14.0991,54.5508) and (13.9012,54.0919) .. (13.9012,53.5275) ..
  controls (13.9012,53.1430) and (14.0046,52.8019) .. (14.2113,52.5042) ..
  controls (14.4180,52.2003) and (14.7016,51.9646) .. (15.0619,51.7972) ..
  controls (14.7252,51.3816) and (14.5568,50.9661) .. (14.5568,50.5506) ..
  controls (14.5568,50.0482) and (14.7252,49.6389) .. (15.0619,49.3226) ..
  controls (15.4045,49.0001) and (15.8357,48.8388) .. (16.3555,48.8388) ..
  controls (16.8575,48.8388) and (17.2651,48.9970) .. (17.5782,49.3133) ..
  controls (17.8971,49.6234) and (18.0862,50.0482) .. (18.1452,50.5878) --
  (17.3301,50.5878) .. controls (17.2888,50.2715) and (17.1795,50.0296) ..
  (17.0023,49.8622) .. controls (16.8310,49.6885) and (16.6124,49.6017) ..
  (16.3466,49.6017) .. controls (16.0690,49.6017) and (15.8416,49.6916) ..
  (15.6644,49.8715) .. controls (15.4872,50.0513) and (15.3986,50.2839) ..
  (15.3986,50.5692) .. controls (15.3986,50.8483) and (15.4901,51.1026) ..
  (15.6732,51.3320) .. controls (15.8622,51.5615) and (16.2167,51.9243) ..
  (16.7365,52.4205) -- (16.8251,52.5042) -- (17.8883,53.5275) --
  (18.7743,52.1507) -- (19.6249,52.1507) -- (18.4819,53.9648) --
  (18.4287,54.0485) -- (19.7312,55.3044) -- cycle(16.0808,54.6718) .. controls
  (16.5652,54.6718) and (17.0082,54.4888) .. (17.4098,54.1229) --
  (16.1162,52.8856) .. controls (15.8741,52.6562) and (15.6910,52.4763) ..
  (15.5669,52.3461) .. controls (14.9880,52.6065) and (14.6986,52.9973) ..
  (14.6986,53.5182) .. controls (14.6986,53.8407) and (14.8286,54.1136) ..
  (15.0885,54.3369) .. controls (15.3484,54.5601) and (15.6791,54.6718) ..
  (16.0808,54.6718) -- cycle;\path[fill=black,line width=0.048pt] (26.8855,53.6903) -- (23.8730,53.6903) --
  (23.1996,55.3183) -- (22.2959,55.3183) -- (24.9805,48.9551) --
  (25.7779,48.9551) -- (28.4626,55.3183) -- (27.5588,55.3183) --
  cycle(26.5488,52.8717) -- (25.3792,50.0250) -- (24.2097,52.8717) -- cycle;\path[fill=black,line width=0.048pt] (33.5665,54.4997) -- (33.5665,55.2997) --
  (29.4553,55.2997) -- (29.4553,48.9737) -- (30.2971,48.9737) --
  (30.2971,54.4997) -- cycle;\path[fill=black,line width=0.048pt] (38.6888,54.4997) -- (38.6888,55.2997) --
  (34.5777,55.2997) -- (34.5777,48.9737) -- (35.4194,48.9737) --
  (35.4194,54.4997) -- cycle;\path[fill=black,line width=0.048pt] (39.6400,49.8110) .. controls
  (40.2543,49.1660) and (40.9956,48.8435) .. (41.8639,48.8435) .. controls
  (42.7263,48.8435) and (43.4647,49.1660) .. (44.0790,49.8110) .. controls
  (44.6933,50.4560) and (45.0004,51.2313) .. (45.0004,52.1367) .. controls
  (45.0004,53.0484) and (44.6933,53.8268) .. (44.0790,54.4718) .. controls
  (43.4706,55.1106) and (42.7322,55.4300) .. (41.8639,55.4300) .. controls
  (40.9897,55.4300) and (40.2484,55.1106) .. (39.6400,54.4718) .. controls
  (39.0316,53.8330) and (38.7274,53.0546) .. (38.7274,52.1367) .. controls
  (38.7274,51.2250) and (39.0316,50.4498) .. (39.6400,49.8110) --
  cycle(41.8639,54.6020) .. controls (42.5019,54.6020) and (43.0394,54.3632) ..
  (43.4765,53.8857) .. controls (43.9195,53.4019) and (44.1410,52.8190) ..
  (44.1410,52.1367) .. controls (44.1410,51.4607) and (43.9195,50.8808) ..
  (43.4765,50.3971) .. controls (43.0335,49.9133) and (42.4960,49.6715) ..
  (41.8639,49.6715) .. controls (41.2201,49.6715) and (40.6767,49.9133) ..
  (40.2336,50.3971) .. controls (39.7965,50.8746) and (39.5780,51.4545) ..
  (39.5780,52.1367) .. controls (39.5780,52.8252) and (39.7965,53.4081) ..
  (40.2336,53.8857) .. controls (40.6707,54.3632) and (41.2142,54.6020) ..
  (41.8639,54.6020) -- cycle;\path[fill=black,line width=0.048pt] (50.9229,48.9737) -- (48.4687,53.1042) --
  (48.4687,55.2997) -- (47.6092,55.2997) -- (47.6092,53.1042) --
  (45.1549,48.9737) -- (46.0941,48.9737) -- (48.0345,52.2484) --
  (49.9838,48.9737) -- cycle;\path[fill=black,line width=0.048pt] (55.6517,48.8388) .. controls
  (56.2779,48.8388) and (56.8154,49.0001) .. (57.2643,49.3226) .. controls
  (57.7191,49.6389) and (57.9820,50.1227) .. (58.0529,50.7739) --
  (57.2111,50.7739) .. controls (57.0812,49.9924) and (56.5791,49.6017) ..
  (55.7049,49.6017) .. controls (55.2973,49.6017) and (54.9695,49.6978) ..
  (54.7214,49.8901) .. controls (54.4733,50.0823) and (54.3493,50.3397) ..
  (54.3493,50.6622) .. controls (54.3493,51.1460) and (54.7185,51.4685) ..
  (55.4568,51.6297) -- (56.2985,51.8065) .. controls (56.8833,51.9305) and
  (57.3175,52.1290) .. (57.6010,52.4019) .. controls (57.8904,52.6686) and
  (58.0351,53.0314) .. (58.0351,53.4903) .. controls (58.0351,54.0671) and
  (57.8255,54.5384) .. (57.4061,54.9044) .. controls (56.9926,55.2579) and
  (56.4403,55.4346) .. (55.7492,55.4346) .. controls (55.0168,55.4346) and
  (54.4231,55.2424) .. (53.9683,54.8578) .. controls (53.5135,54.4733) and
  (53.2565,53.9400) .. (53.1975,53.2577) -- (54.0392,53.2577) .. controls
  (54.1042,53.7043) and (54.2873,54.0485) .. (54.5885,54.2904) .. controls
  (54.8957,54.5260) and (55.2973,54.6439) .. (55.7935,54.6439) .. controls
  (56.2188,54.6439) and (56.5525,54.5508) .. (56.7947,54.3648) .. controls
  (57.0428,54.1725) and (57.1668,53.9182) .. (57.1668,53.6019) .. controls
  (57.1668,53.1120) and (56.8302,52.7988) .. (56.1568,52.6624) --
  (55.1379,52.4577) .. controls (54.0156,52.2344) and (53.4544,51.6576) ..
  (53.4544,50.7273) .. controls (53.4544,50.1630) and (53.6552,49.7071) ..
  (54.0569,49.3598) .. controls (54.4586,49.0125) and (54.9902,48.8388) ..
  (55.6517,48.8388) -- cycle;\path[fill=black,line width=0.048pt] (64.2175,48.9737) -- (61.7632,53.1042) --
  (61.7632,55.2997) -- (60.9038,55.2997) -- (60.9038,53.1042) --
  (58.4495,48.9737) -- (59.3887,48.9737) -- (61.3291,52.2484) --
  (63.2783,48.9737) -- cycle;\path[fill=black,line width=0.048pt] (66.9347,48.8388) .. controls
  (67.5608,48.8388) and (68.0983,49.0001) .. (68.5472,49.3226) .. controls
  (69.0021,49.6389) and (69.2649,50.1227) .. (69.3358,50.7739) --
  (68.4941,50.7739) .. controls (68.3641,49.9924) and (67.8621,49.6017) ..
  (66.9878,49.6017) .. controls (66.5803,49.6017) and (66.2525,49.6978) ..
  (66.0044,49.8901) .. controls (65.7563,50.0823) and (65.6322,50.3397) ..
  (65.6322,50.6622) .. controls (65.6322,51.1460) and (66.0014,51.4685) ..
  (66.7398,51.6297) -- (67.5815,51.8065) .. controls (68.1663,51.9305) and
  (68.6004,52.1290) .. (68.8839,52.4019) .. controls (69.1734,52.6686) and
  (69.3181,53.0314) .. (69.3181,53.4903) .. controls (69.3181,54.0671) and
  (69.1084,54.5384) .. (68.6890,54.9044) .. controls (68.2755,55.2579) and
  (67.7232,55.4346) .. (67.0321,55.4346) .. controls (66.2997,55.4346) and
  (65.7061,55.2424) .. (65.2512,54.8578) .. controls (64.7964,54.4733) and
  (64.5395,53.9400) .. (64.4804,53.2577) -- (65.3221,53.2577) .. controls
  (65.3871,53.7043) and (65.5702,54.0485) .. (65.8715,54.2904) .. controls
  (66.1786,54.5260) and (66.5803,54.6439) .. (67.0764,54.6439) .. controls
  (67.5017,54.6439) and (67.8355,54.5508) .. (68.0777,54.3648) .. controls
  (68.3257,54.1725) and (68.4498,53.9182) .. (68.4498,53.6019) .. controls
  (68.4498,53.1120) and (68.1131,52.7988) .. (67.4397,52.6624) --
  (66.4208,52.4577) .. controls (65.2985,52.2344) and (64.7373,51.6576) ..
  (64.7373,50.7273) .. controls (64.7373,50.1630) and (64.9382,49.7071) ..
  (65.3398,49.3598) .. controls (65.7415,49.0125) and (66.2731,48.8388) ..
  (66.9347,48.8388) -- cycle;\path[fill=black,line width=0.048pt] (75.0140,49.7924) -- (72.8432,49.7924) --
  (72.8432,55.2997) -- (72.0015,55.2997) -- (72.0015,49.7924) --
  (69.8307,49.7924) -- (69.8307,48.9737) -- (75.0140,48.9737) -- cycle;\path[fill=black,line width=0.048pt] (76.5489,49.7924) -- (76.5489,51.7553) --
  (79.6323,51.7553) -- (79.6323,52.5740) -- (76.5489,52.5740) --
  (76.5489,54.4811) -- (79.9867,54.4811) -- (79.9867,55.2997) --
  (75.7072,55.2997) -- (75.7072,48.9737) -- (79.9867,48.9737) --
  (79.9867,49.7924) -- cycle;\path[fill=black,line width=0.051pt] (87.6628,48.9737) -- (87.6628,55.2997) --
  (86.7030,55.2997) -- (86.7030,50.9367) -- (84.5834,54.1741) --
  (83.9735,54.1741) -- (81.8339,50.9087) -- (81.8539,55.2997) --
  (80.8941,55.2997) -- (80.8941,48.9737) -- (81.5940,48.9737) --
  (84.2734,53.0763) -- (86.9529,48.9737) -- cycle;\path[fill=black,line width=0.048pt] (91.1792,48.8388) .. controls
  (91.8053,48.8388) and (92.3428,49.0001) .. (92.7917,49.3226) .. controls
  (93.2465,49.6389) and (93.5094,50.1227) .. (93.5803,50.7739) --
  (92.7386,50.7739) .. controls (92.6086,49.9924) and (92.1065,49.6017) ..
  (91.2323,49.6017) .. controls (90.8248,49.6017) and (90.4969,49.6978) ..
  (90.2488,49.8901) .. controls (90.0008,50.0823) and (89.8767,50.3397) ..
  (89.8767,50.6622) .. controls (89.8767,51.1460) and (90.2459,51.4685) ..
  (90.9842,51.6297) -- (91.8260,51.8065) .. controls (92.4107,51.9305) and
  (92.8449,52.1290) .. (93.1284,52.4019) .. controls (93.4178,52.6686) and
  (93.5626,53.0314) .. (93.5626,53.4903) .. controls (93.5626,54.0671) and
  (93.3529,54.5384) .. (92.9335,54.9044) .. controls (92.5200,55.2579) and
  (91.9677,55.4346) .. (91.2766,55.4346) .. controls (90.5442,55.4346) and
  (89.9505,55.2424) .. (89.4957,54.8578) .. controls (89.0409,54.4733) and
  (88.7839,53.9400) .. (88.7249,53.2577) -- (89.5666,53.2577) .. controls
  (89.6316,53.7043) and (89.8147,54.0485) .. (90.1159,54.2904) .. controls
  (90.4231,54.5260) and (90.8248,54.6439) .. (91.3209,54.6439) .. controls
  (91.7462,54.6439) and (92.0799,54.5508) .. (92.3221,54.3648) .. controls
  (92.5702,54.1725) and (92.6943,53.9182) .. (92.6943,53.6019) .. controls
  (92.6943,53.1120) and (92.3576,52.7988) .. (91.6842,52.6624) --
  (90.6653,52.4577) .. controls (89.5430,52.2344) and (88.9818,51.6576) ..
  (88.9818,50.7273) .. controls (88.9818,50.1630) and (89.1827,49.7071) ..
  (89.5843,49.3598) .. controls (89.9860,49.0125) and (90.5176,48.8388) ..
  (91.1792,48.8388) -- cycle;%
\end{tikzpicture}%
}}
\renewcommand{\headrulewidth}{0.4pt}}

%--Embedded documents overlay header--%
\fancypagestyle{style2}{
\fancyhf{}
\fancyhead[L]{\fontsize{25}{25}\selectfont{\hypersetup{linkcolor=pfblue}\hyperref[sec:firstpage]{\textbf{\the\year \xspace QA}}}}
%\fancyhead[L]{\fontsize{25}{25}\selectfont{\hypersetup{linkcolor=pfblue}\hyperref[sec:firstpage]{\textbf{\input{name}}}}}
\fancyfoot[R]{\raisebox{0em}{
\definecolor{c3486c7}{RGB}{52,134,199}


\begin{tikzpicture}[y=0.80pt, x=0.80pt, yscale=-1.000000, xscale=1.000000, inner sep=0pt, outer sep=0pt]
\path[fill=c3486c7,line width=0.153pt] (127.8015,27.8018) -- (125.1667,27.8018)
  -- (125.1667,33.4979) -- (120.2802,33.4979) -- (120.2802,13.8363) --
  (128.8677,13.8363) .. controls (131.3350,13.8363) and (133.3108,14.4627) ..
  (134.7951,15.7157) .. controls (136.2794,16.9686) and (137.0215,18.6457) ..
  (137.0215,20.7468) .. controls (137.0215,22.0961) and (136.6938,23.3008) ..
  (136.0384,24.3610) .. controls (135.3830,25.4212) and (134.4578,26.2501) ..
  (133.2626,26.8477) -- (137.3918,33.4979) -- (131.5223,33.4979) --
  cycle(125.1667,23.4069) -- (128.7809,23.4069) .. controls (130.1652,23.3346)
  and (131.2985,22.0074) .. (131.3180,20.7468) .. controls (131.2025,19.1994)
  and (129.7217,18.1538) .. (128.7809,18.1445) -- (125.1667,18.1445) --
  cycle(112.1053,18.0867) -- (112.1053,21.6720) -- (117.6624,21.6720) --
  (117.6624,25.8357) -- (112.1053,25.8357) -- (112.1053,29.2475) --
  (122.0807,29.2475) -- (122.0807,33.4979) -- (107.2188,33.4979) --
  (107.2188,13.8363) -- (118.9057,13.8363) -- (118.9057,18.0867) --
  cycle(110.6994,29.1897) -- (110.6994,33.4979) -- (96.0977,33.4979) --
  (96.0977,13.8363) -- (100.9842,13.8363) -- (100.9842,29.1897) --
  cycle(76.3949,13.8363) -- (84.2884,13.8363) .. controls (87.3726,13.8363) and
  (89.9171,14.7519) .. (91.9218,16.5831) .. controls (93.9072,18.3951) and
  (94.9000,20.7371) .. (94.9000,23.6093) .. controls (94.9000,26.5200) and
  (93.9072,28.9006) .. (91.9218,30.7511) .. controls (89.9171,32.5823) and
  (87.3726,33.4979) .. (84.2884,33.4979) -- (76.3949,33.4979) --
  cycle(81.2814,29.1030) -- (84.2884,29.1030) .. controls (85.8884,29.1030) and
  (87.2281,28.5825) .. (88.3075,27.5416) .. controls (89.3870,26.4621) and
  (89.9267,25.1514) .. (89.9267,23.6093) .. controls (89.9267,22.0865) and
  (89.3870,20.8142) .. (88.3075,19.7926) .. controls (87.2473,18.7517) and
  (85.9076,18.2312) .. (84.2884,18.2312) -- (81.2814,18.2312) --
  cycle(74.8185,13.8363) -- (74.8185,24.3321) .. controls (74.8185,27.3585) and
  (73.9993,29.7198) .. (72.3608,31.4161) .. controls (70.7031,33.0931) and
  (68.4670,33.9316) .. (65.6527,33.9316) .. controls (62.7998,33.9316) and
  (60.5638,33.0931) .. (58.9446,31.4161) .. controls (57.3254,29.7391) and
  (56.5158,27.3778) .. (56.5158,24.3321) -- (56.5158,13.8363) --
  (61.4312,13.8363) -- (61.4312,24.4189) .. controls (61.4312,25.9224) and
  (61.8071,27.1272) .. (62.5589,28.0331) .. controls (63.3299,28.9198) and
  (64.3612,29.3632) .. (65.6527,29.3632) .. controls (66.9442,29.3632) and
  (67.9755,28.9198) .. (68.7465,28.0331) .. controls (69.5176,27.1272) and
  (69.9031,25.9224) .. (69.9031,24.4189) -- (69.9031,13.8363) --
  cycle(53.6176,30.1150) -- (45.2903,30.1150) -- (43.9891,33.4979) --
  (38.7846,33.4979) -- (47.1408,13.7205) -- (51.8249,13.7206) --
  (60.1522,33.4979) -- (54.9476,33.4979) -- cycle(51.9984,25.9513) --
  (49.4539,19.4167) -- (46.9095,25.9513) -- cycle(36.2805,18.0261) --
  (36.2805,22.1898) -- (40.4902,22.1898) -- (40.4902,26.4402) --
  (36.2805,26.4402) -- (36.2805,33.3796) -- (31.3940,33.3796) --
  (31.3940,13.7205) -- (47.1408,13.7205) -- (47.1408,18.0261) --
  cycle(13.6824,13.8363) -- (21.6049,13.8363) .. controls (24.1879,13.8363) and
  (26.2119,14.4820) .. (27.6769,15.7735) .. controls (29.1804,17.0650) and
  (29.9322,18.8191) .. (29.9322,21.0359) .. controls (29.9322,23.3298) and
  (29.1515,25.1514) .. (27.5901,26.5007) .. controls (26.0288,27.8307) and
  (23.9084,28.4958) .. (21.2290,28.4958) -- (18.5689,28.4958) --
  (18.5689,33.4979) -- (13.6824,33.4979) -- cycle(18.5689,24.2165) --
  (21.5182,24.2165) .. controls (22.5976,24.2165) and (23.4554,23.9273) ..
  (24.0915,23.3490) .. controls (24.7469,22.7708) and (25.0746,21.9997) ..
  (25.0746,21.0359) .. controls (25.0746,20.1492) and (24.7469,19.4456) ..
  (24.0915,18.9252) .. controls (23.4554,18.4047) and (22.5976,18.1445) ..
  (21.5182,18.1445) -- (18.5689,18.1445) -- cycle;\path[fill=black,line width=0.048pt] (20.1299,41.7333) -- (20.1299,44.8963) --
  (19.4565,44.8963) -- (19.3945,43.7334) .. controls (18.9338,44.5955) and
  (18.1925,45.0265) .. (17.1706,45.0265) .. controls (16.2787,45.0265) and
  (15.5374,44.7102) .. (14.9467,44.0776) .. controls (14.3501,43.4388) and
  (14.0518,42.6760) .. (14.0518,41.7891) .. controls (14.0518,40.8278) and
  (14.3560,40.0309) .. (14.9644,39.3983) .. controls (15.5728,38.7595) and
  (16.3230,38.4401) .. (17.2149,38.4401) .. controls (17.8883,38.4401) and
  (18.4967,38.6416) .. (19.0401,39.0448) .. controls (19.5718,39.4417) and
  (19.9232,39.9937) .. (20.0945,40.7007) -- (19.1908,40.7007) .. controls
  (19.0431,40.2355) and (18.7861,39.8851) .. (18.4199,39.6495) .. controls
  (18.0537,39.4138) and (17.6550,39.2959) .. (17.2238,39.2959) .. controls
  (16.5622,39.2959) and (16.0099,39.5285) .. (15.5669,39.9937) .. controls
  (15.1239,40.4588) and (14.9024,41.0573) .. (14.9024,41.7891) .. controls
  (14.9024,42.4465) and (15.1151,43.0078) .. (15.5403,43.4730) .. controls
  (15.9715,43.9381) and (16.5150,44.1707) .. (17.1706,44.1707) .. controls
  (17.7141,44.1707) and (18.1659,44.0063) .. (18.5262,43.6776) .. controls
  (18.8866,43.3427) and (19.1199,42.8931) .. (19.2262,42.3287) --
  (17.1086,42.3287) -- (17.1086,41.5100) -- (20.1299,41.5100) -- cycle;\path[fill=black,line width=0.048pt] (25.8949,44.0963) -- (25.8949,44.8963) --
  (21.7837,44.8963) -- (21.7837,38.5703) -- (22.6255,38.5703) --
  (22.6255,44.0963) -- cycle;\path[fill=black,line width=0.048pt] (30.9640,43.2683) -- (27.9516,43.2683) --
  (27.2782,44.8963) -- (26.3744,44.8963) -- (29.0591,38.5331) --
  (29.8565,38.5331) -- (32.5412,44.8963) -- (31.6374,44.8963) --
  cycle(30.6273,42.4496) -- (29.4578,39.6029) -- (28.2882,42.4496) -- cycle;\path[fill=black,line width=0.048pt] (35.4886,38.4308) .. controls
  (36.1147,38.4308) and (36.6522,38.5920) .. (37.1011,38.9145) .. controls
  (37.5559,39.2308) and (37.8188,39.7146) .. (37.8897,40.3658) --
  (37.0480,40.3658) .. controls (36.9180,39.5843) and (36.4159,39.1936) ..
  (35.5417,39.1936) .. controls (35.1342,39.1936) and (34.8063,39.2897) ..
  (34.5582,39.4820) .. controls (34.3102,39.6743) and (34.1861,39.9316) ..
  (34.1861,40.2541) .. controls (34.1861,40.7379) and (34.5553,41.0604) ..
  (35.2936,41.2217) -- (36.1354,41.3984) .. controls (36.7201,41.5224) and
  (37.1543,41.7209) .. (37.4378,41.9938) .. controls (37.7272,42.2605) and
  (37.8720,42.6233) .. (37.8720,43.0822) .. controls (37.8720,43.6590) and
  (37.6623,44.1304) .. (37.2429,44.4963) .. controls (36.8294,44.8498) and
  (36.2771,45.0265) .. (35.5860,45.0265) .. controls (34.8536,45.0265) and
  (34.2599,44.8343) .. (33.8051,44.4498) .. controls (33.3503,44.0652) and
  (33.0933,43.5319) .. (33.0343,42.8497) -- (33.8760,42.8497) .. controls
  (33.9410,43.2962) and (34.1241,43.6404) .. (34.4253,43.8823) .. controls
  (34.7325,44.1180) and (35.1342,44.2358) .. (35.6303,44.2358) .. controls
  (36.0556,44.2358) and (36.3893,44.1428) .. (36.6315,43.9567) .. controls
  (36.8796,43.7645) and (37.0037,43.5102) .. (37.0037,43.1939) .. controls
  (37.0037,42.7039) and (36.6670,42.3907) .. (35.9936,42.2543) --
  (34.9747,42.0496) .. controls (33.8524,41.8263) and (33.2912,41.2496) ..
  (33.2912,40.3193) .. controls (33.2912,39.7549) and (33.4921,39.2990) ..
  (33.8937,38.9517) .. controls (34.2954,38.6044) and (34.8270,38.4308) ..
  (35.4886,38.4308) -- cycle;\path[fill=black,line width=0.048pt] (41.1992,38.4308) .. controls
  (41.8254,38.4308) and (42.3629,38.5920) .. (42.8118,38.9145) .. controls
  (43.2666,39.2308) and (43.5295,39.7146) .. (43.6004,40.3658) --
  (42.7586,40.3658) .. controls (42.6287,39.5843) and (42.1266,39.1936) ..
  (41.2524,39.1936) .. controls (40.8448,39.1936) and (40.5170,39.2897) ..
  (40.2689,39.4820) .. controls (40.0208,39.6743) and (39.8968,39.9316) ..
  (39.8968,40.2541) .. controls (39.8968,40.7379) and (40.2660,41.0604) ..
  (41.0043,41.2217) -- (41.8460,41.3984) .. controls (42.4308,41.5224) and
  (42.8650,41.7209) .. (43.1485,41.9938) .. controls (43.4379,42.2605) and
  (43.5826,42.6233) .. (43.5826,43.0822) .. controls (43.5826,43.6590) and
  (43.3730,44.1304) .. (42.9536,44.4963) .. controls (42.5401,44.8498) and
  (41.9878,45.0265) .. (41.2967,45.0265) .. controls (40.5643,45.0265) and
  (39.9706,44.8343) .. (39.5158,44.4498) .. controls (39.0610,44.0652) and
  (38.8040,43.5319) .. (38.7450,42.8497) -- (39.5867,42.8497) .. controls
  (39.6517,43.2962) and (39.8348,43.6404) .. (40.1360,43.8823) .. controls
  (40.4432,44.1180) and (40.8448,44.2358) .. (41.3410,44.2358) .. controls
  (41.7663,44.2358) and (42.1000,44.1428) .. (42.3422,43.9567) .. controls
  (42.5903,43.7645) and (42.7143,43.5102) .. (42.7143,43.1939) .. controls
  (42.7143,42.7039) and (42.3777,42.3907) .. (41.7043,42.2543) --
  (40.6854,42.0496) .. controls (39.5631,41.8263) and (39.0019,41.2496) ..
  (39.0019,40.3193) .. controls (39.0019,39.7549) and (39.2028,39.2990) ..
  (39.6044,38.9517) .. controls (40.0061,38.6044) and (40.5377,38.4308) ..
  (41.1992,38.4308) -- cycle;\path[fill=black,line width=0.048pt] (47.2977,41.9845) -- (47.2977,42.8404) --
  (44.3207,42.8404) -- (44.3207,41.9845) -- cycle;\path[fill=black,line width=0.048pt] (53.1596,44.0963) -- (53.1596,44.8963) --
  (49.0484,44.8963) -- (49.0484,38.5703) -- (49.8902,38.5703) --
  (49.8902,44.0963) -- cycle;\path[fill=black,line width=0.048pt] (55.0125,38.5703) -- (55.0125,44.8963) --
  (54.1708,44.8963) -- (54.1708,38.5703) -- cycle;\path[fill=black,line width=0.048pt] (61.7602,38.5703) -- (61.7602,44.8963) --
  (61.0691,44.8963) -- (57.3833,40.2076) -- (57.3833,44.8963) --
  (56.5416,44.8963) -- (56.5416,38.5703) -- (57.2061,38.5703) --
  (60.9096,43.3334) -- (60.9096,38.5703) -- cycle;\path[fill=black,line width=0.048pt] (64.1323,39.3890) -- (64.1323,41.3519) --
  (67.2156,41.3519) -- (67.2156,42.1706) -- (64.1323,42.1706) --
  (64.1323,44.0776) -- (67.5700,44.0776) -- (67.5700,44.8963) --
  (63.2906,44.8963) -- (63.2906,38.5703) -- (67.5700,38.5703) --
  (67.5700,39.3890) -- cycle;\path[fill=black,line width=0.048pt] (68.8282,38.5703) -- (71.0078,38.5703) ..
  controls (71.9352,38.5703) and (72.6972,38.8649) .. (73.2937,39.4541) ..
  controls (73.8962,40.0433) and (74.1975,40.7906) .. (74.1975,41.6961) ..
  controls (74.1975,42.6264) and (73.8962,43.3923) .. (73.2937,43.9939) ..
  controls (72.7031,44.5955) and (71.9411,44.8963) .. (71.0078,44.8963) --
  (68.8282,44.8963) -- cycle(69.6699,44.0776) -- (71.0078,44.0776) .. controls
  (71.6812,44.0776) and (72.2335,43.8544) .. (72.6647,43.4078) .. controls
  (73.1077,42.9489) and (73.3292,42.3783) .. (73.3292,41.6961) .. controls
  (73.3292,41.0387) and (73.1077,40.4898) .. (72.6647,40.0495) .. controls
  (72.2276,39.6091) and (71.6753,39.3890) .. (71.0078,39.3890) --
  (69.6699,39.3890) -- cycle;\path[fill=black,line width=0.048pt] (18.6503,55.3044) -- (17.9503,54.6346) ..
  controls (17.4246,55.1679) and (16.7985,55.4346) .. (16.0719,55.4346) ..
  controls (15.4163,55.4346) and (14.8906,55.2579) .. (14.4948,54.9044) ..
  controls (14.0991,54.5508) and (13.9012,54.0919) .. (13.9012,53.5275) ..
  controls (13.9012,53.1430) and (14.0046,52.8019) .. (14.2113,52.5042) ..
  controls (14.4180,52.2003) and (14.7016,51.9646) .. (15.0619,51.7972) ..
  controls (14.7252,51.3816) and (14.5568,50.9661) .. (14.5568,50.5506) ..
  controls (14.5568,50.0482) and (14.7252,49.6389) .. (15.0619,49.3226) ..
  controls (15.4045,49.0001) and (15.8357,48.8388) .. (16.3555,48.8388) ..
  controls (16.8575,48.8388) and (17.2651,48.9970) .. (17.5782,49.3133) ..
  controls (17.8971,49.6234) and (18.0862,50.0482) .. (18.1452,50.5878) --
  (17.3301,50.5878) .. controls (17.2888,50.2715) and (17.1795,50.0296) ..
  (17.0023,49.8622) .. controls (16.8310,49.6885) and (16.6124,49.6017) ..
  (16.3466,49.6017) .. controls (16.0690,49.6017) and (15.8416,49.6916) ..
  (15.6644,49.8715) .. controls (15.4872,50.0513) and (15.3986,50.2839) ..
  (15.3986,50.5692) .. controls (15.3986,50.8483) and (15.4901,51.1026) ..
  (15.6732,51.3320) .. controls (15.8622,51.5615) and (16.2167,51.9243) ..
  (16.7365,52.4205) -- (16.8251,52.5042) -- (17.8883,53.5275) --
  (18.7743,52.1507) -- (19.6249,52.1507) -- (18.4819,53.9648) --
  (18.4287,54.0485) -- (19.7312,55.3044) -- cycle(16.0808,54.6718) .. controls
  (16.5652,54.6718) and (17.0082,54.4888) .. (17.4098,54.1229) --
  (16.1162,52.8856) .. controls (15.8741,52.6562) and (15.6910,52.4763) ..
  (15.5669,52.3461) .. controls (14.9880,52.6065) and (14.6986,52.9973) ..
  (14.6986,53.5182) .. controls (14.6986,53.8407) and (14.8286,54.1136) ..
  (15.0885,54.3369) .. controls (15.3484,54.5601) and (15.6791,54.6718) ..
  (16.0808,54.6718) -- cycle;\path[fill=black,line width=0.048pt] (26.8855,53.6903) -- (23.8730,53.6903) --
  (23.1996,55.3183) -- (22.2959,55.3183) -- (24.9805,48.9551) --
  (25.7779,48.9551) -- (28.4626,55.3183) -- (27.5588,55.3183) --
  cycle(26.5488,52.8717) -- (25.3792,50.0250) -- (24.2097,52.8717) -- cycle;\path[fill=black,line width=0.048pt] (33.5665,54.4997) -- (33.5665,55.2997) --
  (29.4553,55.2997) -- (29.4553,48.9737) -- (30.2971,48.9737) --
  (30.2971,54.4997) -- cycle;\path[fill=black,line width=0.048pt] (38.6888,54.4997) -- (38.6888,55.2997) --
  (34.5777,55.2997) -- (34.5777,48.9737) -- (35.4194,48.9737) --
  (35.4194,54.4997) -- cycle;\path[fill=black,line width=0.048pt] (39.6400,49.8110) .. controls
  (40.2543,49.1660) and (40.9956,48.8435) .. (41.8639,48.8435) .. controls
  (42.7263,48.8435) and (43.4647,49.1660) .. (44.0790,49.8110) .. controls
  (44.6933,50.4560) and (45.0004,51.2313) .. (45.0004,52.1367) .. controls
  (45.0004,53.0484) and (44.6933,53.8268) .. (44.0790,54.4718) .. controls
  (43.4706,55.1106) and (42.7322,55.4300) .. (41.8639,55.4300) .. controls
  (40.9897,55.4300) and (40.2484,55.1106) .. (39.6400,54.4718) .. controls
  (39.0316,53.8330) and (38.7274,53.0546) .. (38.7274,52.1367) .. controls
  (38.7274,51.2250) and (39.0316,50.4498) .. (39.6400,49.8110) --
  cycle(41.8639,54.6020) .. controls (42.5019,54.6020) and (43.0394,54.3632) ..
  (43.4765,53.8857) .. controls (43.9195,53.4019) and (44.1410,52.8190) ..
  (44.1410,52.1367) .. controls (44.1410,51.4607) and (43.9195,50.8808) ..
  (43.4765,50.3971) .. controls (43.0335,49.9133) and (42.4960,49.6715) ..
  (41.8639,49.6715) .. controls (41.2201,49.6715) and (40.6767,49.9133) ..
  (40.2336,50.3971) .. controls (39.7965,50.8746) and (39.5780,51.4545) ..
  (39.5780,52.1367) .. controls (39.5780,52.8252) and (39.7965,53.4081) ..
  (40.2336,53.8857) .. controls (40.6707,54.3632) and (41.2142,54.6020) ..
  (41.8639,54.6020) -- cycle;\path[fill=black,line width=0.048pt] (50.9229,48.9737) -- (48.4687,53.1042) --
  (48.4687,55.2997) -- (47.6092,55.2997) -- (47.6092,53.1042) --
  (45.1549,48.9737) -- (46.0941,48.9737) -- (48.0345,52.2484) --
  (49.9838,48.9737) -- cycle;\path[fill=black,line width=0.048pt] (55.6517,48.8388) .. controls
  (56.2779,48.8388) and (56.8154,49.0001) .. (57.2643,49.3226) .. controls
  (57.7191,49.6389) and (57.9820,50.1227) .. (58.0529,50.7739) --
  (57.2111,50.7739) .. controls (57.0812,49.9924) and (56.5791,49.6017) ..
  (55.7049,49.6017) .. controls (55.2973,49.6017) and (54.9695,49.6978) ..
  (54.7214,49.8901) .. controls (54.4733,50.0823) and (54.3493,50.3397) ..
  (54.3493,50.6622) .. controls (54.3493,51.1460) and (54.7185,51.4685) ..
  (55.4568,51.6297) -- (56.2985,51.8065) .. controls (56.8833,51.9305) and
  (57.3175,52.1290) .. (57.6010,52.4019) .. controls (57.8904,52.6686) and
  (58.0351,53.0314) .. (58.0351,53.4903) .. controls (58.0351,54.0671) and
  (57.8255,54.5384) .. (57.4061,54.9044) .. controls (56.9926,55.2579) and
  (56.4403,55.4346) .. (55.7492,55.4346) .. controls (55.0168,55.4346) and
  (54.4231,55.2424) .. (53.9683,54.8578) .. controls (53.5135,54.4733) and
  (53.2565,53.9400) .. (53.1975,53.2577) -- (54.0392,53.2577) .. controls
  (54.1042,53.7043) and (54.2873,54.0485) .. (54.5885,54.2904) .. controls
  (54.8957,54.5260) and (55.2973,54.6439) .. (55.7935,54.6439) .. controls
  (56.2188,54.6439) and (56.5525,54.5508) .. (56.7947,54.3648) .. controls
  (57.0428,54.1725) and (57.1668,53.9182) .. (57.1668,53.6019) .. controls
  (57.1668,53.1120) and (56.8302,52.7988) .. (56.1568,52.6624) --
  (55.1379,52.4577) .. controls (54.0156,52.2344) and (53.4544,51.6576) ..
  (53.4544,50.7273) .. controls (53.4544,50.1630) and (53.6552,49.7071) ..
  (54.0569,49.3598) .. controls (54.4586,49.0125) and (54.9902,48.8388) ..
  (55.6517,48.8388) -- cycle;\path[fill=black,line width=0.048pt] (64.2175,48.9737) -- (61.7632,53.1042) --
  (61.7632,55.2997) -- (60.9038,55.2997) -- (60.9038,53.1042) --
  (58.4495,48.9737) -- (59.3887,48.9737) -- (61.3291,52.2484) --
  (63.2783,48.9737) -- cycle;\path[fill=black,line width=0.048pt] (66.9347,48.8388) .. controls
  (67.5608,48.8388) and (68.0983,49.0001) .. (68.5472,49.3226) .. controls
  (69.0021,49.6389) and (69.2649,50.1227) .. (69.3358,50.7739) --
  (68.4941,50.7739) .. controls (68.3641,49.9924) and (67.8621,49.6017) ..
  (66.9878,49.6017) .. controls (66.5803,49.6017) and (66.2525,49.6978) ..
  (66.0044,49.8901) .. controls (65.7563,50.0823) and (65.6322,50.3397) ..
  (65.6322,50.6622) .. controls (65.6322,51.1460) and (66.0014,51.4685) ..
  (66.7398,51.6297) -- (67.5815,51.8065) .. controls (68.1663,51.9305) and
  (68.6004,52.1290) .. (68.8839,52.4019) .. controls (69.1734,52.6686) and
  (69.3181,53.0314) .. (69.3181,53.4903) .. controls (69.3181,54.0671) and
  (69.1084,54.5384) .. (68.6890,54.9044) .. controls (68.2755,55.2579) and
  (67.7232,55.4346) .. (67.0321,55.4346) .. controls (66.2997,55.4346) and
  (65.7061,55.2424) .. (65.2512,54.8578) .. controls (64.7964,54.4733) and
  (64.5395,53.9400) .. (64.4804,53.2577) -- (65.3221,53.2577) .. controls
  (65.3871,53.7043) and (65.5702,54.0485) .. (65.8715,54.2904) .. controls
  (66.1786,54.5260) and (66.5803,54.6439) .. (67.0764,54.6439) .. controls
  (67.5017,54.6439) and (67.8355,54.5508) .. (68.0777,54.3648) .. controls
  (68.3257,54.1725) and (68.4498,53.9182) .. (68.4498,53.6019) .. controls
  (68.4498,53.1120) and (68.1131,52.7988) .. (67.4397,52.6624) --
  (66.4208,52.4577) .. controls (65.2985,52.2344) and (64.7373,51.6576) ..
  (64.7373,50.7273) .. controls (64.7373,50.1630) and (64.9382,49.7071) ..
  (65.3398,49.3598) .. controls (65.7415,49.0125) and (66.2731,48.8388) ..
  (66.9347,48.8388) -- cycle;\path[fill=black,line width=0.048pt] (75.0140,49.7924) -- (72.8432,49.7924) --
  (72.8432,55.2997) -- (72.0015,55.2997) -- (72.0015,49.7924) --
  (69.8307,49.7924) -- (69.8307,48.9737) -- (75.0140,48.9737) -- cycle;\path[fill=black,line width=0.048pt] (76.5489,49.7924) -- (76.5489,51.7553) --
  (79.6323,51.7553) -- (79.6323,52.5740) -- (76.5489,52.5740) --
  (76.5489,54.4811) -- (79.9867,54.4811) -- (79.9867,55.2997) --
  (75.7072,55.2997) -- (75.7072,48.9737) -- (79.9867,48.9737) --
  (79.9867,49.7924) -- cycle;\path[fill=black,line width=0.051pt] (87.6628,48.9737) -- (87.6628,55.2997) --
  (86.7030,55.2997) -- (86.7030,50.9367) -- (84.5834,54.1741) --
  (83.9735,54.1741) -- (81.8339,50.9087) -- (81.8539,55.2997) --
  (80.8941,55.2997) -- (80.8941,48.9737) -- (81.5940,48.9737) --
  (84.2734,53.0763) -- (86.9529,48.9737) -- cycle;\path[fill=black,line width=0.048pt] (91.1792,48.8388) .. controls
  (91.8053,48.8388) and (92.3428,49.0001) .. (92.7917,49.3226) .. controls
  (93.2465,49.6389) and (93.5094,50.1227) .. (93.5803,50.7739) --
  (92.7386,50.7739) .. controls (92.6086,49.9924) and (92.1065,49.6017) ..
  (91.2323,49.6017) .. controls (90.8248,49.6017) and (90.4969,49.6978) ..
  (90.2488,49.8901) .. controls (90.0008,50.0823) and (89.8767,50.3397) ..
  (89.8767,50.6622) .. controls (89.8767,51.1460) and (90.2459,51.4685) ..
  (90.9842,51.6297) -- (91.8260,51.8065) .. controls (92.4107,51.9305) and
  (92.8449,52.1290) .. (93.1284,52.4019) .. controls (93.4178,52.6686) and
  (93.5626,53.0314) .. (93.5626,53.4903) .. controls (93.5626,54.0671) and
  (93.3529,54.5384) .. (92.9335,54.9044) .. controls (92.5200,55.2579) and
  (91.9677,55.4346) .. (91.2766,55.4346) .. controls (90.5442,55.4346) and
  (89.9505,55.2424) .. (89.4957,54.8578) .. controls (89.0409,54.4733) and
  (88.7839,53.9400) .. (88.7249,53.2577) -- (89.5666,53.2577) .. controls
  (89.6316,53.7043) and (89.8147,54.0485) .. (90.1159,54.2904) .. controls
  (90.4231,54.5260) and (90.8248,54.6439) .. (91.3209,54.6439) .. controls
  (91.7462,54.6439) and (92.0799,54.5508) .. (92.3221,54.3648) .. controls
  (92.5702,54.1725) and (92.6943,53.9182) .. (92.6943,53.6019) .. controls
  (92.6943,53.1120) and (92.3576,52.7988) .. (91.6842,52.6624) --
  (90.6653,52.4577) .. controls (89.5430,52.2344) and (88.9818,51.6576) ..
  (88.9818,50.7273) .. controls (88.9818,50.1630) and (89.1827,49.7071) ..
  (89.5843,49.3598) .. controls (89.9860,49.0125) and (90.5176,48.8388) ..
  (91.1792,48.8388) -- cycle;%
\end{tikzpicture}%
}}
\renewcommand{\headrulewidth}{1pt}}
%-----------------------------------------------------------------------------------%



%---------------Cover page tikz picture----------------%
\newcolumntype{C}[1]{>{\centering\let\newline\\\arraybackslash\hspace{0pt}}m{#1}}
\newcommand{\pflogoart}{
\begin{tikzpicture}[remember picture, overlay]
\draw node[
	fill=pfblue,
	rounded corners=0.45in,
	regular polygon,
	regular polygon sides=6,
	minimum size=14in,
	text=white
](inner) at (current page.west) {};
\path  ($(inner.center)!0.44!(inner.east)$) node[text=white]{\fontsize{50}{50}\selectfont\textbf{
\begin{tabular}{C{6in}}
\input{name}
\end{tabular}}};
\end{tikzpicture}}
%------------------------------------------------------%




%-----------Table of Contents Formatting-----------%
%----(compiler must be ran twice to update TOC)----%
\usepackage{titletoc}

%--section toc format--%
\titlecontents{section}
[6.6in]% Indentation of section name
{}
{\contentslabel{2.3em}}
{}
{\titlerule*[0.7pc]{.}\contentspage}

%--subsecion toc format--%
\titlecontents{subsection}
[6.9in]% Indentation of subsection name
{\hangindent1em}
{}%{\contentslabel{2.3em}}
{}
{\titlerule*[0.7pc]{.}\contentspage}

%--Page number formatting--%
\makeatletter
\renewcommand{\contentspage}[1][\thecontentspage]{\hb@xt@\@pnumwidth{#1\hfil}\hspace*{-\@pnumwidth}}
\renewcommand{\@pnumwidth}{0em} %distance of page number from right margin
\makeatother

%--------------------------------------------------%


\setlength{\parskip}{0pt}


%--newcommand for sections & subsections--%
\newcommand\addsection[1]{\addtocounter{section}{1}\phantomsection\addcontentsline{toc}{section}{{\arabic{section}.\xspace } {#1}}}
\newcommand\addsubsection[1]{\addtocounter{subsection}{1}\phantomsection\addcontentsline{toc}{subsection}{{} {#1}}}
%-----------------------------------------%


%%%%%                                                                                                       %%%%%
%%%%%  ____                   _             _____                                                      _    %%%%%
%%%%% |  _ \                 (_)           |  __ \                                                    | |   %%%%%
%%%%% | |_) |   ___    __ _   _   _ __     | |  | |   ___     ___   _   _   _ __ ___     ___   _ __   | |_  %%%%%
%%%%% |  _ <   / _ \  / _` | | | | '_ \    | |  | |  / _ \   / __| | | | | | '_ ` _ \   / _ \ | '_ \  | __| %%%%%
%%%%% | |_) | |  __/ | (_| | | | | | | |   | |__| | | (_) | | (__  | |_| | | | | | | | |  __/ | | | | | |_  %%%%%
%%%%% |____/   \___|  \__, | |_| |_| |_|   |_____/   \___/   \___|  \__,_| |_| |_| |_|  \___| |_| |_|  \__| %%%%%
%%%%%                  __/ |                                                                                %%%%%
%%%%%                 |___/                                                                                 %%%%%

\begin{document}
\AddToShipoutPictureBG{\pflogoart}
\renewcommand{\headrulewidth}{0pt}
\pagestyle{style1}

\renewcommand{\headrulewidth}{0pt}
\newcolumntype{Y}{>{\centering\arraybackslash}X}



%-----Table of Contents-----%
\setlength\headheight{0pt}
\vspace*{\fill}
\setcounter{tocdepth}{2}
\setcounter{secnumdepth}{-2}
\startcontents[section]

\label{sec:firstpage}
\belowpdfbookmark{Table of Contents}{}


\begin{tabularx}{1\textwidth}{@{}l Y@{}}
\hspace{6.1in} %section indent minus parindent
&%

\definecolor{c3486c7}{RGB}{52,134,199}
\definecolor{cffffff}{RGB}{255,255,255}


\begin{tikzpicture}[y=0.80pt, x=0.80pt, yscale=-1.000000, xscale=1.000000, inner sep=0pt, outer sep=0pt]
\path[fill=c3486c7,line width=0.026pt] (84.8597,0.0874) .. controls
  (85.3842,-0.0525) and (85.9786,0.0175) .. (86.5381,0.0175) .. controls
  (124.3368,0.0175) and (162.1006,0.0175) .. (199.8643,0.0175) .. controls
  (210.4591,-0.0175) and (220.7743,5.9618) .. (226.1591,15.0531) .. controls
  (245.0060,47.9565) and (263.8879,80.8249) .. (282.7348,113.6934) .. controls
  (287.8049,122.6098) and (288.0846,133.9739) .. (283.4691,143.1351) .. controls
  (281.5110,146.8416) and (279.3081,150.4081) .. (277.2800,154.0447) .. controls
  (260.2513,183.6612) and (243.2926,213.3128) .. (226.2640,242.8944) .. controls
  (221.8232,250.5520) and (213.9558,256.1117) .. (205.2492,257.6852) .. controls
  (201.6476,258.4195) and (197.9412,258.2097) .. (194.2697,258.2446) .. controls
  (158.4990,258.2446) and (122.6934,258.2096) .. (86.9227,258.2446) .. controls
  (81.2232,258.1397) and (75.4887,256.4963) .. (70.6633,253.4193) .. controls
  (66.2576,250.6569) and (62.6211,246.7406) .. (60.1035,242.2300) .. controls
  (42.8300,212.0889) and (25.5217,181.9828) .. (8.2482,151.8417) .. controls
  (5.8006,147.3311) and (2.7235,143.1001) .. (1.3249,138.0999) .. controls
  (-0.9130,130.7919) and (-0.2836,122.7147) .. (2.9683,115.7913) .. controls
  (5.8705,110.1617) and (9.2273,104.7769) .. (12.3043,99.2522) .. controls
  (28.0742,71.7686) and (43.8441,44.2850) .. (59.6139,16.8363) .. controls
  (61.0825,14.1089) and (62.8658,11.5214) .. (65.0687,9.3535) .. controls
  (70.2437,3.9687) and (77.4468,0.6119) .. (84.8597,0.0874) -- cycle;\path[color=black,fill=cffffff,line join=miter,line cap=butt,miter
  limit=4.00,fill opacity=0.343,nonzero rule,line width=0.196pt]
  (198.9701,4.7883) -- (198.9701,6.5588) -- (191.5148,6.5588) --
  (191.5148,4.8778) .. controls (191.4332,4.8787) and (191.3517,4.8808) ..
  (191.2701,4.8808) -- (191.2701,6.5584) -- (183.8148,6.5584) --
  (183.8148,4.9396) .. controls (183.7332,4.9401) and (183.6517,4.9405) ..
  (183.5701,4.9406) -- (183.5701,6.5580) -- (176.1148,6.5580) --
  (176.1148,4.9659) .. controls (176.0332,4.9661) and (175.9517,4.9665) ..
  (175.8701,4.9666) -- (175.8701,6.5580) -- (168.4148,6.5580) --
  (168.4148,4.9659) .. controls (168.3332,4.9658) and (168.2517,4.9653) ..
  (168.1701,4.9652) -- (168.1701,6.5580) -- (160.7148,6.5580) --
  (160.7148,4.9447) .. controls (160.6332,4.9445) and (160.5517,4.9443) ..
  (160.4701,4.9440) -- (160.4701,6.5580) -- (153.0148,6.5580) --
  (153.0148,4.9201) .. controls (152.9332,4.9198) and (152.8517,4.9191) ..
  (152.7701,4.9191) -- (152.7701,6.5584) -- (145.3148,6.5584) --
  (145.3148,4.8877) .. controls (145.2332,4.8873) and (145.1517,4.8867) ..
  (145.0701,4.8867) -- (145.0701,6.5587) -- (137.6148,6.5587) --
  (137.6148,4.8614) .. controls (137.5332,4.8611) and (137.4517,4.8610) ..
  (137.3701,4.8607) -- (137.3701,6.5587) -- (129.9148,6.5587) --
  (129.9148,4.8436) .. controls (129.8332,4.8435) and (129.7517,4.8430) ..
  (129.6701,4.8429) -- (129.6701,6.5587) -- (122.2148,6.5587) --
  (122.2148,4.8409) .. controls (122.1332,4.8410) and (122.0517,4.8414) ..
  (121.9701,4.8416) -- (121.9701,6.5587) -- (114.5148,6.5587) --
  (114.5148,4.8682) .. controls (114.4332,4.8687) and (114.3517,4.8691) ..
  (114.2701,4.8692) -- (114.2701,6.5584) -- (106.8148,6.5584) --
  (106.8148,4.9294) .. controls (106.7332,4.9302) and (106.6517,4.9314) ..
  (106.5701,4.9324) -- (106.5701,6.5586) -- (99.1148,6.5586) -- (99.1148,5.0322)
  .. controls (99.0332,5.0332) and (98.9517,5.0342) .. (98.8701,5.0352) --
  (98.8701,6.5582) -- (91.4148,6.5582) -- (91.4148,5.1828) .. controls
  (91.3332,5.1848) and (91.2517,5.1868) .. (91.1701,5.1878) -- (91.1701,6.5578)
  -- (83.7148,6.5578) -- (83.7148,5.3902) .. controls (83.6332,5.3932) and
  (83.5517,5.3952) .. (83.4701,5.3982) -- (83.4701,6.5582) -- (77.8981,6.5582)
  .. controls (77.6957,6.6359) and (77.4952,6.7187) .. (77.2952,6.8030) --
  (83.4701,6.8030) -- (83.4701,14.2582) -- (76.0148,14.2582) -- (76.0148,7.3915)
  .. controls (75.9336,7.4322) and (75.8509,7.4694) .. (75.7701,7.5112) --
  (75.7701,14.2582) -- (68.3148,14.2582) -- (68.3148,13.1399) .. controls
  (68.2381,13.2170) and (68.1589,13.2928) .. (68.0831,13.3702) .. controls
  (68.0786,13.3742) and (68.0746,13.3792) .. (68.0701,13.3839) --
  (68.0701,14.2583) -- (67.2416,14.2583) .. controls (67.1681,14.3394) and
  (67.0958,14.4211) .. (67.0235,14.5030) -- (68.0701,14.5030) --
  (68.0701,21.9583) -- (62.0859,21.9583) .. controls (62.0424,22.0399) and
  (61.9986,22.1212) .. (61.9553,22.2030) -- (68.0701,22.2030) --
  (68.0701,29.6583) -- (60.6148,29.6583) -- (60.6148,24.8075) .. controls
  (60.5336,24.9681) and (60.4512,25.1282) .. (60.3701,25.2887) --
  (60.3701,29.6583) -- (57.9556,29.6583) .. controls (57.9094,29.7399) and
  (57.8631,29.8214) .. (57.8169,29.9030) -- (60.3701,29.9030) --
  (60.3701,37.3583) -- (53.5676,37.3583) .. controls (53.5209,37.4398) and
  (53.4749,37.5215) .. (53.4282,37.6030) -- (60.3701,37.6030) --
  (60.3701,45.0583) -- (52.9148,45.0583) -- (52.9148,38.4985) .. controls
  (52.8334,38.6403) and (52.7516,38.7819) .. (52.6701,38.9237) --
  (52.6701,45.0583) -- (49.1339,45.0583) .. controls (49.0867,45.1399) and
  (49.0395,45.2214) .. (48.9924,45.3030) -- (52.6701,45.3030) --
  (52.6701,52.7583) -- (45.2148,52.7583) -- (45.2148,51.8183) .. controls
  (45.1333,51.9587) and (45.0516,52.0990) .. (44.9701,52.2394) --
  (44.9701,52.7583) -- (44.6686,52.7583) .. controls (44.6212,52.8399) and
  (44.5738,52.9214) .. (44.5264,53.0030) -- (44.9701,53.0030) --
  (44.9701,60.4583) -- (40.1918,60.4583) .. controls (40.1443,60.5398) and
  (40.0970,60.6214) .. (40.0496,60.7030) -- (44.9701,60.7030) --
  (44.9701,68.1583) -- (37.5148,68.1583) -- (37.5148,65.0677) .. controls
  (37.4331,65.2083) and (37.3518,65.3490) .. (37.2701,65.4895) --
  (37.2701,68.1583) -- (35.7217,68.1583) .. controls (35.6744,68.2399) and
  (35.6269,68.3214) .. (35.5796,68.4030) -- (37.2701,68.4030) --
  (37.2701,75.8583) -- (31.2606,75.8583) .. controls (31.2135,75.9399) and
  (31.1662,76.0214) .. (31.1191,76.1030) -- (37.2701,76.1030) --
  (37.2701,83.5583) -- (29.8148,83.5583) -- (29.8148,78.3657) .. controls
  (29.7333,78.5072) and (29.6515,78.6486) .. (29.5701,78.7902) --
  (29.5701,83.5583) -- (26.8364,83.5583) .. controls (26.7897,83.6399) and
  (26.7429,83.7214) .. (26.6963,83.8030) -- (29.5701,83.8030) --
  (29.5701,91.2583) -- (22.4580,91.2583) .. controls (22.4119,91.3400) and
  (22.3653,91.4213) .. (22.3192,91.5030) -- (29.5701,91.5030) --
  (29.5701,98.9583) -- (22.1148,98.9583) -- (22.1148,91.8667) .. controls
  (22.0329,92.0119) and (21.9519,92.1575) .. (21.8701,92.3028) --
  (21.8701,98.9583) -- (18.1493,98.9583) .. controls (18.1040,99.0398) and
  (18.0592,99.1215) .. (18.0139,99.2030) -- (21.8701,99.2030) --
  (21.8701,106.6583) -- (14.4148,106.6583) -- (14.4148,105.7450) .. controls
  (14.3332,105.8946) and (14.2517,106.0444) .. (14.1701,106.1941) --
  (14.1701,106.6583) -- (13.9178,106.6583) .. controls (13.8734,106.7399) and
  (13.8289,106.8214) .. (13.7845,106.9030) -- (14.1701,106.9030) --
  (14.1701,114.3583) -- (9.7739,114.3583) .. controls (9.7305,114.4400) and
  (9.6867,114.5213) .. (9.6433,114.6030) -- (14.1701,114.6030) --
  (14.1701,122.0583) -- (6.7148,122.0583) -- (6.7148,120.1756) .. controls
  (6.6336,120.3317) and (6.5512,120.4871) .. (6.4701,120.6432) --
  (6.4701,122.0583) -- (5.8063,122.0583) .. controls (5.7836,122.1398) and
  (5.7618,122.2214) .. (5.7400,122.3030) -- (6.4701,122.3030) --
  (6.4701,129.7583) -- (4.9457,129.7583) .. controls (4.9491,129.8399) and
  (4.9545,129.9214) .. (4.9587,130.0030) -- (6.4701,130.0030) --
  (6.4701,137.4583) -- (6.3272,137.4583) .. controls (6.3526,137.5399) and
  (6.3771,137.6216) .. (6.4031,137.7030) -- (6.4701,137.7030) --
  (6.4701,137.9094) .. controls (6.5495,138.1524) and (6.6300,138.3947) ..
  (6.7148,138.6361) -- (6.7148,137.7030) -- (14.1701,137.7030) --
  (14.1701,145.1583) -- (9.7397,145.1583) .. controls (9.7872,145.2402) and
  (9.8353,145.3215) .. (9.8833,145.4030) -- (14.1701,145.4030) --
  (14.1701,152.2608) .. controls (14.2516,152.4003) and (14.3334,152.5396) ..
  (14.4148,152.6792) -- (14.4148,145.4030) -- (21.8701,145.4030) --
  (21.8701,152.8583) -- (14.5194,152.8583) .. controls (14.5670,152.9399) and
  (14.6147,153.0214) .. (14.6623,153.1030) -- (21.8701,153.1030) --
  (21.8701,160.5583) -- (18.9853,160.5583) .. controls (19.0324,160.6400) and
  (19.0798,160.7213) .. (19.1268,160.8030) -- (21.8701,160.8030) --
  (21.8701,165.5854) .. controls (21.9516,165.7279) and (22.0333,165.8702) ..
  (22.1148,166.0127) -- (22.1148,160.8030) -- (29.5701,160.8030) --
  (29.5701,168.2583) -- (23.3979,168.2583) .. controls (23.4444,168.3398) and
  (23.4908,168.4214) .. (23.5374,168.5030) -- (29.5701,168.5030) --
  (29.5701,175.9583) -- (27.7763,175.9583) .. controls (27.8226,176.0398) and
  (27.8682,176.1215) .. (27.9144,176.2030) -- (29.5701,176.2030) --
  (29.5701,179.1240) .. controls (29.6517,179.2680) and (29.7332,179.4120) ..
  (29.8148,179.5560) -- (29.8148,176.2030) -- (37.2701,176.2030) --
  (37.2701,183.6583) -- (32.1397,183.6583) .. controls (32.1859,183.7399) and
  (32.2323,183.8214) .. (32.2785,183.9030) -- (37.2701,183.9030) --
  (37.2701,191.3583) -- (36.5051,191.3583) .. controls (36.5514,191.4399) and
  (36.5977,191.5214) .. (36.6439,191.6030) -- (37.2701,191.6030) --
  (37.2701,192.7070) .. controls (37.3517,192.8510) and (37.4331,192.9951) ..
  (37.5148,193.1390) -- (37.5148,191.6030) -- (44.9701,191.6030) --
  (44.9701,199.0583) -- (40.8767,199.0583) .. controls (40.9232,199.1399) and
  (40.9697,199.2214) .. (41.0162,199.3030) -- (44.9701,199.3030) --
  (44.9701,206.2141) .. controls (45.0517,206.3563) and (45.1332,206.4986) ..
  (45.2148,206.6407) -- (45.2148,199.3030) -- (52.6701,199.3030) --
  (52.6701,206.7583) -- (45.2825,206.7583) .. controls (45.3294,206.8399) and
  (45.3764,206.9214) .. (45.4233,207.0030) -- (52.6701,207.0030) --
  (52.6701,214.4583) -- (49.7313,214.4583) .. controls (49.7788,214.5399) and
  (49.8260,214.6214) .. (49.8735,214.7030) -- (52.6701,214.7030) --
  (52.6701,219.4820) .. controls (52.7518,219.6211) and (52.8330,219.7606) ..
  (52.9148,219.8997) -- (52.9148,214.7030) -- (60.3701,214.7030) --
  (60.3701,222.1583) -- (54.2444,222.1583) .. controls (54.2926,222.2399) and
  (54.3411,222.3214) .. (54.3893,222.4030) -- (60.3701,222.4030) --
  (60.3701,229.8583) -- (58.8423,229.8583) .. controls (58.8914,229.9399) and
  (58.9407,230.0214) .. (58.9899,230.1030) -- (60.3701,230.1030) --
  (60.3701,232.3896) .. controls (60.4520,232.5246) and (60.5329,232.6601) ..
  (60.6148,232.7950) -- (60.6148,230.1030) -- (68.0701,230.1030) --
  (68.0701,237.5583) -- (63.5338,237.5583) .. controls (63.5840,237.6398) and
  (63.6339,237.7215) .. (63.6841,237.8030) -- (68.0701,237.8030) --
  (68.0701,244.8358) .. controls (68.1522,244.9658) and (68.2326,245.0969) ..
  (68.3148,245.2268) -- (68.3148,237.8030) -- (75.7701,237.8030) --
  (75.7701,245.2583) -- (68.3346,245.2583) .. controls (68.3535,245.2881) and
  (68.3718,245.3181) .. (68.3907,245.3479) .. controls (68.4418,245.4008) and
  (68.4951,245.4507) .. (68.5466,245.5030) -- (75.7701,245.5030) --
  (75.7701,250.5917) .. controls (75.8513,250.6302) and (75.9333,250.6668) ..
  (76.0148,250.7045) -- (76.0148,245.5030) -- (83.4701,245.5030) --
  (83.4701,252.9583) -- (83.0975,252.9583) .. controls (83.5768,253.0507) and
  (84.0587,253.1310) .. (84.5420,253.2030) -- (91.1701,253.2030) --
  (91.1701,253.5831) .. controls (91.2518,253.5801) and (91.3332,253.5761) ..
  (91.4148,253.5731) -- (91.4148,253.2026) -- (98.8701,253.2026) --
  (98.8701,253.2826) .. controls (98.9517,253.2820) and (99.0332,253.2816) ..
  (99.1148,253.2806) -- (99.1148,253.2027) -- (106.5701,253.2027) --
  (106.5701,253.2430) .. controls (106.6517,253.2428) and (106.7332,253.2425) ..
  (106.8148,253.2423) -- (106.8148,253.2027) -- (114.2701,253.2027) --
  (114.2701,253.2396) .. controls (114.3517,253.2398) and (114.4332,253.2395) ..
  (114.5148,253.2396) -- (114.5148,253.2027) -- (121.9701,253.2027) --
  (121.9701,253.2622) .. controls (122.0517,253.2626) and (122.1332,253.2632) ..
  (122.2148,253.2632) -- (122.2148,253.2024) -- (129.6701,253.2024) --
  (129.6701,253.3050) .. controls (129.7517,253.3055) and (129.8332,253.3059) ..
  (129.9148,253.3060) -- (129.9148,253.2021) -- (137.3701,253.2021) --
  (137.3701,253.3490) .. controls (137.4517,253.3496) and (137.5332,253.3499) ..
  (137.6148,253.3500) -- (137.6148,253.2017) -- (145.0701,253.2017) --
  (145.0701,253.3911) .. controls (145.1517,253.3915) and (145.2333,253.3920) ..
  (145.3148,253.3921) -- (145.3148,253.2013) -- (152.7701,253.2013) --
  (152.7701,253.4262) .. controls (152.8517,253.4265) and (152.9333,253.4267) ..
  (153.0148,253.4269) -- (153.0148,253.2013) -- (160.4701,253.2013) --
  (160.4701,253.4372) .. controls (160.5517,253.4372) and (160.6333,253.4372) ..
  (160.7148,253.4372) -- (160.7148,253.2013) -- (168.1701,253.2013) --
  (168.1701,253.4174) .. controls (168.2517,253.4170) and (168.3332,253.4171) ..
  (168.4148,253.4167) -- (168.4148,253.2013) -- (175.8701,253.2013) --
  (175.8701,253.3586) .. controls (175.9517,253.3577) and (176.0333,253.3576) ..
  (176.1148,253.3566) -- (176.1148,253.2014) -- (183.5701,253.2014) --
  (183.5701,253.2499) .. controls (183.6517,253.2489) and (183.7333,253.2479) ..
  (183.8148,253.2469) -- (183.8148,253.2018) -- (186.1582,253.2018) .. controls
  (189.2754,253.1380) and (192.3922,253.0545) .. (195.5084,252.9571) --
  (191.5148,252.9571) -- (191.5148,245.5018) -- (198.9701,245.5018) --
  (198.9701,252.8484) .. controls (199.0517,252.8454) and (199.1332,252.8414) ..
  (199.2148,252.8384) -- (199.2148,245.5014) -- (206.6701,245.5014) --
  (206.6701,252.3681) .. controls (206.7516,252.3449) and (206.8337,252.3231) ..
  (206.9148,252.2990) -- (206.9148,245.5013) -- (214.3701,245.5013) --
  (214.3701,248.5973) .. controls (214.4524,248.5380) and (214.5332,248.4770) ..
  (214.6148,248.4168) -- (214.6148,245.5013) -- (217.9952,245.5013) .. controls
  (218.0767,245.4200) and (218.1586,245.3388) .. (218.2392,245.2566) --
  (214.6148,245.2566) -- (214.6148,237.8013) -- (222.0701,237.8013) --
  (222.0701,240.6068) .. controls (222.1522,240.4889) and (222.2340,240.3708) ..
  (222.3148,240.2520) -- (222.3148,237.8013) -- (223.8611,237.8013) .. controls
  (223.9085,237.7200) and (223.9550,237.6382) .. (224.0019,237.5566) --
  (222.3148,237.5566) -- (222.3148,230.1013) -- (227.9032,230.1013) .. controls
  (227.9493,230.0197) and (227.9958,229.9382) .. (228.0420,229.8566) --
  (222.3148,229.8566) -- (222.3148,222.4013) -- (229.7701,222.4013) --
  (229.7701,226.8002) .. controls (229.8514,226.6572) and (229.9335,226.5146) ..
  (230.0148,226.3716) -- (230.0148,222.4013) -- (232.2823,222.4013) .. controls
  (232.3290,222.3197) and (232.3757,222.2382) .. (232.4224,222.1566) --
  (230.0148,222.1566) -- (230.0148,214.7013) -- (236.7079,214.7013) .. controls
  (236.7550,214.6197) and (236.8023,214.5382) .. (236.8494,214.4566) --
  (230.0148,214.4566) -- (230.0148,207.0013) -- (237.4701,207.0013) --
  (237.4701,213.3840) .. controls (237.5516,213.2431) and (237.6333,213.1024) ..
  (237.7148,212.9616) -- (237.7148,207.0013) -- (241.1697,207.0013) .. controls
  (241.2171,206.9197) and (241.2645,206.8382) .. (241.3119,206.7566) --
  (237.7148,206.7566) -- (237.7148,199.3013) -- (245.1701,199.3013) --
  (245.1701,200.1202) .. controls (245.2517,199.9799) and (245.3332,199.8395) ..
  (245.4148,199.6992) -- (245.4148,199.3013) -- (245.6459,199.3013) .. controls
  (245.6933,199.2198) and (245.7406,199.1381) .. (245.7881,199.0566) --
  (245.4148,199.0566) -- (245.4148,191.6013) -- (250.1179,191.6013) .. controls
  (250.1653,191.5197) and (250.2127,191.4382) .. (250.2601,191.3566) --
  (245.4148,191.3566) -- (245.4148,183.9013) -- (252.8701,183.9013) --
  (252.8701,186.8606) .. controls (252.9515,186.7201) and (253.0335,186.5800) ..
  (253.1148,186.4395) -- (253.1148,183.9013) -- (254.5846,183.9013) .. controls
  (254.6317,183.8197) and (254.6790,183.7382) .. (254.7261,183.6566) --
  (253.1148,183.6566) -- (253.1148,176.2013) -- (259.0149,176.2013) .. controls
  (259.0616,176.1197) and (259.1083,176.0381) .. (259.1551,175.9566) --
  (253.1148,175.9566) -- (253.1148,168.5013) -- (260.5701,168.5013) --
  (260.5701,173.4861) .. controls (260.6515,173.3435) and (260.7335,173.2014) ..
  (260.8148,173.0588) -- (260.8148,168.5013) -- (263.4022,168.5013) .. controls
  (263.4484,168.4197) and (263.4948,168.3382) .. (263.5410,168.2566) --
  (260.8148,168.2566) -- (260.8148,160.8013) -- (267.7321,160.8013) .. controls
  (267.7775,160.7199) and (267.8221,160.6380) .. (267.8675,160.5566) --
  (260.8148,160.5566) -- (260.8148,153.1013) -- (268.2701,153.1013) --
  (268.2701,159.8292) .. controls (268.3520,159.6820) and (268.4330,159.5342) ..
  (268.5148,159.3869) -- (268.5148,153.1013) -- (271.9786,153.1013) .. controls
  (272.0232,153.0198) and (272.0674,152.9381) .. (272.1119,152.8566) --
  (268.5148,152.8566) -- (268.5148,145.4013) -- (275.9701,145.4013) --
  (275.9701,145.7164) .. controls (276.0699,145.5298) and (276.1698,145.3433) ..
  (276.2695,145.1566) -- (276.2148,145.1566) -- (276.2148,137.7013) --
  (280.2008,137.7013) .. controls (280.2433,137.6195) and (280.2868,137.5383) ..
  (280.3294,137.4566) -- (276.2148,137.4566) -- (276.2148,130.0013) --
  (281.8217,130.0013) .. controls (281.8227,129.9197) and (281.8218,129.8381) ..
  (281.8223,129.7566) -- (276.2148,129.7566) -- (276.2148,122.3013) --
  (280.8605,122.3013) .. controls (280.8393,122.2196) and (280.8180,122.1381) ..
  (280.7962,122.0566) -- (276.2148,122.0566) -- (276.2148,114.6013) --
  (277.8793,114.6013) .. controls (277.8371,114.5195) and (277.7942,114.4381) ..
  (277.7515,114.3566) -- (276.2148,114.3566) -- (276.2148,111.6584) .. controls
  (276.1341,111.5278) and (276.0520,111.3984) .. (275.9700,111.2688) --
  (275.9700,114.3566) -- (268.5148,114.3566) -- (268.5148,106.9013) --
  (273.1502,106.9013) .. controls (273.1027,106.8198) and (273.0555,106.7381) ..
  (273.0080,106.6566) -- (268.5148,106.6566) -- (268.5148,99.2013) --
  (268.6870,99.2013) .. controls (268.6399,99.1197) and (268.5925,99.0382) ..
  (268.5455,98.9566) -- (268.5147,98.9566) -- (268.5147,98.9033) .. controls
  (268.4332,98.7617) and (268.3515,98.6203) .. (268.2700,98.4788) --
  (268.2700,98.9566) -- (260.8147,98.9566) -- (260.8147,91.5013) --
  (264.2669,91.5013) .. controls (264.2203,91.4197) and (264.1734,91.3382) ..
  (264.1267,91.2566) -- (260.8147,91.2566) -- (260.8147,85.4460) .. controls
  (260.7332,85.3027) and (260.6515,85.1594) .. (260.5700,85.0161) --
  (260.5700,91.2566) -- (253.1147,91.2566) -- (253.1147,83.8013) --
  (259.8803,83.8013) .. controls (259.8339,83.7197) and (259.7872,83.6382) ..
  (259.7408,83.5566) -- (253.1147,83.5566) -- (253.1147,76.1013) --
  (255.5121,76.1013) .. controls (255.4659,76.0197) and (255.4196,75.9382) ..
  (255.3733,75.8566) -- (253.1147,75.8566) -- (253.1147,71.8760) .. controls
  (253.0332,71.7322) and (252.9515,71.5885) .. (252.8700,71.4447) --
  (252.8700,75.8566) -- (245.4147,75.8566) -- (245.4147,68.4013) --
  (251.1419,68.4013) .. controls (251.0956,68.3197) and (251.0494,68.2382) ..
  (251.0031,68.1566) -- (245.4147,68.1566) -- (245.4147,60.7013) --
  (246.7703,60.7013) .. controls (246.7238,60.6197) and (246.6773,60.5382) ..
  (246.6308,60.4566) -- (245.4147,60.4566) -- (245.4147,58.3265) .. controls
  (245.3332,58.1836) and (245.2516,58.0408) .. (245.1700,57.8979) --
  (245.1700,60.4566) -- (237.7147,60.4566) -- (237.7147,53.0013) --
  (242.3686,53.0013) .. controls (242.3218,52.9197) and (242.2746,52.8382) ..
  (242.2278,52.7566) -- (237.7147,52.7566) -- (237.7147,45.3013) --
  (237.9287,45.3013) .. controls (237.8815,45.2198) and (237.8345,45.1381) ..
  (237.7872,45.0566) -- (237.7147,45.0566) -- (237.7147,44.9315) .. controls
  (237.6331,44.7907) and (237.5516,44.6498) .. (237.4700,44.5090) --
  (237.4700,45.0566) -- (230.0147,45.0566) -- (230.0147,37.6013) --
  (233.4395,37.6013) .. controls (233.3917,37.5198) and (233.3445,37.4380) ..
  (233.2966,37.3566) -- (230.0147,37.3566) -- (230.0147,31.8058) .. controls
  (229.9328,31.6681) and (229.8519,31.5299) .. (229.7700,31.3922) --
  (229.7700,37.3566) -- (222.3147,37.3566) -- (222.3147,29.9013) --
  (228.8799,29.9013) .. controls (228.8311,29.8196) and (228.7817,29.7383) ..
  (228.7330,29.6566) -- (222.3147,29.6566) -- (222.3147,22.2013) --
  (224.2335,22.2013) .. controls (224.1838,22.1197) and (224.1342,22.0381) ..
  (224.0845,21.9566) -- (222.3147,21.9566) -- (222.3147,19.0534) .. controls
  (222.2336,18.9214) and (222.1511,18.7902) .. (222.0700,18.6582) --
  (222.0700,21.9566) -- (214.6147,21.9566) -- (214.6147,14.5013) --
  (219.4928,14.5013) .. controls (219.4420,14.4198) and (219.3918,14.3380) ..
  (219.3411,14.2566) -- (214.6147,14.2566) -- (214.6147,9.6779) .. controls
  (214.5332,9.6213) and (214.4521,9.5640) .. (214.3700,9.5083) --
  (214.3700,14.2566) -- (206.9147,14.2566) -- (206.9147,6.8013) --
  (209.2423,6.8013) .. controls (209.0312,6.7171) and (208.8190,6.6356) ..
  (208.6059,6.5566) -- (206.9147,6.5566) -- (206.9147,5.9885) .. controls
  (206.8335,5.9641) and (206.7514,5.9424) .. (206.6700,5.9188) --
  (206.6700,6.5566) -- (199.2147,6.5566) -- (199.2147,4.7875) .. controls
  (199.1331,4.7865) and (199.0516,4.7867) .. (198.9700,4.7865) --
  cycle(83.7148,6.8035) -- (91.1701,6.8035) -- (91.1701,14.2588) --
  (83.7148,14.2588) -- cycle(91.4148,6.8035) -- (98.8701,6.8035) --
  (98.8701,14.2588) -- (91.4148,14.2588) -- cycle(99.1148,6.8035) --
  (106.5701,6.8035) -- (106.5701,14.2588) -- (99.1148,14.2588) --
  cycle(106.8148,6.8035) -- (114.2701,6.8035) -- (114.2701,14.2588) --
  (106.8148,14.2588) -- cycle(114.5148,6.8035) -- (121.9701,6.8035) --
  (121.9701,14.2588) -- (114.5148,14.2588) -- cycle(122.2148,6.8035) --
  (129.6701,6.8035) -- (129.6701,14.2588) -- (122.2148,14.2588) --
  cycle(129.9148,6.8035) -- (137.3701,6.8035) -- (137.3701,14.2588) --
  (129.9148,14.2588) -- cycle(137.6148,6.8035) -- (145.0701,6.8035) --
  (145.0701,14.2588) -- (137.6148,14.2588) -- cycle(145.3148,6.8035) --
  (152.7701,6.8035) -- (152.7701,14.2588) -- (145.3148,14.2588) --
  cycle(153.0148,6.8035) -- (160.4701,6.8035) -- (160.4701,14.2588) --
  (153.0148,14.2588) -- cycle(160.7148,6.8035) -- (168.1701,6.8035) --
  (168.1701,14.2588) -- (160.7148,14.2588) -- cycle(168.4148,6.8035) --
  (175.8701,6.8035) -- (175.8701,14.2588) -- (168.4148,14.2588) --
  cycle(176.1148,6.8035) -- (183.5701,6.8035) -- (183.5701,14.2588) --
  (176.1148,14.2588) -- cycle(183.8148,6.8035) -- (191.2701,6.8035) --
  (191.2701,14.2588) -- (183.8148,14.2588) -- cycle(191.5148,6.8035) --
  (198.9701,6.8035) -- (198.9701,14.2588) -- (191.5148,14.2588) --
  cycle(199.2148,6.8035) -- (206.6701,6.8035) -- (206.6701,14.2588) --
  (199.2148,14.2588) -- cycle(68.3148,14.5035) -- (75.7701,14.5035) --
  (75.7701,21.9588) -- (68.3148,21.9588) -- cycle(76.0148,14.5035) --
  (83.4701,14.5035) -- (83.4701,21.9588) -- (76.0148,21.9588) --
  cycle(83.7148,14.5035) -- (91.1701,14.5035) -- (91.1701,21.9588) --
  (83.7148,21.9588) -- cycle(91.4148,14.5035) -- (98.8701,14.5035) --
  (98.8701,21.9588) -- (91.4148,21.9588) -- cycle(99.1148,14.5035) --
  (106.5701,14.5035) -- (106.5701,21.9588) -- (99.1148,21.9588) --
  cycle(106.8148,14.5035) -- (114.2701,14.5035) -- (114.2701,21.9588) --
  (106.8148,21.9588) -- cycle(114.5148,14.5035) -- (121.9701,14.5035) --
  (121.9701,21.9588) -- (114.5148,21.9588) -- cycle(122.2148,14.5035) --
  (129.6701,14.5035) -- (129.6701,21.9588) -- (122.2148,21.9588) --
  cycle(129.9148,14.5035) -- (137.3701,14.5035) -- (137.3701,21.9588) --
  (129.9148,21.9588) -- cycle(137.6148,14.5035) -- (145.0701,14.5035) --
  (145.0701,21.9588) -- (137.6148,21.9588) -- cycle(145.3148,14.5035) --
  (152.7701,14.5035) -- (152.7701,21.9588) -- (145.3148,21.9588) --
  cycle(153.0148,14.5035) -- (160.4701,14.5035) -- (160.4701,21.9588) --
  (153.0148,21.9588) -- cycle(160.7148,14.5035) -- (168.1701,14.5035) --
  (168.1701,21.9588) -- (160.7148,21.9588) -- cycle(168.4148,14.5035) --
  (175.8701,14.5035) -- (175.8701,21.9588) -- (168.4148,21.9588) --
  cycle(176.1148,14.5035) -- (183.5701,14.5035) -- (183.5701,21.9588) --
  (176.1148,21.9588) -- cycle(183.8148,14.5035) -- (191.2701,14.5035) --
  (191.2701,21.9588) -- (183.8148,21.9588) -- cycle(191.5148,14.5035) --
  (198.9701,14.5035) -- (198.9701,21.9588) -- (191.5148,21.9588) --
  cycle(199.2148,14.5035) -- (206.6701,14.5035) -- (206.6701,21.9588) --
  (199.2148,21.9588) -- cycle(206.9148,14.5035) -- (214.3701,14.5035) --
  (214.3701,21.9588) -- (206.9148,21.9588) -- cycle(68.3148,22.2035) --
  (75.7701,22.2035) -- (75.7701,29.6588) -- (68.3148,29.6588) --
  cycle(76.0148,22.2035) -- (83.4701,22.2035) -- (83.4701,29.6588) --
  (76.0148,29.6588) -- cycle(83.7148,22.2035) -- (91.1701,22.2035) --
  (91.1701,29.6588) -- (83.7148,29.6588) -- cycle(91.4148,22.2035) --
  (98.8701,22.2035) -- (98.8701,29.6588) -- (91.4148,29.6588) --
  cycle(99.1148,22.2035) -- (106.5701,22.2035) -- (106.5701,29.6588) --
  (99.1148,29.6588) -- cycle(106.8148,22.2035) -- (114.2701,22.2035) --
  (114.2701,29.6588) -- (106.8148,29.6588) -- cycle(114.5148,22.2035) --
  (121.9701,22.2035) -- (121.9701,29.6588) -- (114.5148,29.6588) --
  cycle(122.2148,22.2035) -- (129.6701,22.2035) -- (129.6701,29.6588) --
  (122.2148,29.6588) -- cycle(129.9148,22.2035) -- (137.3701,22.2035) --
  (137.3701,29.6588) -- (129.9148,29.6588) -- cycle(137.6148,22.2035) --
  (145.0701,22.2035) -- (145.0701,29.6588) -- (137.6148,29.6588) --
  cycle(145.3148,22.2035) -- (152.7701,22.2035) -- (152.7701,29.6588) --
  (145.3148,29.6588) -- cycle(153.0148,22.2035) -- (160.4701,22.2035) --
  (160.4701,29.6588) -- (153.0148,29.6588) -- cycle(160.7148,22.2035) --
  (168.1701,22.2035) -- (168.1701,29.6588) -- (160.7148,29.6588) --
  cycle(168.4148,22.2035) -- (175.8701,22.2035) -- (175.8701,29.6588) --
  (168.4148,29.6588) -- cycle(176.1148,22.2035) -- (183.5701,22.2035) --
  (183.5701,29.6588) -- (176.1148,29.6588) -- cycle(183.8148,22.2035) --
  (191.2701,22.2035) -- (191.2701,29.6588) -- (183.8148,29.6588) --
  cycle(191.5148,22.2035) -- (198.9701,22.2035) -- (198.9701,29.6588) --
  (191.5148,29.6588) -- cycle(199.2148,22.2035) -- (206.6701,22.2035) --
  (206.6701,29.6588) -- (199.2148,29.6588) -- cycle(206.9148,22.2035) --
  (214.3701,22.2035) -- (214.3701,29.6588) -- (206.9148,29.6588) --
  cycle(214.6148,22.2035) -- (222.0701,22.2035) -- (222.0701,29.6588) --
  (214.6148,29.6588) -- cycle(60.6148,29.9035) -- (68.0701,29.9035) --
  (68.0701,37.3588) -- (60.6148,37.3588) -- cycle(68.3148,29.9035) --
  (75.7701,29.9035) -- (75.7701,37.3588) -- (68.3148,37.3588) --
  cycle(76.0148,29.9035) -- (83.4701,29.9035) -- (83.4701,37.3588) --
  (76.0148,37.3588) -- cycle(83.7148,29.9035) -- (91.1701,29.9035) --
  (91.1701,37.3588) -- (83.7148,37.3588) -- cycle(91.4148,29.9035) --
  (98.8701,29.9035) -- (98.8701,37.3588) -- (91.4148,37.3588) --
  cycle(99.1148,29.9035) -- (106.5701,29.9035) -- (106.5701,37.3588) --
  (99.1148,37.3588) -- cycle(106.8148,29.9035) -- (114.2701,29.9035) --
  (114.2701,37.3588) -- (106.8148,37.3588) -- cycle(114.5148,29.9035) --
  (121.9701,29.9035) -- (121.9701,37.3588) -- (114.5148,37.3588) --
  cycle(122.2148,29.9035) -- (129.6701,29.9035) -- (129.6701,37.3588) --
  (122.2148,37.3588) -- cycle(129.9148,29.9035) -- (137.3701,29.9035) --
  (137.3701,37.3588) -- (129.9148,37.3588) -- cycle(137.6148,29.9035) --
  (145.0701,29.9035) -- (145.0701,37.3588) -- (137.6148,37.3588) --
  cycle(145.3148,29.9035) -- (152.7701,29.9035) -- (152.7701,37.3588) --
  (145.3148,37.3588) -- cycle(153.0148,29.9035) -- (160.4701,29.9035) --
  (160.4701,37.3588) -- (153.0148,37.3588) -- cycle(160.7148,29.9035) --
  (168.1701,29.9035) -- (168.1701,37.3588) -- (160.7148,37.3588) --
  cycle(168.4148,29.9035) -- (175.8701,29.9035) -- (175.8701,37.3588) --
  (168.4148,37.3588) -- cycle(176.1148,29.9035) -- (183.5701,29.9035) --
  (183.5701,37.3588) -- (176.1148,37.3588) -- cycle(183.8148,29.9035) --
  (191.2701,29.9035) -- (191.2701,37.3588) -- (183.8148,37.3588) --
  cycle(191.5148,29.9035) -- (198.9701,29.9035) -- (198.9701,37.3588) --
  (191.5148,37.3588) -- cycle(199.2148,29.9035) -- (206.6701,29.9035) --
  (206.6701,37.3588) -- (199.2148,37.3588) -- cycle(206.9148,29.9035) --
  (214.3701,29.9035) -- (214.3701,37.3588) -- (206.9148,37.3588) --
  cycle(214.6148,29.9035) -- (222.0701,29.9035) -- (222.0701,37.3588) --
  (214.6148,37.3588) -- cycle(60.6148,37.6035) -- (68.0701,37.6035) --
  (68.0701,45.0588) -- (60.6148,45.0588) -- cycle(68.3148,37.6035) --
  (75.7701,37.6035) -- (75.7701,45.0588) -- (68.3148,45.0588) --
  cycle(76.0148,37.6035) -- (83.4701,37.6035) -- (83.4701,45.0588) --
  (76.0148,45.0588) -- cycle(83.7148,37.6035) -- (91.1701,37.6035) --
  (91.1701,45.0588) -- (83.7148,45.0588) -- cycle(91.4148,37.6035) --
  (98.8701,37.6035) -- (98.8701,45.0588) -- (91.4148,45.0588) --
  cycle(99.1148,37.6035) -- (106.5701,37.6035) -- (106.5701,45.0588) --
  (99.1148,45.0588) -- cycle(106.8148,37.6035) -- (114.2701,37.6035) --
  (114.2701,45.0588) -- (106.8148,45.0588) -- cycle(114.5148,37.6035) --
  (121.9701,37.6035) -- (121.9701,45.0588) -- (114.5148,45.0588) --
  cycle(122.2148,37.6035) -- (129.6701,37.6035) -- (129.6701,45.0588) --
  (122.2148,45.0588) -- cycle(129.9148,37.6035) -- (137.3701,37.6035) --
  (137.3701,45.0588) -- (129.9148,45.0588) -- cycle(137.6148,37.6035) --
  (145.0701,37.6035) -- (145.0701,45.0588) -- (137.6148,45.0588) --
  cycle(145.3148,37.6035) -- (152.7701,37.6035) -- (152.7701,45.0588) --
  (145.3148,45.0588) -- cycle(153.0148,37.6035) -- (160.4701,37.6035) --
  (160.4701,45.0588) -- (153.0148,45.0588) -- cycle(160.7148,37.6035) --
  (168.1701,37.6035) -- (168.1701,45.0588) -- (160.7148,45.0588) --
  cycle(168.4148,37.6035) -- (175.8701,37.6035) -- (175.8701,45.0588) --
  (168.4148,45.0588) -- cycle(176.1148,37.6035) -- (183.5701,37.6035) --
  (183.5701,45.0588) -- (176.1148,45.0588) -- cycle(183.8148,37.6035) --
  (191.2701,37.6035) -- (191.2701,45.0588) -- (183.8148,45.0588) --
  cycle(191.5148,37.6035) -- (198.9701,37.6035) -- (198.9701,45.0588) --
  (191.5148,45.0588) -- cycle(199.2148,37.6035) -- (206.6701,37.6035) --
  (206.6701,45.0588) -- (199.2148,45.0588) -- cycle(206.9148,37.6035) --
  (214.3701,37.6035) -- (214.3701,45.0588) -- (206.9148,45.0588) --
  cycle(214.6148,37.6035) -- (222.0701,37.6035) -- (222.0701,45.0588) --
  (214.6148,45.0588) -- cycle(222.3148,37.6035) -- (229.7701,37.6035) --
  (229.7701,45.0588) -- (222.3148,45.0588) -- cycle(52.9148,45.3035) --
  (60.3701,45.3035) -- (60.3701,52.7588) -- (52.9148,52.7588) --
  cycle(60.6148,45.3035) -- (68.0701,45.3035) -- (68.0701,52.7588) --
  (60.6148,52.7588) -- cycle(68.3148,45.3035) -- (75.7701,45.3035) --
  (75.7701,52.7588) -- (68.3148,52.7588) -- cycle(76.0148,45.3035) --
  (83.4701,45.3035) -- (83.4701,52.7588) -- (76.0148,52.7588) --
  cycle(83.7148,45.3035) -- (91.1701,45.3035) -- (91.1701,52.7588) --
  (83.7148,52.7588) -- cycle(91.4148,45.3035) -- (98.8701,45.3035) --
  (98.8701,52.7588) -- (91.4148,52.7588) -- cycle(99.1148,45.3035) --
  (106.5701,45.3035) -- (106.5701,52.7588) -- (99.1148,52.7588) --
  cycle(106.8148,45.3035) -- (114.2701,45.3035) -- (114.2701,52.7588) --
  (106.8148,52.7588) -- cycle(114.5148,45.3035) -- (121.9701,45.3035) --
  (121.9701,52.7588) -- (114.5148,52.7588) -- cycle(122.2148,45.3035) --
  (129.6701,45.3035) -- (129.6701,52.7588) -- (122.2148,52.7588) --
  cycle(129.9148,45.3035) -- (137.3701,45.3035) -- (137.3701,52.7588) --
  (129.9148,52.7588) -- cycle(137.6148,45.3035) -- (145.0701,45.3035) --
  (145.0701,52.7588) -- (137.6148,52.7588) -- cycle(145.3148,45.3035) --
  (152.7701,45.3035) -- (152.7701,52.7588) -- (145.3148,52.7588) --
  cycle(153.0148,45.3035) -- (160.4701,45.3035) -- (160.4701,52.7588) --
  (153.0148,52.7588) -- cycle(160.7148,45.3035) -- (168.1701,45.3035) --
  (168.1701,52.7588) -- (160.7148,52.7588) -- cycle(168.4148,45.3035) --
  (175.8701,45.3035) -- (175.8701,52.7588) -- (168.4148,52.7588) --
  cycle(176.1148,45.3035) -- (183.5701,45.3035) -- (183.5701,52.7588) --
  (176.1148,52.7588) -- cycle(183.8148,45.3035) -- (191.2701,45.3035) --
  (191.2701,52.7588) -- (183.8148,52.7588) -- cycle(191.5148,45.3035) --
  (198.9701,45.3035) -- (198.9701,52.7588) -- (191.5148,52.7588) --
  cycle(199.2148,45.3035) -- (206.6701,45.3035) -- (206.6701,52.7588) --
  (199.2148,52.7588) -- cycle(206.9148,45.3035) -- (214.3701,45.3035) --
  (214.3701,52.7588) -- (206.9148,52.7588) -- cycle(214.6148,45.3035) --
  (222.0701,45.3035) -- (222.0701,52.7588) -- (214.6148,52.7588) --
  cycle(222.3148,45.3035) -- (229.7701,45.3035) -- (229.7701,52.7588) --
  (222.3148,52.7588) -- cycle(230.0148,45.3035) -- (237.4701,45.3035) --
  (237.4701,52.7588) -- (230.0148,52.7588) -- cycle(45.2148,53.0035) --
  (52.6701,53.0035) -- (52.6701,60.4588) -- (45.2148,60.4588) --
  cycle(52.9148,53.0035) -- (60.3701,53.0035) -- (60.3701,60.4588) --
  (52.9148,60.4588) -- cycle(60.6148,53.0035) -- (68.0701,53.0035) --
  (68.0701,60.4588) -- (60.6148,60.4588) -- cycle(68.3148,53.0035) --
  (75.7701,53.0035) -- (75.7701,60.4588) -- (68.3148,60.4588) --
  cycle(76.0148,53.0035) -- (83.4701,53.0035) -- (83.4701,60.4588) --
  (76.0148,60.4588) -- cycle(83.7148,53.0035) -- (91.1701,53.0035) --
  (91.1701,60.4588) -- (83.7148,60.4588) -- cycle(91.4148,53.0035) --
  (98.8701,53.0035) -- (98.8701,60.4588) -- (91.4148,60.4588) --
  cycle(99.1148,53.0035) -- (106.5701,53.0035) -- (106.5701,60.4588) --
  (99.1148,60.4588) -- cycle(106.8148,53.0035) -- (114.2701,53.0035) --
  (114.2701,60.4588) -- (106.8148,60.4588) -- cycle(114.5148,53.0035) --
  (121.9701,53.0035) -- (121.9701,60.4588) -- (114.5148,60.4588) --
  cycle(122.2148,53.0035) -- (129.6701,53.0035) -- (129.6701,60.4588) --
  (122.2148,60.4588) -- cycle(129.9148,53.0035) -- (137.3701,53.0035) --
  (137.3701,60.4588) -- (129.9148,60.4588) -- cycle(137.6148,53.0035) --
  (145.0701,53.0035) -- (145.0701,60.4588) -- (137.6148,60.4588) --
  cycle(145.3148,53.0035) -- (152.7701,53.0035) -- (152.7701,60.4588) --
  (145.3148,60.4588) -- cycle(153.0148,53.0035) -- (160.4701,53.0035) --
  (160.4701,60.4588) -- (153.0148,60.4588) -- cycle(160.7148,53.0035) --
  (168.1701,53.0035) -- (168.1701,60.4588) -- (160.7148,60.4588) --
  cycle(168.4148,53.0035) -- (175.8701,53.0035) -- (175.8701,60.4588) --
  (168.4148,60.4588) -- cycle(176.1148,53.0035) -- (183.5701,53.0035) --
  (183.5701,60.4588) -- (176.1148,60.4588) -- cycle(183.8148,53.0035) --
  (191.2701,53.0035) -- (191.2701,60.4588) -- (183.8148,60.4588) --
  cycle(191.5148,53.0035) -- (198.9701,53.0035) -- (198.9701,60.4588) --
  (191.5148,60.4588) -- cycle(199.2148,53.0035) -- (206.6701,53.0035) --
  (206.6701,60.4588) -- (199.2148,60.4588) -- cycle(206.9148,53.0035) --
  (214.3701,53.0035) -- (214.3701,60.4588) -- (206.9148,60.4588) --
  cycle(214.6148,53.0035) -- (222.0701,53.0035) -- (222.0701,60.4588) --
  (214.6148,60.4588) -- cycle(222.3148,53.0035) -- (229.7701,53.0035) --
  (229.7701,60.4588) -- (222.3148,60.4588) -- cycle(230.0148,53.0035) --
  (237.4701,53.0035) -- (237.4701,60.4588) -- (230.0148,60.4588) --
  cycle(45.2148,60.7035) -- (52.6701,60.7035) -- (52.6701,68.1588) --
  (45.2148,68.1588) -- cycle(52.9148,60.7035) -- (60.3701,60.7035) --
  (60.3701,68.1588) -- (52.9148,68.1588) -- cycle(60.6148,60.7035) --
  (68.0701,60.7035) -- (68.0701,68.1588) -- (60.6148,68.1588) --
  cycle(68.3148,60.7035) -- (75.7701,60.7035) -- (75.7701,68.1588) --
  (68.3148,68.1588) -- cycle(76.0148,60.7035) -- (83.4701,60.7035) --
  (83.4701,68.1588) -- (76.0148,68.1588) -- cycle(83.7148,60.7035) --
  (91.1701,60.7035) -- (91.1701,68.1588) -- (83.7148,68.1588) --
  cycle(91.4148,60.7035) -- (98.8701,60.7035) -- (98.8701,68.1588) --
  (91.4148,68.1588) -- cycle(99.1148,60.7035) -- (106.5701,60.7035) --
  (106.5701,68.1588) -- (99.1148,68.1588) -- cycle(106.8148,60.7035) --
  (114.2701,60.7035) -- (114.2701,68.1588) -- (106.8148,68.1588) --
  cycle(114.5148,60.7035) -- (121.9701,60.7035) -- (121.9701,68.1588) --
  (114.5148,68.1588) -- cycle(122.2148,60.7035) -- (129.6701,60.7035) --
  (129.6701,68.1588) -- (122.2148,68.1588) -- cycle(129.9148,60.7035) --
  (137.3701,60.7035) -- (137.3701,68.1588) -- (129.9148,68.1588) --
  cycle(137.6148,60.7035) -- (145.0701,60.7035) -- (145.0701,68.1588) --
  (137.6148,68.1588) -- cycle(145.3148,60.7035) -- (152.7701,60.7035) --
  (152.7701,68.1588) -- (145.3148,68.1588) -- cycle(153.0148,60.7035) --
  (160.4701,60.7035) -- (160.4701,68.1588) -- (153.0148,68.1588) --
  cycle(160.7148,60.7035) -- (168.1701,60.7035) -- (168.1701,68.1588) --
  (160.7148,68.1588) -- cycle(168.4148,60.7035) -- (175.8701,60.7035) --
  (175.8701,68.1588) -- (168.4148,68.1588) -- cycle(176.1148,60.7035) --
  (183.5701,60.7035) -- (183.5701,68.1588) -- (176.1148,68.1588) --
  cycle(183.8148,60.7035) -- (191.2701,60.7035) -- (191.2701,68.1588) --
  (183.8148,68.1588) -- cycle(191.5148,60.7035) -- (198.9701,60.7035) --
  (198.9701,68.1588) -- (191.5148,68.1588) -- cycle(199.2148,60.7035) --
  (206.6701,60.7035) -- (206.6701,68.1588) -- (199.2148,68.1588) --
  cycle(206.9148,60.7035) -- (214.3701,60.7035) -- (214.3701,68.1588) --
  (206.9148,68.1588) -- cycle(214.6148,60.7035) -- (222.0701,60.7035) --
  (222.0701,68.1588) -- (214.6148,68.1588) -- cycle(222.3148,60.7035) --
  (229.7701,60.7035) -- (229.7701,68.1588) -- (222.3148,68.1588) --
  cycle(230.0148,60.7035) -- (237.4701,60.7035) -- (237.4701,68.1588) --
  (230.0148,68.1588) -- cycle(237.7148,60.7035) -- (245.1701,60.7035) --
  (245.1701,68.1588) -- (237.7148,68.1588) -- cycle(37.5148,68.4035) --
  (44.9701,68.4035) -- (44.9701,75.8588) -- (37.5148,75.8588) --
  cycle(45.2148,68.4035) -- (52.6701,68.4035) -- (52.6701,75.8588) --
  (45.2148,75.8588) -- cycle(52.9148,68.4035) -- (60.3701,68.4035) --
  (60.3701,75.8588) -- (52.9148,75.8588) -- cycle(60.6148,68.4035) --
  (68.0701,68.4035) -- (68.0701,75.8588) -- (60.6148,75.8588) --
  cycle(68.3148,68.4035) -- (75.7701,68.4035) -- (75.7701,75.8588) --
  (68.3148,75.8588) -- cycle(76.0148,68.4035) -- (83.4701,68.4035) --
  (83.4701,75.8588) -- (76.0148,75.8588) -- cycle(83.7148,68.4035) --
  (91.1701,68.4035) -- (91.1701,75.8588) -- (83.7148,75.8588) --
  cycle(91.4148,68.4035) -- (98.8701,68.4035) -- (98.8701,75.8588) --
  (91.4148,75.8588) -- cycle(99.1148,68.4035) -- (106.5701,68.4035) --
  (106.5701,75.8588) -- (99.1148,75.8588) -- cycle(106.8148,68.4035) --
  (114.2701,68.4035) -- (114.2701,75.8588) -- (106.8148,75.8588) --
  cycle(114.5148,68.4035) -- (121.9701,68.4035) -- (121.9701,75.8588) --
  (114.5148,75.8588) -- cycle(122.2148,68.4035) -- (129.6701,68.4035) --
  (129.6701,75.8588) -- (122.2148,75.8588) -- cycle(129.9148,68.4035) --
  (137.3701,68.4035) -- (137.3701,75.8588) -- (129.9148,75.8588) --
  cycle(137.6148,68.4035) -- (145.0701,68.4035) -- (145.0701,75.8588) --
  (137.6148,75.8588) -- cycle(145.3148,68.4035) -- (152.7701,68.4035) --
  (152.7701,75.8588) -- (145.3148,75.8588) -- cycle(153.0148,68.4035) --
  (160.4701,68.4035) -- (160.4701,75.8588) -- (153.0148,75.8588) --
  cycle(160.7148,68.4035) -- (168.1701,68.4035) -- (168.1701,75.8588) --
  (160.7148,75.8588) -- cycle(168.4148,68.4035) -- (175.8701,68.4035) --
  (175.8701,75.8588) -- (168.4148,75.8588) -- cycle(176.1148,68.4035) --
  (183.5701,68.4035) -- (183.5701,75.8588) -- (176.1148,75.8588) --
  cycle(183.8148,68.4035) -- (191.2701,68.4035) -- (191.2701,75.8588) --
  (183.8148,75.8588) -- cycle(191.5148,68.4035) -- (198.9701,68.4035) --
  (198.9701,75.8588) -- (191.5148,75.8588) -- cycle(199.2148,68.4035) --
  (206.6701,68.4035) -- (206.6701,75.8588) -- (199.2148,75.8588) --
  cycle(206.9148,68.4035) -- (214.3701,68.4035) -- (214.3701,75.8588) --
  (206.9148,75.8588) -- cycle(214.6148,68.4035) -- (222.0701,68.4035) --
  (222.0701,75.8588) -- (214.6148,75.8588) -- cycle(222.3148,68.4035) --
  (229.7701,68.4035) -- (229.7701,75.8588) -- (222.3148,75.8588) --
  cycle(230.0148,68.4035) -- (237.4701,68.4035) -- (237.4701,75.8588) --
  (230.0148,75.8588) -- cycle(237.7148,68.4035) -- (245.1701,68.4035) --
  (245.1701,75.8588) -- (237.7148,75.8588) -- cycle(37.5148,76.1035) --
  (44.9701,76.1035) -- (44.9701,83.5588) -- (37.5148,83.5588) --
  cycle(45.2148,76.1035) -- (52.6701,76.1035) -- (52.6701,83.5588) --
  (45.2148,83.5588) -- cycle(52.9148,76.1035) -- (60.3701,76.1035) --
  (60.3701,83.5588) -- (52.9148,83.5588) -- cycle(60.6148,76.1035) --
  (68.0701,76.1035) -- (68.0701,83.5588) -- (60.6148,83.5588) --
  cycle(68.3148,76.1035) -- (75.7701,76.1035) -- (75.7701,83.5588) --
  (68.3148,83.5588) -- cycle(76.0148,76.1035) -- (83.4701,76.1035) --
  (83.4701,83.5588) -- (76.0148,83.5588) -- cycle(83.7148,76.1035) --
  (91.1701,76.1035) -- (91.1701,83.5588) -- (83.7148,83.5588) --
  cycle(91.4148,76.1035) -- (98.8701,76.1035) -- (98.8701,83.5588) --
  (91.4148,83.5588) -- cycle(99.1148,76.1035) -- (106.5701,76.1035) --
  (106.5701,83.5588) -- (99.1148,83.5588) -- cycle(106.8148,76.1035) --
  (114.2701,76.1035) -- (114.2701,83.5588) -- (106.8148,83.5588) --
  cycle(114.5148,76.1035) -- (121.9701,76.1035) -- (121.9701,83.5588) --
  (114.5148,83.5588) -- cycle(122.2148,76.1035) -- (129.6701,76.1035) --
  (129.6701,83.5588) -- (122.2148,83.5588) -- cycle(129.9148,76.1035) --
  (137.3701,76.1035) -- (137.3701,83.5588) -- (129.9148,83.5588) --
  cycle(137.6148,76.1035) -- (145.0701,76.1035) -- (145.0701,83.5588) --
  (137.6148,83.5588) -- cycle(145.3148,76.1035) -- (152.7701,76.1035) --
  (152.7701,83.5588) -- (145.3148,83.5588) -- cycle(153.0148,76.1035) --
  (160.4701,76.1035) -- (160.4701,83.5588) -- (153.0148,83.5588) --
  cycle(160.7148,76.1035) -- (168.1701,76.1035) -- (168.1701,83.5588) --
  (160.7148,83.5588) -- cycle(168.4148,76.1035) -- (175.8701,76.1035) --
  (175.8701,83.5588) -- (168.4148,83.5588) -- cycle(176.1148,76.1035) --
  (183.5701,76.1035) -- (183.5701,83.5588) -- (176.1148,83.5588) --
  cycle(183.8148,76.1035) -- (191.2701,76.1035) -- (191.2701,83.5588) --
  (183.8148,83.5588) -- cycle(191.5148,76.1035) -- (198.9701,76.1035) --
  (198.9701,83.5588) -- (191.5148,83.5588) -- cycle(199.2148,76.1035) --
  (206.6701,76.1035) -- (206.6701,83.5588) -- (199.2148,83.5588) --
  cycle(206.9148,76.1035) -- (214.3701,76.1035) -- (214.3701,83.5588) --
  (206.9148,83.5588) -- cycle(214.6148,76.1035) -- (222.0701,76.1035) --
  (222.0701,83.5588) -- (214.6148,83.5588) -- cycle(222.3148,76.1035) --
  (229.7701,76.1035) -- (229.7701,83.5588) -- (222.3148,83.5588) --
  cycle(230.0148,76.1035) -- (237.4701,76.1035) -- (237.4701,83.5588) --
  (230.0148,83.5588) -- cycle(237.7148,76.1035) -- (245.1701,76.1035) --
  (245.1701,83.5588) -- (237.7148,83.5588) -- cycle(245.4148,76.1035) --
  (252.8701,76.1035) -- (252.8701,83.5588) -- (245.4148,83.5588) --
  cycle(29.8148,83.8035) -- (37.2701,83.8035) -- (37.2701,91.2588) --
  (29.8148,91.2588) -- cycle(37.5148,83.8035) -- (44.9701,83.8035) --
  (44.9701,91.2588) -- (37.5148,91.2588) -- cycle(45.2148,83.8035) --
  (52.6701,83.8035) -- (52.6701,91.2588) -- (45.2148,91.2588) --
  cycle(52.9148,83.8035) -- (60.3701,83.8035) -- (60.3701,91.2588) --
  (52.9148,91.2588) -- cycle(60.6148,83.8035) -- (68.0701,83.8035) --
  (68.0701,91.2588) -- (60.6148,91.2588) -- cycle(68.3148,83.8035) --
  (75.7701,83.8035) -- (75.7701,91.2588) -- (68.3148,91.2588) --
  cycle(76.0148,83.8035) -- (83.4701,83.8035) -- (83.4701,91.2588) --
  (76.0148,91.2588) -- cycle(83.7148,83.8035) -- (91.1701,83.8035) --
  (91.1701,91.2588) -- (83.7148,91.2588) -- cycle(91.4148,83.8035) --
  (98.8701,83.8035) -- (98.8701,91.2588) -- (91.4148,91.2588) --
  cycle(99.1148,83.8035) -- (106.5701,83.8035) -- (106.5701,91.2588) --
  (99.1148,91.2588) -- cycle(106.8148,83.8035) -- (114.2701,83.8035) --
  (114.2701,91.2588) -- (106.8148,91.2588) -- cycle(114.5148,83.8035) --
  (121.9701,83.8035) -- (121.9701,91.2588) -- (114.5148,91.2588) --
  cycle(122.2148,83.8035) -- (129.6701,83.8035) -- (129.6701,91.2588) --
  (122.2148,91.2588) -- cycle(129.9148,83.8035) -- (137.3701,83.8035) --
  (137.3701,91.2588) -- (129.9148,91.2588) -- cycle(137.6148,83.8035) --
  (145.0701,83.8035) -- (145.0701,91.2588) -- (137.6148,91.2588) --
  cycle(145.3148,83.8035) -- (152.7701,83.8035) -- (152.7701,91.2588) --
  (145.3148,91.2588) -- cycle(153.0148,83.8035) -- (160.4701,83.8035) --
  (160.4701,91.2588) -- (153.0148,91.2588) -- cycle(160.7148,83.8035) --
  (168.1701,83.8035) -- (168.1701,91.2588) -- (160.7148,91.2588) --
  cycle(168.4148,83.8035) -- (175.8701,83.8035) -- (175.8701,91.2588) --
  (168.4148,91.2588) -- cycle(176.1148,83.8035) -- (183.5701,83.8035) --
  (183.5701,91.2588) -- (176.1148,91.2588) -- cycle(183.8148,83.8035) --
  (191.2701,83.8035) -- (191.2701,91.2588) -- (183.8148,91.2588) --
  cycle(191.5148,83.8035) -- (198.9701,83.8035) -- (198.9701,91.2588) --
  (191.5148,91.2588) -- cycle(199.2148,83.8035) -- (206.6701,83.8035) --
  (206.6701,91.2588) -- (199.2148,91.2588) -- cycle(206.9148,83.8035) --
  (214.3701,83.8035) -- (214.3701,91.2588) -- (206.9148,91.2588) --
  cycle(214.6148,83.8035) -- (222.0701,83.8035) -- (222.0701,91.2588) --
  (214.6148,91.2588) -- cycle(222.3148,83.8035) -- (229.7701,83.8035) --
  (229.7701,91.2588) -- (222.3148,91.2588) -- cycle(230.0148,83.8035) --
  (237.4701,83.8035) -- (237.4701,91.2588) -- (230.0148,91.2588) --
  cycle(237.7148,83.8035) -- (245.1701,83.8035) -- (245.1701,91.2588) --
  (237.7148,91.2588) -- cycle(245.4148,83.8035) -- (252.8701,83.8035) --
  (252.8701,91.2588) -- (245.4148,91.2588) -- cycle(29.8148,91.5035) --
  (37.2701,91.5035) -- (37.2701,98.9588) -- (29.8148,98.9588) --
  cycle(37.5148,91.5035) -- (44.9701,91.5035) -- (44.9701,98.9588) --
  (37.5148,98.9588) -- cycle(45.2148,91.5035) -- (52.6701,91.5035) --
  (52.6701,98.9588) -- (45.2148,98.9588) -- cycle(52.9148,91.5035) --
  (60.3701,91.5035) -- (60.3701,98.9588) -- (52.9148,98.9588) --
  cycle(60.6148,91.5035) -- (68.0701,91.5035) -- (68.0701,98.9588) --
  (60.6148,98.9588) -- cycle(68.3148,91.5035) -- (75.7701,91.5035) --
  (75.7701,98.9588) -- (68.3148,98.9588) -- cycle(76.0148,91.5035) --
  (83.4701,91.5035) -- (83.4701,98.9588) -- (76.0148,98.9588) --
  cycle(83.7148,91.5035) -- (91.1701,91.5035) -- (91.1701,98.9588) --
  (83.7148,98.9588) -- cycle(91.4148,91.5035) -- (98.8701,91.5035) --
  (98.8701,98.9588) -- (91.4148,98.9588) -- cycle(99.1148,91.5035) --
  (106.5701,91.5035) -- (106.5701,98.9588) -- (99.1148,98.9588) --
  cycle(106.8148,91.5035) -- (114.2701,91.5035) -- (114.2701,98.9588) --
  (106.8148,98.9588) -- cycle(114.5148,91.5035) -- (121.9701,91.5035) --
  (121.9701,98.9588) -- (114.5148,98.9588) -- cycle(122.2148,91.5035) --
  (129.6701,91.5035) -- (129.6701,98.9588) -- (122.2148,98.9588) --
  cycle(129.9148,91.5035) -- (137.3701,91.5035) -- (137.3701,98.9588) --
  (129.9148,98.9588) -- cycle(137.6148,91.5035) -- (145.0701,91.5035) --
  (145.0701,98.9588) -- (137.6148,98.9588) -- cycle(145.3148,91.5035) --
  (152.7701,91.5035) -- (152.7701,98.9588) -- (145.3148,98.9588) --
  cycle(153.0148,91.5035) -- (160.4701,91.5035) -- (160.4701,98.9588) --
  (153.0148,98.9588) -- cycle(160.7148,91.5035) -- (168.1701,91.5035) --
  (168.1701,98.9588) -- (160.7148,98.9588) -- cycle(168.4148,91.5035) --
  (175.8701,91.5035) -- (175.8701,98.9588) -- (168.4148,98.9588) --
  cycle(176.1148,91.5035) -- (183.5701,91.5035) -- (183.5701,98.9588) --
  (176.1148,98.9588) -- cycle(183.8148,91.5035) -- (191.2701,91.5035) --
  (191.2701,98.9588) -- (183.8148,98.9588) -- cycle(191.5148,91.5035) --
  (198.9701,91.5035) -- (198.9701,98.9588) -- (191.5148,98.9588) --
  cycle(199.2148,91.5035) -- (206.6701,91.5035) -- (206.6701,98.9588) --
  (199.2148,98.9588) -- cycle(206.9148,91.5035) -- (214.3701,91.5035) --
  (214.3701,98.9588) -- (206.9148,98.9588) -- cycle(214.6148,91.5035) --
  (222.0701,91.5035) -- (222.0701,98.9588) -- (214.6148,98.9588) --
  cycle(222.3148,91.5035) -- (229.7701,91.5035) -- (229.7701,98.9588) --
  (222.3148,98.9588) -- cycle(230.0148,91.5035) -- (237.4701,91.5035) --
  (237.4701,98.9588) -- (230.0148,98.9588) -- cycle(237.7148,91.5035) --
  (245.1701,91.5035) -- (245.1701,98.9588) -- (237.7148,98.9588) --
  cycle(245.4148,91.5035) -- (252.8701,91.5035) -- (252.8701,98.9588) --
  (245.4148,98.9588) -- cycle(253.1148,91.5035) -- (260.5701,91.5035) --
  (260.5701,98.9588) -- (253.1148,98.9588) -- cycle(22.1148,99.2035) --
  (29.5701,99.2035) -- (29.5701,106.6588) -- (22.1148,106.6588) --
  cycle(29.8148,99.2035) -- (37.2701,99.2035) -- (37.2701,106.6588) --
  (29.8148,106.6588) -- cycle(37.5148,99.2035) -- (44.9701,99.2035) --
  (44.9701,106.6588) -- (37.5148,106.6588) -- cycle(45.2148,99.2035) --
  (52.6701,99.2035) -- (52.6701,106.6588) -- (45.2148,106.6588) --
  cycle(52.9148,99.2035) -- (60.3701,99.2035) -- (60.3701,106.6588) --
  (52.9148,106.6588) -- cycle(60.6148,99.2035) -- (68.0701,99.2035) --
  (68.0701,106.6588) -- (60.6148,106.6588) -- cycle(68.3148,99.2035) --
  (75.7701,99.2035) -- (75.7701,106.6588) -- (68.3148,106.6588) --
  cycle(76.0148,99.2035) -- (83.4701,99.2035) -- (83.4701,106.6588) --
  (76.0148,106.6588) -- cycle(83.7148,99.2035) -- (91.1701,99.2035) --
  (91.1701,106.6588) -- (83.7148,106.6588) -- cycle(91.4148,99.2035) --
  (98.8701,99.2035) -- (98.8701,106.6588) -- (91.4148,106.6588) --
  cycle(99.1148,99.2035) -- (106.5701,99.2035) -- (106.5701,106.6588) --
  (99.1148,106.6588) -- cycle(106.8148,99.2035) -- (114.2701,99.2035) --
  (114.2701,106.6588) -- (106.8148,106.6588) -- cycle(114.5148,99.2035) --
  (121.9701,99.2035) -- (121.9701,106.6588) -- (114.5148,106.6588) --
  cycle(122.2148,99.2035) -- (129.6701,99.2035) -- (129.6701,106.6588) --
  (122.2148,106.6588) -- cycle(129.9148,99.2035) -- (137.3701,99.2035) --
  (137.3701,106.6588) -- (129.9148,106.6588) -- cycle(137.6148,99.2035) --
  (145.0701,99.2035) -- (145.0701,106.6588) -- (137.6148,106.6588) --
  cycle(145.3148,99.2035) -- (152.7701,99.2035) -- (152.7701,106.6588) --
  (145.3148,106.6588) -- cycle(153.0148,99.2035) -- (160.4701,99.2035) --
  (160.4701,106.6588) -- (153.0148,106.6588) -- cycle(160.7148,99.2035) --
  (168.1701,99.2035) -- (168.1701,106.6588) -- (160.7148,106.6588) --
  cycle(168.4148,99.2035) -- (175.8701,99.2035) -- (175.8701,106.6588) --
  (168.4148,106.6588) -- cycle(176.1148,99.2035) -- (183.5701,99.2035) --
  (183.5701,106.6588) -- (176.1148,106.6588) -- cycle(183.8148,99.2035) --
  (191.2701,99.2035) -- (191.2701,106.6588) -- (183.8148,106.6588) --
  cycle(191.5148,99.2035) -- (198.9701,99.2035) -- (198.9701,106.6588) --
  (191.5148,106.6588) -- cycle(199.2148,99.2035) -- (206.6701,99.2035) --
  (206.6701,106.6588) -- (199.2148,106.6588) -- cycle(206.9148,99.2035) --
  (214.3701,99.2035) -- (214.3701,106.6588) -- (206.9148,106.6588) --
  cycle(214.6148,99.2035) -- (222.0701,99.2035) -- (222.0701,106.6588) --
  (214.6148,106.6588) -- cycle(222.3148,99.2035) -- (229.7701,99.2035) --
  (229.7701,106.6588) -- (222.3148,106.6588) -- cycle(230.0148,99.2035) --
  (237.4701,99.2035) -- (237.4701,106.6588) -- (230.0148,106.6588) --
  cycle(237.7148,99.2035) -- (245.1701,99.2035) -- (245.1701,106.6588) --
  (237.7148,106.6588) -- cycle(245.4148,99.2035) -- (252.8701,99.2035) --
  (252.8701,106.6588) -- (245.4148,106.6588) -- cycle(253.1148,99.2035) --
  (260.5701,99.2035) -- (260.5701,106.6588) -- (253.1148,106.6588) --
  cycle(260.8148,99.2035) -- (268.2701,99.2035) -- (268.2701,106.6588) --
  (260.8148,106.6588) -- cycle(14.4148,106.9035) -- (21.8701,106.9035) --
  (21.8701,114.3588) -- (14.4148,114.3588) -- cycle(22.1148,106.9035) --
  (29.5701,106.9035) -- (29.5701,114.3588) -- (22.1148,114.3588) --
  cycle(29.8148,106.9035) -- (37.2701,106.9035) -- (37.2701,114.3588) --
  (29.8148,114.3588) -- cycle(37.5148,106.9035) -- (44.9701,106.9035) --
  (44.9701,114.3588) -- (37.5148,114.3588) -- cycle(45.2148,106.9035) --
  (52.6701,106.9035) -- (52.6701,114.3588) -- (45.2148,114.3588) --
  cycle(52.9148,106.9035) -- (60.3701,106.9035) -- (60.3701,114.3588) --
  (52.9148,114.3588) -- cycle(60.6148,106.9035) -- (68.0701,106.9035) --
  (68.0701,114.3588) -- (60.6148,114.3588) -- cycle(68.3148,106.9035) --
  (75.7701,106.9035) -- (75.7701,114.3588) -- (68.3148,114.3588) --
  cycle(76.0148,106.9035) -- (83.4701,106.9035) -- (83.4701,114.3588) --
  (76.0148,114.3588) -- cycle(83.7148,106.9035) -- (91.1701,106.9035) --
  (91.1701,114.3588) -- (83.7148,114.3588) -- cycle(91.4148,106.9035) --
  (98.8701,106.9035) -- (98.8701,114.3588) -- (91.4148,114.3588) --
  cycle(99.1148,106.9035) -- (106.5701,106.9035) -- (106.5701,114.3588) --
  (99.1148,114.3588) -- cycle(106.8148,106.9035) -- (114.2701,106.9035) --
  (114.2701,114.3588) -- (106.8148,114.3588) -- cycle(114.5148,106.9035) --
  (121.9701,106.9035) -- (121.9701,114.3588) -- (114.5148,114.3588) --
  cycle(122.2148,106.9035) -- (129.6701,106.9035) -- (129.6701,114.3588) --
  (122.2148,114.3588) -- cycle(129.9148,106.9035) -- (137.3701,106.9035) --
  (137.3701,114.3588) -- (129.9148,114.3588) -- cycle(137.6148,106.9035) --
  (145.0701,106.9035) -- (145.0701,114.3588) -- (137.6148,114.3588) --
  cycle(145.3148,106.9035) -- (152.7701,106.9035) -- (152.7701,114.3588) --
  (145.3148,114.3588) -- cycle(153.0148,106.9035) -- (160.4701,106.9035) --
  (160.4701,114.3588) -- (153.0148,114.3588) -- cycle(160.7148,106.9035) --
  (168.1701,106.9035) -- (168.1701,114.3588) -- (160.7148,114.3588) --
  cycle(168.4148,106.9035) -- (175.8701,106.9035) -- (175.8701,114.3588) --
  (168.4148,114.3588) -- cycle(176.1148,106.9035) -- (183.5701,106.9035) --
  (183.5701,114.3588) -- (176.1148,114.3588) -- cycle(183.8148,106.9035) --
  (191.2701,106.9035) -- (191.2701,114.3588) -- (183.8148,114.3588) --
  cycle(191.5148,106.9035) -- (198.9701,106.9035) -- (198.9701,114.3588) --
  (191.5148,114.3588) -- cycle(199.2148,106.9035) -- (206.6701,106.9035) --
  (206.6701,114.3588) -- (199.2148,114.3588) -- cycle(206.9148,106.9035) --
  (214.3701,106.9035) -- (214.3701,114.3588) -- (206.9148,114.3588) --
  cycle(214.6148,106.9035) -- (222.0701,106.9035) -- (222.0701,114.3588) --
  (214.6148,114.3588) -- cycle(222.3148,106.9035) -- (229.7701,106.9035) --
  (229.7701,114.3588) -- (222.3148,114.3588) -- cycle(230.0148,106.9035) --
  (237.4701,106.9035) -- (237.4701,114.3588) -- (230.0148,114.3588) --
  cycle(237.7148,106.9035) -- (245.1701,106.9035) -- (245.1701,114.3588) --
  (237.7148,114.3588) -- cycle(245.4148,106.9035) -- (252.8701,106.9035) --
  (252.8701,114.3588) -- (245.4148,114.3588) -- cycle(253.1148,106.9035) --
  (260.5701,106.9035) -- (260.5701,114.3588) -- (253.1148,114.3588) --
  cycle(260.8148,106.9035) -- (268.2701,106.9035) -- (268.2701,114.3588) --
  (260.8148,114.3588) -- cycle(14.4148,114.6035) -- (21.8701,114.6035) --
  (21.8701,122.0588) -- (14.4148,122.0588) -- cycle(22.1148,114.6035) --
  (29.5701,114.6035) -- (29.5701,122.0588) -- (22.1148,122.0588) --
  cycle(29.8148,114.6035) -- (37.2701,114.6035) -- (37.2701,122.0588) --
  (29.8148,122.0588) -- cycle(37.5148,114.6035) -- (44.9701,114.6035) --
  (44.9701,122.0588) -- (37.5148,122.0588) -- cycle(45.2148,114.6035) --
  (52.6701,114.6035) -- (52.6701,122.0588) -- (45.2148,122.0588) --
  cycle(52.9148,114.6035) -- (60.3701,114.6035) -- (60.3701,122.0588) --
  (52.9148,122.0588) -- cycle(60.6148,114.6035) -- (68.0701,114.6035) --
  (68.0701,122.0588) -- (60.6148,122.0588) -- cycle(68.3148,114.6035) --
  (75.7701,114.6035) -- (75.7701,122.0588) -- (68.3148,122.0588) --
  cycle(76.0148,114.6035) -- (83.4701,114.6035) -- (83.4701,122.0588) --
  (76.0148,122.0588) -- cycle(83.7148,114.6035) -- (91.1701,114.6035) --
  (91.1701,122.0588) -- (83.7148,122.0588) -- cycle(91.4148,114.6035) --
  (98.8701,114.6035) -- (98.8701,122.0588) -- (91.4148,122.0588) --
  cycle(99.1148,114.6035) -- (106.5701,114.6035) -- (106.5701,122.0588) --
  (99.1148,122.0588) -- cycle(106.8148,114.6035) -- (114.2701,114.6035) --
  (114.2701,122.0588) -- (106.8148,122.0588) -- cycle(114.5148,114.6035) --
  (121.9701,114.6035) -- (121.9701,122.0588) -- (114.5148,122.0588) --
  cycle(122.2148,114.6035) -- (129.6701,114.6035) -- (129.6701,122.0588) --
  (122.2148,122.0588) -- cycle(129.9148,114.6035) -- (137.3701,114.6035) --
  (137.3701,122.0588) -- (129.9148,122.0588) -- cycle(137.6148,114.6035) --
  (145.0701,114.6035) -- (145.0701,122.0588) -- (137.6148,122.0588) --
  cycle(145.3148,114.6035) -- (152.7701,114.6035) -- (152.7701,122.0588) --
  (145.3148,122.0588) -- cycle(153.0148,114.6035) -- (160.4701,114.6035) --
  (160.4701,122.0588) -- (153.0148,122.0588) -- cycle(160.7148,114.6035) --
  (168.1701,114.6035) -- (168.1701,122.0588) -- (160.7148,122.0588) --
  cycle(168.4148,114.6035) -- (175.8701,114.6035) -- (175.8701,122.0588) --
  (168.4148,122.0588) -- cycle(176.1148,114.6035) -- (183.5701,114.6035) --
  (183.5701,122.0588) -- (176.1148,122.0588) -- cycle(183.8148,114.6035) --
  (191.2701,114.6035) -- (191.2701,122.0588) -- (183.8148,122.0588) --
  cycle(191.5148,114.6035) -- (198.9701,114.6035) -- (198.9701,122.0588) --
  (191.5148,122.0588) -- cycle(199.2148,114.6035) -- (206.6701,114.6035) --
  (206.6701,122.0588) -- (199.2148,122.0588) -- cycle(206.9148,114.6035) --
  (214.3701,114.6035) -- (214.3701,122.0588) -- (206.9148,122.0588) --
  cycle(214.6148,114.6035) -- (222.0701,114.6035) -- (222.0701,122.0588) --
  (214.6148,122.0588) -- cycle(222.3148,114.6035) -- (229.7701,114.6035) --
  (229.7701,122.0588) -- (222.3148,122.0588) -- cycle(230.0148,114.6035) --
  (237.4701,114.6035) -- (237.4701,122.0588) -- (230.0148,122.0588) --
  cycle(237.7148,114.6035) -- (245.1701,114.6035) -- (245.1701,122.0588) --
  (237.7148,122.0588) -- cycle(245.4148,114.6035) -- (252.8701,114.6035) --
  (252.8701,122.0588) -- (245.4148,122.0588) -- cycle(253.1148,114.6035) --
  (260.5701,114.6035) -- (260.5701,122.0588) -- (253.1148,122.0588) --
  cycle(260.8148,114.6035) -- (268.2701,114.6035) -- (268.2701,122.0588) --
  (260.8148,122.0588) -- cycle(268.5148,114.6035) -- (275.9701,114.6035) --
  (275.9701,122.0588) -- (268.5148,122.0588) -- cycle(6.7148,122.3035) --
  (14.1701,122.3035) -- (14.1701,129.7588) -- (6.7148,129.7588) --
  cycle(14.4148,122.3035) -- (21.8701,122.3035) -- (21.8701,129.7588) --
  (14.4148,129.7588) -- cycle(22.1148,122.3035) -- (29.5701,122.3035) --
  (29.5701,129.7588) -- (22.1148,129.7588) -- cycle(29.8148,122.3035) --
  (37.2701,122.3035) -- (37.2701,129.7588) -- (29.8148,129.7588) --
  cycle(37.5148,122.3035) -- (44.9701,122.3035) -- (44.9701,129.7588) --
  (37.5148,129.7588) -- cycle(45.2148,122.3035) -- (52.6701,122.3035) --
  (52.6701,129.7588) -- (45.2148,129.7588) -- cycle(52.9148,122.3035) --
  (60.3701,122.3035) -- (60.3701,129.7588) -- (52.9148,129.7588) --
  cycle(60.6148,122.3035) -- (68.0701,122.3035) -- (68.0701,129.7588) --
  (60.6148,129.7588) -- cycle(68.3148,122.3035) -- (75.7701,122.3035) --
  (75.7701,129.7588) -- (68.3148,129.7588) -- cycle(76.0148,122.3035) --
  (83.4701,122.3035) -- (83.4701,129.7588) -- (76.0148,129.7588) --
  cycle(83.7148,122.3035) -- (91.1701,122.3035) -- (91.1701,129.7588) --
  (83.7148,129.7588) -- cycle(91.4148,122.3035) -- (98.8701,122.3035) --
  (98.8701,129.7588) -- (91.4148,129.7588) -- cycle(99.1148,122.3035) --
  (106.5701,122.3035) -- (106.5701,129.7588) -- (99.1148,129.7588) --
  cycle(106.8148,122.3035) -- (114.2701,122.3035) -- (114.2701,129.7588) --
  (106.8148,129.7588) -- cycle(114.5148,122.3035) -- (121.9701,122.3035) --
  (121.9701,129.7588) -- (114.5148,129.7588) -- cycle(122.2148,122.3035) --
  (129.6701,122.3035) -- (129.6701,129.7588) -- (122.2148,129.7588) --
  cycle(129.9148,122.3035) -- (137.3701,122.3035) -- (137.3701,129.7588) --
  (129.9148,129.7588) -- cycle(137.6148,122.3035) -- (145.0701,122.3035) --
  (145.0701,129.7588) -- (137.6148,129.7588) -- cycle(145.3148,122.3035) --
  (152.7701,122.3035) -- (152.7701,129.7588) -- (145.3148,129.7588) --
  cycle(153.0148,122.3035) -- (160.4701,122.3035) -- (160.4701,129.7588) --
  (153.0148,129.7588) -- cycle(160.7148,122.3035) -- (168.1701,122.3035) --
  (168.1701,129.7588) -- (160.7148,129.7588) -- cycle(168.4148,122.3035) --
  (175.8701,122.3035) -- (175.8701,129.7588) -- (168.4148,129.7588) --
  cycle(176.1148,122.3035) -- (183.5701,122.3035) -- (183.5701,129.7588) --
  (176.1148,129.7588) -- cycle(183.8148,122.3035) -- (191.2701,122.3035) --
  (191.2701,129.7588) -- (183.8148,129.7588) -- cycle(191.5148,122.3035) --
  (198.9701,122.3035) -- (198.9701,129.7588) -- (191.5148,129.7588) --
  cycle(199.2148,122.3035) -- (206.6701,122.3035) -- (206.6701,129.7588) --
  (199.2148,129.7588) -- cycle(206.9148,122.3035) -- (214.3701,122.3035) --
  (214.3701,129.7588) -- (206.9148,129.7588) -- cycle(214.6148,122.3035) --
  (222.0701,122.3035) -- (222.0701,129.7588) -- (214.6148,129.7588) --
  cycle(222.3148,122.3035) -- (229.7701,122.3035) -- (229.7701,129.7588) --
  (222.3148,129.7588) -- cycle(230.0148,122.3035) -- (237.4701,122.3035) --
  (237.4701,129.7588) -- (230.0148,129.7588) -- cycle(237.7148,122.3035) --
  (245.1701,122.3035) -- (245.1701,129.7588) -- (237.7148,129.7588) --
  cycle(245.4148,122.3035) -- (252.8701,122.3035) -- (252.8701,129.7588) --
  (245.4148,129.7588) -- cycle(253.1148,122.3035) -- (260.5701,122.3035) --
  (260.5701,129.7588) -- (253.1148,129.7588) -- cycle(260.8148,122.3035) --
  (268.2701,122.3035) -- (268.2701,129.7588) -- (260.8148,129.7588) --
  cycle(268.5148,122.3035) -- (275.9701,122.3035) -- (275.9701,129.7588) --
  (268.5148,129.7588) -- cycle(6.7148,130.0035) -- (14.1701,130.0035) --
  (14.1701,137.4588) -- (6.7148,137.4588) -- cycle(14.4148,130.0035) --
  (21.8701,130.0035) -- (21.8701,137.4588) -- (14.4148,137.4588) --
  cycle(22.1148,130.0035) -- (29.5701,130.0035) -- (29.5701,137.4588) --
  (22.1148,137.4588) -- cycle(29.8148,130.0035) -- (37.2701,130.0035) --
  (37.2701,137.4588) -- (29.8148,137.4588) -- cycle(37.5148,130.0035) --
  (44.9701,130.0035) -- (44.9701,137.4588) -- (37.5148,137.4588) --
  cycle(45.2148,130.0035) -- (52.6701,130.0035) -- (52.6701,137.4588) --
  (45.2148,137.4588) -- cycle(52.9148,130.0035) -- (60.3701,130.0035) --
  (60.3701,137.4588) -- (52.9148,137.4588) -- cycle(60.6148,130.0035) --
  (68.0701,130.0035) -- (68.0701,137.4588) -- (60.6148,137.4588) --
  cycle(68.3148,130.0035) -- (75.7701,130.0035) -- (75.7701,137.4588) --
  (68.3148,137.4588) -- cycle(76.0148,130.0035) -- (83.4701,130.0035) --
  (83.4701,137.4588) -- (76.0148,137.4588) -- cycle(83.7148,130.0035) --
  (91.1701,130.0035) -- (91.1701,137.4588) -- (83.7148,137.4588) --
  cycle(91.4148,130.0035) -- (98.8701,130.0035) -- (98.8701,137.4588) --
  (91.4148,137.4588) -- cycle(99.1148,130.0035) -- (106.5701,130.0035) --
  (106.5701,137.4588) -- (99.1148,137.4588) -- cycle(106.8148,130.0035) --
  (114.2701,130.0035) -- (114.2701,137.4588) -- (106.8148,137.4588) --
  cycle(114.5148,130.0035) -- (121.9701,130.0035) -- (121.9701,137.4588) --
  (114.5148,137.4588) -- cycle(122.2148,130.0035) -- (129.6701,130.0035) --
  (129.6701,137.4588) -- (122.2148,137.4588) -- cycle(129.9148,130.0035) --
  (137.3701,130.0035) -- (137.3701,137.4588) -- (129.9148,137.4588) --
  cycle(137.6148,130.0035) -- (145.0701,130.0035) -- (145.0701,137.4588) --
  (137.6148,137.4588) -- cycle(145.3148,130.0035) -- (152.7701,130.0035) --
  (152.7701,137.4588) -- (145.3148,137.4588) -- cycle(153.0148,130.0035) --
  (160.4701,130.0035) -- (160.4701,137.4588) -- (153.0148,137.4588) --
  cycle(160.7148,130.0035) -- (168.1701,130.0035) -- (168.1701,137.4588) --
  (160.7148,137.4588) -- cycle(168.4148,130.0035) -- (175.8701,130.0035) --
  (175.8701,137.4588) -- (168.4148,137.4588) -- cycle(176.1148,130.0035) --
  (183.5701,130.0035) -- (183.5701,137.4588) -- (176.1148,137.4588) --
  cycle(183.8148,130.0035) -- (191.2701,130.0035) -- (191.2701,137.4588) --
  (183.8148,137.4588) -- cycle(191.5148,130.0035) -- (198.9701,130.0035) --
  (198.9701,137.4588) -- (191.5148,137.4588) -- cycle(199.2148,130.0035) --
  (206.6701,130.0035) -- (206.6701,137.4588) -- (199.2148,137.4588) --
  cycle(206.9148,130.0035) -- (214.3701,130.0035) -- (214.3701,137.4588) --
  (206.9148,137.4588) -- cycle(214.6148,130.0035) -- (222.0701,130.0035) --
  (222.0701,137.4588) -- (214.6148,137.4588) -- cycle(222.3148,130.0035) --
  (229.7701,130.0035) -- (229.7701,137.4588) -- (222.3148,137.4588) --
  cycle(230.0148,130.0035) -- (237.4701,130.0035) -- (237.4701,137.4588) --
  (230.0148,137.4588) -- cycle(237.7148,130.0035) -- (245.1701,130.0035) --
  (245.1701,137.4588) -- (237.7148,137.4588) -- cycle(245.4148,130.0035) --
  (252.8701,130.0035) -- (252.8701,137.4588) -- (245.4148,137.4588) --
  cycle(253.1148,130.0035) -- (260.5701,130.0035) -- (260.5701,137.4588) --
  (253.1148,137.4588) -- cycle(260.8148,130.0035) -- (268.2701,130.0035) --
  (268.2701,137.4588) -- (260.8148,137.4588) -- cycle(268.5148,130.0035) --
  (275.9701,130.0035) -- (275.9701,137.4588) -- (268.5148,137.4588) --
  cycle(14.4148,137.7035) -- (21.8701,137.7035) -- (21.8701,145.1588) --
  (14.4148,145.1588) -- cycle(22.1148,137.7035) -- (29.5701,137.7035) --
  (29.5701,145.1588) -- (22.1148,145.1588) -- cycle(29.8148,137.7035) --
  (37.2701,137.7035) -- (37.2701,145.1588) -- (29.8148,145.1588) --
  cycle(37.5148,137.7035) -- (44.9701,137.7035) -- (44.9701,145.1588) --
  (37.5148,145.1588) -- cycle(45.2148,137.7035) -- (52.6701,137.7035) --
  (52.6701,145.1588) -- (45.2148,145.1588) -- cycle(52.9148,137.7035) --
  (60.3701,137.7035) -- (60.3701,145.1588) -- (52.9148,145.1588) --
  cycle(60.6148,137.7035) -- (68.0701,137.7035) -- (68.0701,145.1588) --
  (60.6148,145.1588) -- cycle(68.3148,137.7035) -- (75.7701,137.7035) --
  (75.7701,145.1588) -- (68.3148,145.1588) -- cycle(76.0148,137.7035) --
  (83.4701,137.7035) -- (83.4701,145.1588) -- (76.0148,145.1588) --
  cycle(83.7148,137.7035) -- (91.1701,137.7035) -- (91.1701,145.1588) --
  (83.7148,145.1588) -- cycle(91.4148,137.7035) -- (98.8701,137.7035) --
  (98.8701,145.1588) -- (91.4148,145.1588) -- cycle(99.1148,137.7035) --
  (106.5701,137.7035) -- (106.5701,145.1588) -- (99.1148,145.1588) --
  cycle(106.8148,137.7035) -- (114.2701,137.7035) -- (114.2701,145.1588) --
  (106.8148,145.1588) -- cycle(114.5148,137.7035) -- (121.9701,137.7035) --
  (121.9701,145.1588) -- (114.5148,145.1588) -- cycle(122.2148,137.7035) --
  (129.6701,137.7035) -- (129.6701,145.1588) -- (122.2148,145.1588) --
  cycle(129.9148,137.7035) -- (137.3701,137.7035) -- (137.3701,145.1588) --
  (129.9148,145.1588) -- cycle(137.6148,137.7035) -- (145.0701,137.7035) --
  (145.0701,145.1588) -- (137.6148,145.1588) -- cycle(145.3148,137.7035) --
  (152.7701,137.7035) -- (152.7701,145.1588) -- (145.3148,145.1588) --
  cycle(153.0148,137.7035) -- (160.4701,137.7035) -- (160.4701,145.1588) --
  (153.0148,145.1588) -- cycle(160.7148,137.7035) -- (168.1701,137.7035) --
  (168.1701,145.1588) -- (160.7148,145.1588) -- cycle(168.4148,137.7035) --
  (175.8701,137.7035) -- (175.8701,145.1588) -- (168.4148,145.1588) --
  cycle(176.1148,137.7035) -- (183.5701,137.7035) -- (183.5701,145.1588) --
  (176.1148,145.1588) -- cycle(183.8148,137.7035) -- (191.2701,137.7035) --
  (191.2701,145.1588) -- (183.8148,145.1588) -- cycle(191.5148,137.7035) --
  (198.9701,137.7035) -- (198.9701,145.1588) -- (191.5148,145.1588) --
  cycle(199.2148,137.7035) -- (206.6701,137.7035) -- (206.6701,145.1588) --
  (199.2148,145.1588) -- cycle(206.9148,137.7035) -- (214.3701,137.7035) --
  (214.3701,145.1588) -- (206.9148,145.1588) -- cycle(214.6148,137.7035) --
  (222.0701,137.7035) -- (222.0701,145.1588) -- (214.6148,145.1588) --
  cycle(222.3148,137.7035) -- (229.7701,137.7035) -- (229.7701,145.1588) --
  (222.3148,145.1588) -- cycle(230.0148,137.7035) -- (237.4701,137.7035) --
  (237.4701,145.1588) -- (230.0148,145.1588) -- cycle(237.7148,137.7035) --
  (245.1701,137.7035) -- (245.1701,145.1588) -- (237.7148,145.1588) --
  cycle(245.4148,137.7035) -- (252.8701,137.7035) -- (252.8701,145.1588) --
  (245.4148,145.1588) -- cycle(253.1148,137.7035) -- (260.5701,137.7035) --
  (260.5701,145.1588) -- (253.1148,145.1588) -- cycle(260.8148,137.7035) --
  (268.2701,137.7035) -- (268.2701,145.1588) -- (260.8148,145.1588) --
  cycle(268.5148,137.7035) -- (275.9701,137.7035) -- (275.9701,145.1588) --
  (268.5148,145.1588) -- cycle(22.1148,145.4035) -- (29.5701,145.4035) --
  (29.5701,152.8588) -- (22.1148,152.8588) -- cycle(29.8148,145.4035) --
  (37.2701,145.4035) -- (37.2701,152.8588) -- (29.8148,152.8588) --
  cycle(37.5148,145.4035) -- (44.9701,145.4035) -- (44.9701,152.8588) --
  (37.5148,152.8588) -- cycle(45.2148,145.4035) -- (52.6701,145.4035) --
  (52.6701,152.8588) -- (45.2148,152.8588) -- cycle(52.9148,145.4035) --
  (60.3701,145.4035) -- (60.3701,152.8588) -- (52.9148,152.8588) --
  cycle(60.6148,145.4035) -- (68.0701,145.4035) -- (68.0701,152.8588) --
  (60.6148,152.8588) -- cycle(68.3148,145.4035) -- (75.7701,145.4035) --
  (75.7701,152.8588) -- (68.3148,152.8588) -- cycle(76.0148,145.4035) --
  (83.4701,145.4035) -- (83.4701,152.8588) -- (76.0148,152.8588) --
  cycle(83.7148,145.4035) -- (91.1701,145.4035) -- (91.1701,152.8588) --
  (83.7148,152.8588) -- cycle(91.4148,145.4035) -- (98.8701,145.4035) --
  (98.8701,152.8588) -- (91.4148,152.8588) -- cycle(99.1148,145.4035) --
  (106.5701,145.4035) -- (106.5701,152.8588) -- (99.1148,152.8588) --
  cycle(106.8148,145.4035) -- (114.2701,145.4035) -- (114.2701,152.8588) --
  (106.8148,152.8588) -- cycle(114.5148,145.4035) -- (121.9701,145.4035) --
  (121.9701,152.8588) -- (114.5148,152.8588) -- cycle(122.2148,145.4035) --
  (129.6701,145.4035) -- (129.6701,152.8588) -- (122.2148,152.8588) --
  cycle(129.9148,145.4035) -- (137.3701,145.4035) -- (137.3701,152.8588) --
  (129.9148,152.8588) -- cycle(137.6148,145.4035) -- (145.0701,145.4035) --
  (145.0701,152.8588) -- (137.6148,152.8588) -- cycle(145.3148,145.4035) --
  (152.7701,145.4035) -- (152.7701,152.8588) -- (145.3148,152.8588) --
  cycle(153.0148,145.4035) -- (160.4701,145.4035) -- (160.4701,152.8588) --
  (153.0148,152.8588) -- cycle(160.7148,145.4035) -- (168.1701,145.4035) --
  (168.1701,152.8588) -- (160.7148,152.8588) -- cycle(168.4148,145.4035) --
  (175.8701,145.4035) -- (175.8701,152.8588) -- (168.4148,152.8588) --
  cycle(176.1148,145.4035) -- (183.5701,145.4035) -- (183.5701,152.8588) --
  (176.1148,152.8588) -- cycle(183.8148,145.4035) -- (191.2701,145.4035) --
  (191.2701,152.8588) -- (183.8148,152.8588) -- cycle(191.5148,145.4035) --
  (198.9701,145.4035) -- (198.9701,152.8588) -- (191.5148,152.8588) --
  cycle(199.2148,145.4035) -- (206.6701,145.4035) -- (206.6701,152.8588) --
  (199.2148,152.8588) -- cycle(206.9148,145.4035) -- (214.3701,145.4035) --
  (214.3701,152.8588) -- (206.9148,152.8588) -- cycle(214.6148,145.4035) --
  (222.0701,145.4035) -- (222.0701,152.8588) -- (214.6148,152.8588) --
  cycle(222.3148,145.4035) -- (229.7701,145.4035) -- (229.7701,152.8588) --
  (222.3148,152.8588) -- cycle(230.0148,145.4035) -- (237.4701,145.4035) --
  (237.4701,152.8588) -- (230.0148,152.8588) -- cycle(237.7148,145.4035) --
  (245.1701,145.4035) -- (245.1701,152.8588) -- (237.7148,152.8588) --
  cycle(245.4148,145.4035) -- (252.8701,145.4035) -- (252.8701,152.8588) --
  (245.4148,152.8588) -- cycle(253.1148,145.4035) -- (260.5701,145.4035) --
  (260.5701,152.8588) -- (253.1148,152.8588) -- cycle(260.8148,145.4035) --
  (268.2701,145.4035) -- (268.2701,152.8588) -- (260.8148,152.8588) --
  cycle(22.1148,153.1035) -- (29.5701,153.1035) -- (29.5701,160.5588) --
  (22.1148,160.5588) -- cycle(29.8148,153.1035) -- (37.2701,153.1035) --
  (37.2701,160.5588) -- (29.8148,160.5588) -- cycle(37.5148,153.1035) --
  (44.9701,153.1035) -- (44.9701,160.5588) -- (37.5148,160.5588) --
  cycle(45.2148,153.1035) -- (52.6701,153.1035) -- (52.6701,160.5588) --
  (45.2148,160.5588) -- cycle(52.9148,153.1035) -- (60.3701,153.1035) --
  (60.3701,160.5588) -- (52.9148,160.5588) -- cycle(60.6148,153.1035) --
  (68.0701,153.1035) -- (68.0701,160.5588) -- (60.6148,160.5588) --
  cycle(68.3148,153.1035) -- (75.7701,153.1035) -- (75.7701,160.5588) --
  (68.3148,160.5588) -- cycle(76.0148,153.1035) -- (83.4701,153.1035) --
  (83.4701,160.5588) -- (76.0148,160.5588) -- cycle(83.7148,153.1035) --
  (91.1701,153.1035) -- (91.1701,160.5588) -- (83.7148,160.5588) --
  cycle(91.4148,153.1035) -- (98.8701,153.1035) -- (98.8701,160.5588) --
  (91.4148,160.5588) -- cycle(99.1148,153.1035) -- (106.5701,153.1035) --
  (106.5701,160.5588) -- (99.1148,160.5588) -- cycle(106.8148,153.1035) --
  (114.2701,153.1035) -- (114.2701,160.5588) -- (106.8148,160.5588) --
  cycle(114.5148,153.1035) -- (121.9701,153.1035) -- (121.9701,160.5588) --
  (114.5148,160.5588) -- cycle(122.2148,153.1035) -- (129.6701,153.1035) --
  (129.6701,160.5588) -- (122.2148,160.5588) -- cycle(129.9148,153.1035) --
  (137.3701,153.1035) -- (137.3701,160.5588) -- (129.9148,160.5588) --
  cycle(137.6148,153.1035) -- (145.0701,153.1035) -- (145.0701,160.5588) --
  (137.6148,160.5588) -- cycle(145.3148,153.1035) -- (152.7701,153.1035) --
  (152.7701,160.5588) -- (145.3148,160.5588) -- cycle(153.0148,153.1035) --
  (160.4701,153.1035) -- (160.4701,160.5588) -- (153.0148,160.5588) --
  cycle(160.7148,153.1035) -- (168.1701,153.1035) -- (168.1701,160.5588) --
  (160.7148,160.5588) -- cycle(168.4148,153.1035) -- (175.8701,153.1035) --
  (175.8701,160.5588) -- (168.4148,160.5588) -- cycle(176.1148,153.1035) --
  (183.5701,153.1035) -- (183.5701,160.5588) -- (176.1148,160.5588) --
  cycle(183.8148,153.1035) -- (191.2701,153.1035) -- (191.2701,160.5588) --
  (183.8148,160.5588) -- cycle(191.5148,153.1035) -- (198.9701,153.1035) --
  (198.9701,160.5588) -- (191.5148,160.5588) -- cycle(199.2148,153.1035) --
  (206.6701,153.1035) -- (206.6701,160.5588) -- (199.2148,160.5588) --
  cycle(206.9148,153.1035) -- (214.3701,153.1035) -- (214.3701,160.5588) --
  (206.9148,160.5588) -- cycle(214.6148,153.1035) -- (222.0701,153.1035) --
  (222.0701,160.5588) -- (214.6148,160.5588) -- cycle(222.3148,153.1035) --
  (229.7701,153.1035) -- (229.7701,160.5588) -- (222.3148,160.5588) --
  cycle(230.0148,153.1035) -- (237.4701,153.1035) -- (237.4701,160.5588) --
  (230.0148,160.5588) -- cycle(237.7148,153.1035) -- (245.1701,153.1035) --
  (245.1701,160.5588) -- (237.7148,160.5588) -- cycle(245.4148,153.1035) --
  (252.8701,153.1035) -- (252.8701,160.5588) -- (245.4148,160.5588) --
  cycle(253.1148,153.1035) -- (260.5701,153.1035) -- (260.5701,160.5588) --
  (253.1148,160.5588) -- cycle(29.8148,160.8035) -- (37.2701,160.8035) --
  (37.2701,168.2588) -- (29.8148,168.2588) -- cycle(37.5148,160.8035) --
  (44.9701,160.8035) -- (44.9701,168.2588) -- (37.5148,168.2588) --
  cycle(45.2148,160.8035) -- (52.6701,160.8035) -- (52.6701,168.2588) --
  (45.2148,168.2588) -- cycle(52.9148,160.8035) -- (60.3701,160.8035) --
  (60.3701,168.2588) -- (52.9148,168.2588) -- cycle(60.6148,160.8035) --
  (68.0701,160.8035) -- (68.0701,168.2588) -- (60.6148,168.2588) --
  cycle(68.3148,160.8035) -- (75.7701,160.8035) -- (75.7701,168.2588) --
  (68.3148,168.2588) -- cycle(76.0148,160.8035) -- (83.4701,160.8035) --
  (83.4701,168.2588) -- (76.0148,168.2588) -- cycle(83.7148,160.8035) --
  (91.1701,160.8035) -- (91.1701,168.2588) -- (83.7148,168.2588) --
  cycle(91.4148,160.8035) -- (98.8701,160.8035) -- (98.8701,168.2588) --
  (91.4148,168.2588) -- cycle(99.1148,160.8035) -- (106.5701,160.8035) --
  (106.5701,168.2588) -- (99.1148,168.2588) -- cycle(106.8148,160.8035) --
  (114.2701,160.8035) -- (114.2701,168.2588) -- (106.8148,168.2588) --
  cycle(114.5148,160.8035) -- (121.9701,160.8035) -- (121.9701,168.2588) --
  (114.5148,168.2588) -- cycle(122.2148,160.8035) -- (129.6701,160.8035) --
  (129.6701,168.2588) -- (122.2148,168.2588) -- cycle(129.9148,160.8035) --
  (137.3701,160.8035) -- (137.3701,168.2588) -- (129.9148,168.2588) --
  cycle(137.6148,160.8035) -- (145.0701,160.8035) -- (145.0701,168.2588) --
  (137.6148,168.2588) -- cycle(145.3148,160.8035) -- (152.7701,160.8035) --
  (152.7701,168.2588) -- (145.3148,168.2588) -- cycle(153.0148,160.8035) --
  (160.4701,160.8035) -- (160.4701,168.2588) -- (153.0148,168.2588) --
  cycle(160.7148,160.8035) -- (168.1701,160.8035) -- (168.1701,168.2588) --
  (160.7148,168.2588) -- cycle(168.4148,160.8035) -- (175.8701,160.8035) --
  (175.8701,168.2588) -- (168.4148,168.2588) -- cycle(176.1148,160.8035) --
  (183.5701,160.8035) -- (183.5701,168.2588) -- (176.1148,168.2588) --
  cycle(183.8148,160.8035) -- (191.2701,160.8035) -- (191.2701,168.2588) --
  (183.8148,168.2588) -- cycle(191.5148,160.8035) -- (198.9701,160.8035) --
  (198.9701,168.2588) -- (191.5148,168.2588) -- cycle(199.2148,160.8035) --
  (206.6701,160.8035) -- (206.6701,168.2588) -- (199.2148,168.2588) --
  cycle(206.9148,160.8035) -- (214.3701,160.8035) -- (214.3701,168.2588) --
  (206.9148,168.2588) -- cycle(214.6148,160.8035) -- (222.0701,160.8035) --
  (222.0701,168.2588) -- (214.6148,168.2588) -- cycle(222.3148,160.8035) --
  (229.7701,160.8035) -- (229.7701,168.2588) -- (222.3148,168.2588) --
  cycle(230.0148,160.8035) -- (237.4701,160.8035) -- (237.4701,168.2588) --
  (230.0148,168.2588) -- cycle(237.7148,160.8035) -- (245.1701,160.8035) --
  (245.1701,168.2588) -- (237.7148,168.2588) -- cycle(245.4148,160.8035) --
  (252.8701,160.8035) -- (252.8701,168.2588) -- (245.4148,168.2588) --
  cycle(253.1148,160.8035) -- (260.5701,160.8035) -- (260.5701,168.2588) --
  (253.1148,168.2588) -- cycle(29.8148,168.5035) -- (37.2701,168.5035) --
  (37.2701,175.9588) -- (29.8148,175.9588) -- cycle(37.5148,168.5035) --
  (44.9701,168.5035) -- (44.9701,175.9588) -- (37.5148,175.9588) --
  cycle(45.2148,168.5035) -- (52.6701,168.5035) -- (52.6701,175.9588) --
  (45.2148,175.9588) -- cycle(52.9148,168.5035) -- (60.3701,168.5035) --
  (60.3701,175.9588) -- (52.9148,175.9588) -- cycle(60.6148,168.5035) --
  (68.0701,168.5035) -- (68.0701,175.9588) -- (60.6148,175.9588) --
  cycle(68.3148,168.5035) -- (75.7701,168.5035) -- (75.7701,175.9588) --
  (68.3148,175.9588) -- cycle(76.0148,168.5035) -- (83.4701,168.5035) --
  (83.4701,175.9588) -- (76.0148,175.9588) -- cycle(83.7148,168.5035) --
  (91.1701,168.5035) -- (91.1701,175.9588) -- (83.7148,175.9588) --
  cycle(91.4148,168.5035) -- (98.8701,168.5035) -- (98.8701,175.9588) --
  (91.4148,175.9588) -- cycle(99.1148,168.5035) -- (106.5701,168.5035) --
  (106.5701,175.9588) -- (99.1148,175.9588) -- cycle(106.8148,168.5035) --
  (114.2701,168.5035) -- (114.2701,175.9588) -- (106.8148,175.9588) --
  cycle(114.5148,168.5035) -- (121.9701,168.5035) -- (121.9701,175.9588) --
  (114.5148,175.9588) -- cycle(122.2148,168.5035) -- (129.6701,168.5035) --
  (129.6701,175.9588) -- (122.2148,175.9588) -- cycle(129.9148,168.5035) --
  (137.3701,168.5035) -- (137.3701,175.9588) -- (129.9148,175.9588) --
  cycle(137.6148,168.5035) -- (145.0701,168.5035) -- (145.0701,175.9588) --
  (137.6148,175.9588) -- cycle(145.3148,168.5035) -- (152.7701,168.5035) --
  (152.7701,175.9588) -- (145.3148,175.9588) -- cycle(153.0148,168.5035) --
  (160.4701,168.5035) -- (160.4701,175.9588) -- (153.0148,175.9588) --
  cycle(160.7148,168.5035) -- (168.1701,168.5035) -- (168.1701,175.9588) --
  (160.7148,175.9588) -- cycle(168.4148,168.5035) -- (175.8701,168.5035) --
  (175.8701,175.9588) -- (168.4148,175.9588) -- cycle(176.1148,168.5035) --
  (183.5701,168.5035) -- (183.5701,175.9588) -- (176.1148,175.9588) --
  cycle(183.8148,168.5035) -- (191.2701,168.5035) -- (191.2701,175.9588) --
  (183.8148,175.9588) -- cycle(191.5148,168.5035) -- (198.9701,168.5035) --
  (198.9701,175.9588) -- (191.5148,175.9588) -- cycle(199.2148,168.5035) --
  (206.6701,168.5035) -- (206.6701,175.9588) -- (199.2148,175.9588) --
  cycle(206.9148,168.5035) -- (214.3701,168.5035) -- (214.3701,175.9588) --
  (206.9148,175.9588) -- cycle(214.6148,168.5035) -- (222.0701,168.5035) --
  (222.0701,175.9588) -- (214.6148,175.9588) -- cycle(222.3148,168.5035) --
  (229.7701,168.5035) -- (229.7701,175.9588) -- (222.3148,175.9588) --
  cycle(230.0148,168.5035) -- (237.4701,168.5035) -- (237.4701,175.9588) --
  (230.0148,175.9588) -- cycle(237.7148,168.5035) -- (245.1701,168.5035) --
  (245.1701,175.9588) -- (237.7148,175.9588) -- cycle(245.4148,168.5035) --
  (252.8701,168.5035) -- (252.8701,175.9588) -- (245.4148,175.9588) --
  cycle(37.5148,176.2035) -- (44.9701,176.2035) -- (44.9701,183.6588) --
  (37.5148,183.6588) -- cycle(45.2148,176.2035) -- (52.6701,176.2035) --
  (52.6701,183.6588) -- (45.2148,183.6588) -- cycle(52.9148,176.2035) --
  (60.3701,176.2035) -- (60.3701,183.6588) -- (52.9148,183.6588) --
  cycle(60.6148,176.2035) -- (68.0701,176.2035) -- (68.0701,183.6588) --
  (60.6148,183.6588) -- cycle(68.3148,176.2035) -- (75.7701,176.2035) --
  (75.7701,183.6588) -- (68.3148,183.6588) -- cycle(76.0148,176.2035) --
  (83.4701,176.2035) -- (83.4701,183.6588) -- (76.0148,183.6588) --
  cycle(83.7148,176.2035) -- (91.1701,176.2035) -- (91.1701,183.6588) --
  (83.7148,183.6588) -- cycle(91.4148,176.2035) -- (98.8701,176.2035) --
  (98.8701,183.6588) -- (91.4148,183.6588) -- cycle(99.1148,176.2035) --
  (106.5701,176.2035) -- (106.5701,183.6588) -- (99.1148,183.6588) --
  cycle(106.8148,176.2035) -- (114.2701,176.2035) -- (114.2701,183.6588) --
  (106.8148,183.6588) -- cycle(114.5148,176.2035) -- (121.9701,176.2035) --
  (121.9701,183.6588) -- (114.5148,183.6588) -- cycle(122.2148,176.2035) --
  (129.6701,176.2035) -- (129.6701,183.6588) -- (122.2148,183.6588) --
  cycle(129.9148,176.2035) -- (137.3701,176.2035) -- (137.3701,183.6588) --
  (129.9148,183.6588) -- cycle(137.6148,176.2035) -- (145.0701,176.2035) --
  (145.0701,183.6588) -- (137.6148,183.6588) -- cycle(145.3148,176.2035) --
  (152.7701,176.2035) -- (152.7701,183.6588) -- (145.3148,183.6588) --
  cycle(153.0148,176.2035) -- (160.4701,176.2035) -- (160.4701,183.6588) --
  (153.0148,183.6588) -- cycle(160.7148,176.2035) -- (168.1701,176.2035) --
  (168.1701,183.6588) -- (160.7148,183.6588) -- cycle(168.4148,176.2035) --
  (175.8701,176.2035) -- (175.8701,183.6588) -- (168.4148,183.6588) --
  cycle(176.1148,176.2035) -- (183.5701,176.2035) -- (183.5701,183.6588) --
  (176.1148,183.6588) -- cycle(183.8148,176.2035) -- (191.2701,176.2035) --
  (191.2701,183.6588) -- (183.8148,183.6588) -- cycle(191.5148,176.2035) --
  (198.9701,176.2035) -- (198.9701,183.6588) -- (191.5148,183.6588) --
  cycle(199.2148,176.2035) -- (206.6701,176.2035) -- (206.6701,183.6588) --
  (199.2148,183.6588) -- cycle(206.9148,176.2035) -- (214.3701,176.2035) --
  (214.3701,183.6588) -- (206.9148,183.6588) -- cycle(214.6148,176.2035) --
  (222.0701,176.2035) -- (222.0701,183.6588) -- (214.6148,183.6588) --
  cycle(222.3148,176.2035) -- (229.7701,176.2035) -- (229.7701,183.6588) --
  (222.3148,183.6588) -- cycle(230.0148,176.2035) -- (237.4701,176.2035) --
  (237.4701,183.6588) -- (230.0148,183.6588) -- cycle(237.7148,176.2035) --
  (245.1701,176.2035) -- (245.1701,183.6588) -- (237.7148,183.6588) --
  cycle(245.4148,176.2035) -- (252.8701,176.2035) -- (252.8701,183.6588) --
  (245.4148,183.6588) -- cycle(37.5148,183.9035) -- (44.9701,183.9035) --
  (44.9701,191.3588) -- (37.5148,191.3588) -- cycle(45.2148,183.9035) --
  (52.6701,183.9035) -- (52.6701,191.3588) -- (45.2148,191.3588) --
  cycle(52.9148,183.9035) -- (60.3701,183.9035) -- (60.3701,191.3588) --
  (52.9148,191.3588) -- cycle(60.6148,183.9035) -- (68.0701,183.9035) --
  (68.0701,191.3588) -- (60.6148,191.3588) -- cycle(68.3148,183.9035) --
  (75.7701,183.9035) -- (75.7701,191.3588) -- (68.3148,191.3588) --
  cycle(76.0148,183.9035) -- (83.4701,183.9035) -- (83.4701,191.3588) --
  (76.0148,191.3588) -- cycle(83.7148,183.9035) -- (91.1701,183.9035) --
  (91.1701,191.3588) -- (83.7148,191.3588) -- cycle(91.4148,183.9035) --
  (98.8701,183.9035) -- (98.8701,191.3588) -- (91.4148,191.3588) --
  cycle(99.1148,183.9035) -- (106.5701,183.9035) -- (106.5701,191.3588) --
  (99.1148,191.3588) -- cycle(106.8148,183.9035) -- (114.2701,183.9035) --
  (114.2701,191.3588) -- (106.8148,191.3588) -- cycle(114.5148,183.9035) --
  (121.9701,183.9035) -- (121.9701,191.3588) -- (114.5148,191.3588) --
  cycle(122.2148,183.9035) -- (129.6701,183.9035) -- (129.6701,191.3588) --
  (122.2148,191.3588) -- cycle(129.9148,183.9035) -- (137.3701,183.9035) --
  (137.3701,191.3588) -- (129.9148,191.3588) -- cycle(137.6148,183.9035) --
  (145.0701,183.9035) -- (145.0701,191.3588) -- (137.6148,191.3588) --
  cycle(145.3148,183.9035) -- (152.7701,183.9035) -- (152.7701,191.3588) --
  (145.3148,191.3588) -- cycle(153.0148,183.9035) -- (160.4701,183.9035) --
  (160.4701,191.3588) -- (153.0148,191.3588) -- cycle(160.7148,183.9035) --
  (168.1701,183.9035) -- (168.1701,191.3588) -- (160.7148,191.3588) --
  cycle(168.4148,183.9035) -- (175.8701,183.9035) -- (175.8701,191.3588) --
  (168.4148,191.3588) -- cycle(176.1148,183.9035) -- (183.5701,183.9035) --
  (183.5701,191.3588) -- (176.1148,191.3588) -- cycle(183.8148,183.9035) --
  (191.2701,183.9035) -- (191.2701,191.3588) -- (183.8148,191.3588) --
  cycle(191.5148,183.9035) -- (198.9701,183.9035) -- (198.9701,191.3588) --
  (191.5148,191.3588) -- cycle(199.2148,183.9035) -- (206.6701,183.9035) --
  (206.6701,191.3588) -- (199.2148,191.3588) -- cycle(206.9148,183.9035) --
  (214.3701,183.9035) -- (214.3701,191.3588) -- (206.9148,191.3588) --
  cycle(214.6148,183.9035) -- (222.0701,183.9035) -- (222.0701,191.3588) --
  (214.6148,191.3588) -- cycle(222.3148,183.9035) -- (229.7701,183.9035) --
  (229.7701,191.3588) -- (222.3148,191.3588) -- cycle(230.0148,183.9035) --
  (237.4701,183.9035) -- (237.4701,191.3588) -- (230.0148,191.3588) --
  cycle(237.7148,183.9035) -- (245.1701,183.9035) -- (245.1701,191.3588) --
  (237.7148,191.3588) -- cycle(45.2148,191.6035) -- (52.6701,191.6035) --
  (52.6701,199.0588) -- (45.2148,199.0588) -- cycle(52.9148,191.6035) --
  (60.3701,191.6035) -- (60.3701,199.0588) -- (52.9148,199.0588) --
  cycle(60.6148,191.6035) -- (68.0701,191.6035) -- (68.0701,199.0588) --
  (60.6148,199.0588) -- cycle(68.3148,191.6035) -- (75.7701,191.6035) --
  (75.7701,199.0588) -- (68.3148,199.0588) -- cycle(76.0148,191.6035) --
  (83.4701,191.6035) -- (83.4701,199.0588) -- (76.0148,199.0588) --
  cycle(83.7148,191.6035) -- (91.1701,191.6035) -- (91.1701,199.0588) --
  (83.7148,199.0588) -- cycle(91.4148,191.6035) -- (98.8701,191.6035) --
  (98.8701,199.0588) -- (91.4148,199.0588) -- cycle(99.1148,191.6035) --
  (106.5701,191.6035) -- (106.5701,199.0588) -- (99.1148,199.0588) --
  cycle(106.8148,191.6035) -- (114.2701,191.6035) -- (114.2701,199.0588) --
  (106.8148,199.0588) -- cycle(114.5148,191.6035) -- (121.9701,191.6035) --
  (121.9701,199.0588) -- (114.5148,199.0588) -- cycle(122.2148,191.6035) --
  (129.6701,191.6035) -- (129.6701,199.0588) -- (122.2148,199.0588) --
  cycle(129.9148,191.6035) -- (137.3701,191.6035) -- (137.3701,199.0588) --
  (129.9148,199.0588) -- cycle(137.6148,191.6035) -- (145.0701,191.6035) --
  (145.0701,199.0588) -- (137.6148,199.0588) -- cycle(145.3148,191.6035) --
  (152.7701,191.6035) -- (152.7701,199.0588) -- (145.3148,199.0588) --
  cycle(153.0148,191.6035) -- (160.4701,191.6035) -- (160.4701,199.0588) --
  (153.0148,199.0588) -- cycle(160.7148,191.6035) -- (168.1701,191.6035) --
  (168.1701,199.0588) -- (160.7148,199.0588) -- cycle(168.4148,191.6035) --
  (175.8701,191.6035) -- (175.8701,199.0588) -- (168.4148,199.0588) --
  cycle(176.1148,191.6035) -- (183.5701,191.6035) -- (183.5701,199.0588) --
  (176.1148,199.0588) -- cycle(183.8148,191.6035) -- (191.2701,191.6035) --
  (191.2701,199.0588) -- (183.8148,199.0588) -- cycle(191.5148,191.6035) --
  (198.9701,191.6035) -- (198.9701,199.0588) -- (191.5148,199.0588) --
  cycle(199.2148,191.6035) -- (206.6701,191.6035) -- (206.6701,199.0588) --
  (199.2148,199.0588) -- cycle(206.9148,191.6035) -- (214.3701,191.6035) --
  (214.3701,199.0588) -- (206.9148,199.0588) -- cycle(214.6148,191.6035) --
  (222.0701,191.6035) -- (222.0701,199.0588) -- (214.6148,199.0588) --
  cycle(222.3148,191.6035) -- (229.7701,191.6035) -- (229.7701,199.0588) --
  (222.3148,199.0588) -- cycle(230.0148,191.6035) -- (237.4701,191.6035) --
  (237.4701,199.0588) -- (230.0148,199.0588) -- cycle(237.7148,191.6035) --
  (245.1701,191.6035) -- (245.1701,199.0588) -- (237.7148,199.0588) --
  cycle(52.9148,199.3035) -- (60.3701,199.3035) -- (60.3701,206.7588) --
  (52.9148,206.7588) -- cycle(60.6148,199.3035) -- (68.0701,199.3035) --
  (68.0701,206.7588) -- (60.6148,206.7588) -- cycle(68.3148,199.3035) --
  (75.7701,199.3035) -- (75.7701,206.7588) -- (68.3148,206.7588) --
  cycle(76.0148,199.3035) -- (83.4701,199.3035) -- (83.4701,206.7588) --
  (76.0148,206.7588) -- cycle(83.7148,199.3035) -- (91.1701,199.3035) --
  (91.1701,206.7588) -- (83.7148,206.7588) -- cycle(91.4148,199.3035) --
  (98.8701,199.3035) -- (98.8701,206.7588) -- (91.4148,206.7588) --
  cycle(99.1148,199.3035) -- (106.5701,199.3035) -- (106.5701,206.7588) --
  (99.1148,206.7588) -- cycle(106.8148,199.3035) -- (114.2701,199.3035) --
  (114.2701,206.7588) -- (106.8148,206.7588) -- cycle(114.5148,199.3035) --
  (121.9701,199.3035) -- (121.9701,206.7588) -- (114.5148,206.7588) --
  cycle(122.2148,199.3035) -- (129.6701,199.3035) -- (129.6701,206.7588) --
  (122.2148,206.7588) -- cycle(129.9148,199.3035) -- (137.3701,199.3035) --
  (137.3701,206.7588) -- (129.9148,206.7588) -- cycle(137.6148,199.3035) --
  (145.0701,199.3035) -- (145.0701,206.7588) -- (137.6148,206.7588) --
  cycle(145.3148,199.3035) -- (152.7701,199.3035) -- (152.7701,206.7588) --
  (145.3148,206.7588) -- cycle(153.0148,199.3035) -- (160.4701,199.3035) --
  (160.4701,206.7588) -- (153.0148,206.7588) -- cycle(160.7148,199.3035) --
  (168.1701,199.3035) -- (168.1701,206.7588) -- (160.7148,206.7588) --
  cycle(168.4148,199.3035) -- (175.8701,199.3035) -- (175.8701,206.7588) --
  (168.4148,206.7588) -- cycle(176.1148,199.3035) -- (183.5701,199.3035) --
  (183.5701,206.7588) -- (176.1148,206.7588) -- cycle(183.8148,199.3035) --
  (191.2701,199.3035) -- (191.2701,206.7588) -- (183.8148,206.7588) --
  cycle(191.5148,199.3035) -- (198.9701,199.3035) -- (198.9701,206.7588) --
  (191.5148,206.7588) -- cycle(199.2148,199.3035) -- (206.6701,199.3035) --
  (206.6701,206.7588) -- (199.2148,206.7588) -- cycle(206.9148,199.3035) --
  (214.3701,199.3035) -- (214.3701,206.7588) -- (206.9148,206.7588) --
  cycle(214.6148,199.3035) -- (222.0701,199.3035) -- (222.0701,206.7588) --
  (214.6148,206.7588) -- cycle(222.3148,199.3035) -- (229.7701,199.3035) --
  (229.7701,206.7588) -- (222.3148,206.7588) -- cycle(230.0148,199.3035) --
  (237.4701,199.3035) -- (237.4701,206.7588) -- (230.0148,206.7588) --
  cycle(52.9148,207.0035) -- (60.3701,207.0035) -- (60.3701,214.4588) --
  (52.9148,214.4588) -- cycle(60.6148,207.0035) -- (68.0701,207.0035) --
  (68.0701,214.4588) -- (60.6148,214.4588) -- cycle(68.3148,207.0035) --
  (75.7701,207.0035) -- (75.7701,214.4588) -- (68.3148,214.4588) --
  cycle(76.0148,207.0035) -- (83.4701,207.0035) -- (83.4701,214.4588) --
  (76.0148,214.4588) -- cycle(83.7148,207.0035) -- (91.1701,207.0035) --
  (91.1701,214.4588) -- (83.7148,214.4588) -- cycle(91.4148,207.0035) --
  (98.8701,207.0035) -- (98.8701,214.4588) -- (91.4148,214.4588) --
  cycle(99.1148,207.0035) -- (106.5701,207.0035) -- (106.5701,214.4588) --
  (99.1148,214.4588) -- cycle(106.8148,207.0035) -- (114.2701,207.0035) --
  (114.2701,214.4588) -- (106.8148,214.4588) -- cycle(114.5148,207.0035) --
  (121.9701,207.0035) -- (121.9701,214.4588) -- (114.5148,214.4588) --
  cycle(122.2148,207.0035) -- (129.6701,207.0035) -- (129.6701,214.4588) --
  (122.2148,214.4588) -- cycle(129.9148,207.0035) -- (137.3701,207.0035) --
  (137.3701,214.4588) -- (129.9148,214.4588) -- cycle(137.6148,207.0035) --
  (145.0701,207.0035) -- (145.0701,214.4588) -- (137.6148,214.4588) --
  cycle(145.3148,207.0035) -- (152.7701,207.0035) -- (152.7701,214.4588) --
  (145.3148,214.4588) -- cycle(153.0148,207.0035) -- (160.4701,207.0035) --
  (160.4701,214.4588) -- (153.0148,214.4588) -- cycle(160.7148,207.0035) --
  (168.1701,207.0035) -- (168.1701,214.4588) -- (160.7148,214.4588) --
  cycle(168.4148,207.0035) -- (175.8701,207.0035) -- (175.8701,214.4588) --
  (168.4148,214.4588) -- cycle(176.1148,207.0035) -- (183.5701,207.0035) --
  (183.5701,214.4588) -- (176.1148,214.4588) -- cycle(183.8148,207.0035) --
  (191.2701,207.0035) -- (191.2701,214.4588) -- (183.8148,214.4588) --
  cycle(191.5148,207.0035) -- (198.9701,207.0035) -- (198.9701,214.4588) --
  (191.5148,214.4588) -- cycle(199.2148,207.0035) -- (206.6701,207.0035) --
  (206.6701,214.4588) -- (199.2148,214.4588) -- cycle(206.9148,207.0035) --
  (214.3701,207.0035) -- (214.3701,214.4588) -- (206.9148,214.4588) --
  cycle(214.6148,207.0035) -- (222.0701,207.0035) -- (222.0701,214.4588) --
  (214.6148,214.4588) -- cycle(222.3148,207.0035) -- (229.7701,207.0035) --
  (229.7701,214.4588) -- (222.3148,214.4588) -- cycle(60.6148,214.7035) --
  (68.0701,214.7035) -- (68.0701,222.1588) -- (60.6148,222.1588) --
  cycle(68.3148,214.7035) -- (75.7701,214.7035) -- (75.7701,222.1588) --
  (68.3148,222.1588) -- cycle(76.0148,214.7035) -- (83.4701,214.7035) --
  (83.4701,222.1588) -- (76.0148,222.1588) -- cycle(83.7148,214.7035) --
  (91.1701,214.7035) -- (91.1701,222.1588) -- (83.7148,222.1588) --
  cycle(91.4148,214.7035) -- (98.8701,214.7035) -- (98.8701,222.1588) --
  (91.4148,222.1588) -- cycle(99.1148,214.7035) -- (106.5701,214.7035) --
  (106.5701,222.1588) -- (99.1148,222.1588) -- cycle(106.8148,214.7035) --
  (114.2701,214.7035) -- (114.2701,222.1588) -- (106.8148,222.1588) --
  cycle(114.5148,214.7035) -- (121.9701,214.7035) -- (121.9701,222.1588) --
  (114.5148,222.1588) -- cycle(122.2148,214.7035) -- (129.6701,214.7035) --
  (129.6701,222.1588) -- (122.2148,222.1588) -- cycle(129.9148,214.7035) --
  (137.3701,214.7035) -- (137.3701,222.1588) -- (129.9148,222.1588) --
  cycle(137.6148,214.7035) -- (145.0701,214.7035) -- (145.0701,222.1588) --
  (137.6148,222.1588) -- cycle(145.3148,214.7035) -- (152.7701,214.7035) --
  (152.7701,222.1588) -- (145.3148,222.1588) -- cycle(153.0148,214.7035) --
  (160.4701,214.7035) -- (160.4701,222.1588) -- (153.0148,222.1588) --
  cycle(160.7148,214.7035) -- (168.1701,214.7035) -- (168.1701,222.1588) --
  (160.7148,222.1588) -- cycle(168.4148,214.7035) -- (175.8701,214.7035) --
  (175.8701,222.1588) -- (168.4148,222.1588) -- cycle(176.1148,214.7035) --
  (183.5701,214.7035) -- (183.5701,222.1588) -- (176.1148,222.1588) --
  cycle(183.8148,214.7035) -- (191.2701,214.7035) -- (191.2701,222.1588) --
  (183.8148,222.1588) -- cycle(191.5148,214.7035) -- (198.9701,214.7035) --
  (198.9701,222.1588) -- (191.5148,222.1588) -- cycle(199.2148,214.7035) --
  (206.6701,214.7035) -- (206.6701,222.1588) -- (199.2148,222.1588) --
  cycle(206.9148,214.7035) -- (214.3701,214.7035) -- (214.3701,222.1588) --
  (206.9148,222.1588) -- cycle(214.6148,214.7035) -- (222.0701,214.7035) --
  (222.0701,222.1588) -- (214.6148,222.1588) -- cycle(222.3148,214.7035) --
  (229.7701,214.7035) -- (229.7701,222.1588) -- (222.3148,222.1588) --
  cycle(60.6148,222.4035) -- (68.0701,222.4035) -- (68.0701,229.8588) --
  (60.6148,229.8588) -- cycle(68.3148,222.4035) -- (75.7701,222.4035) --
  (75.7701,229.8588) -- (68.3148,229.8588) -- cycle(76.0148,222.4035) --
  (83.4701,222.4035) -- (83.4701,229.8588) -- (76.0148,229.8588) --
  cycle(83.7148,222.4035) -- (91.1701,222.4035) -- (91.1701,229.8588) --
  (83.7148,229.8588) -- cycle(91.4148,222.4035) -- (98.8701,222.4035) --
  (98.8701,229.8588) -- (91.4148,229.8588) -- cycle(99.1148,222.4035) --
  (106.5701,222.4035) -- (106.5701,229.8588) -- (99.1148,229.8588) --
  cycle(106.8148,222.4035) -- (114.2701,222.4035) -- (114.2701,229.8588) --
  (106.8148,229.8588) -- cycle(114.5148,222.4035) -- (121.9701,222.4035) --
  (121.9701,229.8588) -- (114.5148,229.8588) -- cycle(122.2148,222.4035) --
  (129.6701,222.4035) -- (129.6701,229.8588) -- (122.2148,229.8588) --
  cycle(129.9148,222.4035) -- (137.3701,222.4035) -- (137.3701,229.8588) --
  (129.9148,229.8588) -- cycle(137.6148,222.4035) -- (145.0701,222.4035) --
  (145.0701,229.8588) -- (137.6148,229.8588) -- cycle(145.3148,222.4035) --
  (152.7701,222.4035) -- (152.7701,229.8588) -- (145.3148,229.8588) --
  cycle(153.0148,222.4035) -- (160.4701,222.4035) -- (160.4701,229.8588) --
  (153.0148,229.8588) -- cycle(160.7148,222.4035) -- (168.1701,222.4035) --
  (168.1701,229.8588) -- (160.7148,229.8588) -- cycle(168.4148,222.4035) --
  (175.8701,222.4035) -- (175.8701,229.8588) -- (168.4148,229.8588) --
  cycle(176.1148,222.4035) -- (183.5701,222.4035) -- (183.5701,229.8588) --
  (176.1148,229.8588) -- cycle(183.8148,222.4035) -- (191.2701,222.4035) --
  (191.2701,229.8588) -- (183.8148,229.8588) -- cycle(191.5148,222.4035) --
  (198.9701,222.4035) -- (198.9701,229.8588) -- (191.5148,229.8588) --
  cycle(199.2148,222.4035) -- (206.6701,222.4035) -- (206.6701,229.8588) --
  (199.2148,229.8588) -- cycle(206.9148,222.4035) -- (214.3701,222.4035) --
  (214.3701,229.8588) -- (206.9148,229.8588) -- cycle(214.6148,222.4035) --
  (222.0701,222.4035) -- (222.0701,229.8588) -- (214.6148,229.8588) --
  cycle(68.3148,230.1035) -- (75.7701,230.1035) -- (75.7701,237.5588) --
  (68.3148,237.5588) -- cycle(76.0148,230.1035) -- (83.4701,230.1035) --
  (83.4701,237.5588) -- (76.0148,237.5588) -- cycle(83.7148,230.1035) --
  (91.1701,230.1035) -- (91.1701,237.5588) -- (83.7148,237.5588) --
  cycle(91.4148,230.1035) -- (98.8701,230.1035) -- (98.8701,237.5588) --
  (91.4148,237.5588) -- cycle(99.1148,230.1035) -- (106.5701,230.1035) --
  (106.5701,237.5588) -- (99.1148,237.5588) -- cycle(106.8148,230.1035) --
  (114.2701,230.1035) -- (114.2701,237.5588) -- (106.8148,237.5588) --
  cycle(114.5148,230.1035) -- (121.9701,230.1035) -- (121.9701,237.5588) --
  (114.5148,237.5588) -- cycle(122.2148,230.1035) -- (129.6701,230.1035) --
  (129.6701,237.5588) -- (122.2148,237.5588) -- cycle(129.9148,230.1035) --
  (137.3701,230.1035) -- (137.3701,237.5588) -- (129.9148,237.5588) --
  cycle(137.6148,230.1035) -- (145.0701,230.1035) -- (145.0701,237.5588) --
  (137.6148,237.5588) -- cycle(145.3148,230.1035) -- (152.7701,230.1035) --
  (152.7701,237.5588) -- (145.3148,237.5588) -- cycle(153.0148,230.1035) --
  (160.4701,230.1035) -- (160.4701,237.5588) -- (153.0148,237.5588) --
  cycle(160.7148,230.1035) -- (168.1701,230.1035) -- (168.1701,237.5588) --
  (160.7148,237.5588) -- cycle(168.4148,230.1035) -- (175.8701,230.1035) --
  (175.8701,237.5588) -- (168.4148,237.5588) -- cycle(176.1148,230.1035) --
  (183.5701,230.1035) -- (183.5701,237.5588) -- (176.1148,237.5588) --
  cycle(183.8148,230.1035) -- (191.2701,230.1035) -- (191.2701,237.5588) --
  (183.8148,237.5588) -- cycle(191.5148,230.1035) -- (198.9701,230.1035) --
  (198.9701,237.5588) -- (191.5148,237.5588) -- cycle(199.2148,230.1035) --
  (206.6701,230.1035) -- (206.6701,237.5588) -- (199.2148,237.5588) --
  cycle(206.9148,230.1035) -- (214.3701,230.1035) -- (214.3701,237.5588) --
  (206.9148,237.5588) -- cycle(214.6148,230.1035) -- (222.0701,230.1035) --
  (222.0701,237.5588) -- (214.6148,237.5588) -- cycle(76.0148,237.8035) --
  (83.4701,237.8035) -- (83.4701,245.2588) -- (76.0148,245.2588) --
  cycle(83.7148,237.8035) -- (91.1701,237.8035) -- (91.1701,245.2588) --
  (83.7148,245.2588) -- cycle(91.4148,237.8035) -- (98.8701,237.8035) --
  (98.8701,245.2588) -- (91.4148,245.2588) -- cycle(99.1148,237.8035) --
  (106.5701,237.8035) -- (106.5701,245.2588) -- (99.1148,245.2588) --
  cycle(106.8148,237.8035) -- (114.2701,237.8035) -- (114.2701,245.2588) --
  (106.8148,245.2588) -- cycle(114.5148,237.8035) -- (121.9701,237.8035) --
  (121.9701,245.2588) -- (114.5148,245.2588) -- cycle(122.2148,237.8035) --
  (129.6701,237.8035) -- (129.6701,245.2588) -- (122.2148,245.2588) --
  cycle(129.9148,237.8035) -- (137.3701,237.8035) -- (137.3701,245.2588) --
  (129.9148,245.2588) -- cycle(137.6148,237.8035) -- (145.0701,237.8035) --
  (145.0701,245.2588) -- (137.6148,245.2588) -- cycle(145.3148,237.8035) --
  (152.7701,237.8035) -- (152.7701,245.2588) -- (145.3148,245.2588) --
  cycle(153.0148,237.8035) -- (160.4701,237.8035) -- (160.4701,245.2588) --
  (153.0148,245.2588) -- cycle(160.7148,237.8035) -- (168.1701,237.8035) --
  (168.1701,245.2588) -- (160.7148,245.2588) -- cycle(168.4148,237.8035) --
  (175.8701,237.8035) -- (175.8701,245.2588) -- (168.4148,245.2588) --
  cycle(176.1148,237.8035) -- (183.5701,237.8035) -- (183.5701,245.2588) --
  (176.1148,245.2588) -- cycle(183.8148,237.8035) -- (191.2701,237.8035) --
  (191.2701,245.2588) -- (183.8148,245.2588) -- cycle(191.5148,237.8035) --
  (198.9701,237.8035) -- (198.9701,245.2588) -- (191.5148,245.2588) --
  cycle(199.2148,237.8035) -- (206.6701,237.8035) -- (206.6701,245.2588) --
  (199.2148,245.2588) -- cycle(206.9148,237.8035) -- (214.3701,237.8035) --
  (214.3701,245.2588) -- (206.9148,245.2588) -- cycle(83.7148,245.5035) --
  (91.1701,245.5035) -- (91.1701,252.9588) -- (83.7148,252.9588) --
  cycle(91.4148,245.5035) -- (98.8701,245.5035) -- (98.8701,252.9588) --
  (91.4148,252.9588) -- cycle(99.1148,245.5035) -- (106.5701,245.5035) --
  (106.5701,252.9588) -- (99.1148,252.9588) -- cycle(106.8148,245.5035) --
  (114.2701,245.5035) -- (114.2701,252.9588) -- (106.8148,252.9588) --
  cycle(114.5148,245.5035) -- (121.9701,245.5035) -- (121.9701,252.9588) --
  (114.5148,252.9588) -- cycle(122.2148,245.5035) -- (129.6701,245.5035) --
  (129.6701,252.9588) -- (122.2148,252.9588) -- cycle(129.9148,245.5035) --
  (137.3701,245.5035) -- (137.3701,252.9588) -- (129.9148,252.9588) --
  cycle(137.6148,245.5035) -- (145.0701,245.5035) -- (145.0701,252.9588) --
  (137.6148,252.9588) -- cycle(145.3148,245.5035) -- (152.7701,245.5035) --
  (152.7701,252.9588) -- (145.3148,252.9588) -- cycle(153.0148,245.5035) --
  (160.4701,245.5035) -- (160.4701,252.9588) -- (153.0148,252.9588) --
  cycle(160.7148,245.5035) -- (168.1701,245.5035) -- (168.1701,252.9588) --
  (160.7148,252.9588) -- cycle(168.4148,245.5035) -- (175.8701,245.5035) --
  (175.8701,252.9588) -- (168.4148,252.9588) -- cycle(176.1148,245.5035) --
  (183.5701,245.5035) -- (183.5701,252.9588) -- (176.1148,252.9588) --
  cycle(183.8148,245.5035) -- (191.2701,245.5035) -- (191.2701,252.9588) --
  (183.8148,252.9588) -- cycle;\path[fill=cffffff,line width=0.026pt] (155.0404,73.7884) .. controls
  (155.0404,90.4175) and (155.0054,107.0466) .. (155.0404,123.7107) .. controls
  (171.7746,123.6406) and (188.5087,123.6757) .. (205.2078,123.6757) .. controls
  (205.2428,113.8383) and (205.2428,94.1284) .. (205.2428,94.1284) --
  (202.3721,97.3492) .. controls (202.3721,97.3492) and (202.4071,112.9280) ..
  (202.3021,120.7350) .. controls (187.5284,120.7700) and (172.7198,120.7700) ..
  (157.9461,120.7350) .. controls (157.8411,112.3329) and (157.9811,103.9308) ..
  (157.8761,95.5287) .. controls (157.5610,93.6383) and (157.4210,91.7128) ..
  (157.8411,89.7873) .. controls (158.0162,85.4463) and (157.8411,81.0702) ..
  (157.9111,76.7291) .. controls (169.8841,76.6241) and (182.5135,76.6941) ..
  (194.4514,76.6941) -- (197.6809,73.7884) .. controls (183.3314,73.7174) and
  (168.2406,73.6447) .. (155.0404,73.7884)(201.3569,73.2983) .. controls
  (194.3551,80.3000) and (187.3884,87.3717) .. (180.3867,94.4085) .. controls
  (177.5510,91.6778) and (174.8903,88.8071) .. (172.0546,86.0764) .. controls
  (169.8141,84.0459) and (166.1382,83.8709) .. (163.7926,85.7964) .. controls
  (159.9060,89.1651) and (157.9667,93.4351) .. (161.8321,97.3492) .. controls
  (166.3482,101.9353) and (170.8993,106.4865) .. (175.4505,111.0376) .. controls
  (179.8290,114.1880) and (183.4091,113.0674) .. (186.6182,109.6723) .. controls
  (194.5302,101.7253) and (202.4771,93.7783) .. (210.3891,85.8314) .. controls
  (211.5094,84.6411) and (212.8747,83.5208) .. (213.5049,81.9454) .. controls
  (214.4851,79.7048) and (214.0650,76.9392) .. (212.4196,75.1537) .. controls
  (209.2917,71.5843) and (204.8904,69.8555) ..
  (201.3569,73.2983)(141.0720,37.9396) .. controls (135.4706,38.4647) and
  (130.8844,43.7160) .. (130.9895,49.3174) .. controls (130.9545,80.7202) and
  (130.9895,112.0879) .. (130.9895,143.4907) .. controls (130.9215,150.1629) and
  (136.4866,155.5595) .. (142.2973,155.6387) -- (217.9860,155.6387) .. controls
  (227.0791,154.5830) and (229.3894,147.7597) .. (229.2938,140.3049) .. controls
  (229.2588,110.0924) and (229.2938,79.8449) .. (229.2938,49.6325) .. controls
  (228.7508,43.0410) and (224.5355,38.1300) .. (218.3011,37.8696) .. controls
  (215.5354,37.8346) and (212.7697,37.8696) .. (210.0390,37.8696) --
  (210.0040,45.4664) .. controls (210.0040,45.4664) and (215.2903,45.4664) ..
  (217.9160,45.5365) .. controls (220.0515,45.5015) and (221.8019,47.5320) ..
  (221.6619,49.6326) -- (221.6619,143.8058) .. controls (221.4372,146.9180) and
  (219.7073,148.0013) .. (216.9007,148.0069) -- (142.6474,148.0069) .. controls
  (140.5468,148.1119) and (138.5514,146.2914) .. (138.6214,144.1559) --
  (138.7264,48.5823) .. controls (139.0065,46.8669) and (140.5819,45.5365) ..
  (142.3323,45.5365) .. controls (144.8179,45.4664) and (147.3035,45.6066) ..
  (149.7891,45.5015) .. controls (149.7891,42.9459) and (149.8942,40.3903) ..
  (149.7191,37.8346) .. controls (146.8484,37.9047) and (143.9427,37.7645) ..
  (141.0720,37.9397)(167.6786,25.9667) .. controls (166.3132,26.4918) and
  (166.2432,28.1022) .. (166.2082,29.3275) .. controls (166.1732,32.0582) and
  (163.8976,34.5088) .. (161.2020,34.8939) .. controls (159.3815,35.1390) and
  (157.4911,34.8939) .. (155.6356,34.9989) .. controls (154.5853,35.0339) and
  (153.6401,35.9792) .. (153.5701,37.0294) -- (153.6051,47.8471) .. controls
  (153.7451,48.9674) and (154.8654,49.8776) .. (156.0207,49.8076) --
  (203.5974,49.8076) .. controls (204.4026,49.8426) and (205.2078,49.5975) ..
  (205.6980,48.9324) .. controls (206.3981,48.1972) and (206.2581,47.0769) ..
  (206.2931,46.1317) .. controls (206.2581,43.2260) and (206.2931,40.3202) ..
  (206.2581,37.3795) .. controls (206.3281,36.1192) and (205.2079,34.9639) ..
  (203.9475,34.9989) .. controls (201.4969,34.8239) and (198.8713,35.4190) ..
  (196.6307,34.2287) .. controls (194.6002,33.1435) and (193.4450,30.8679) ..
  (193.5500,28.6274) .. controls (193.7250,27.1570) and (192.5347,25.6866) ..
  (190.9943,25.7566) .. controls (183.9926,25.7216) and (176.9909,25.7566) ..
  (169.9892,25.7216) .. controls (169.2190,25.7566) and (168.4138,25.7216) ..
  (167.6786,25.9667);\begin{scope}[fill=cffffff,line width=1.108pt]
\path[fill=cffffff,line width=1.108pt] (86.1058,134.7520) .. controls
  (86.1058,140.5883) and (83.9818,145.4458) .. (79.7339,149.3243) --
  (85.6625,156.3057) -- (79.0136,162.3452) -- (71.5889,153.6461) .. controls
  (70.0005,154.1263) and (68.1721,154.3664) .. (66.1035,154.3664) .. controls
  (60.4149,154.3664) and (55.6683,152.4826) .. (51.8636,148.7148) .. controls
  (48.0590,144.9101) and (46.1566,140.2559) .. (46.1566,134.7520) .. controls
  (46.1566,129.2482) and (48.0590,124.6123) .. (51.8636,120.8446) .. controls
  (55.6683,117.0399) and (60.4149,115.1376) .. (66.1035,115.1376) .. controls
  (71.7921,115.1376) and (76.5387,117.0399) .. (80.3434,120.8446) .. controls
  (84.1850,124.6123) and (86.1058,129.2482) .. (86.1058,134.7520) --
  cycle(58.6234,127.0503) .. controls (56.6287,129.1558) and (55.6314,131.7230)
  .. (55.6314,134.7520) .. controls (55.6314,137.7810) and (56.6287,140.3667) ..
  (58.6234,142.5091) .. controls (60.6181,144.6146) and (63.1115,145.6674) ..
  (66.1035,145.6674) .. controls (69.0955,145.6674) and (71.5889,144.6146) ..
  (73.5836,142.5091) .. controls (75.5783,140.3667) and (76.5756,137.7810) ..
  (76.5756,134.7520) .. controls (76.5756,131.7230) and (75.5783,129.1558) ..
  (73.5836,127.0503) .. controls (71.5889,124.9079) and (69.0955,123.8366) ..
  (66.1035,123.8366) .. controls (63.1115,123.8366) and (60.6181,124.9079) ..
  (58.6234,127.0503) -- cycle;\path[fill=cffffff,line width=1.108pt] (114.5803,147.1080) -- (98.6228,147.1080)
  -- (96.1294,153.5907) -- (86.1560,153.5907) -- (102.1689,115.6917) --
  (111.1450,115.6917) -- (127.1025,153.5907) -- (117.1291,153.5907) --
  cycle(111.4774,139.1292) -- (106.6015,126.6070) -- (101.7256,139.1292) --
  cycle;\end{scope}
\begin{scope}[fill=cffffff,line width=0.897pt]
\path[fill,line width=0.897pt] (66.3104,192.7883) .. controls (68.2844,192.7883)
  and (69.9294,192.2499) .. (71.2454,191.1732) .. controls (72.5614,190.0965)
  and (73.4288,188.5711) .. (73.8475,186.5971) -- (81.6538,186.5971) .. controls
  (81.1753,190.5451) and (79.5303,193.7603) .. (76.7188,196.2428) .. controls
  (73.9372,198.6654) and (70.4827,199.8768) .. (66.3553,199.8768) .. controls
  (61.7493,199.8768) and (57.9059,198.3813) .. (54.8253,195.3904) .. controls
  (51.7446,192.3995) and (50.2043,188.6010) .. (50.2043,183.9950) .. controls
  (50.2043,179.4189) and (51.7446,175.6354) .. (54.8253,172.6445) .. controls
  (57.9358,169.6236) and (61.7792,168.1132) .. (66.3553,168.1132) .. controls
  (70.5126,168.1132) and (73.9672,169.3395) .. (76.7188,171.7920) .. controls
  (79.5303,174.2446) and (81.1753,177.4449) .. (81.6538,181.3929) --
  (73.8475,181.3929) .. controls (73.4288,179.4189) and (72.5614,177.8935) ..
  (71.2454,176.8168) .. controls (69.9294,175.7401) and (68.2844,175.2017) ..
  (66.3104,175.2017) .. controls (63.8279,175.2017) and (61.7941,176.0391) ..
  (60.2089,177.7141) .. controls (58.6537,179.3591) and (57.8760,181.4527) ..
  (57.8760,183.9950) .. controls (57.8760,186.5672) and (58.6537,188.6758) ..
  (60.2089,190.3208) .. controls (61.7642,191.9658) and (63.7980,192.7883) ..
  (66.3104,192.7883) -- cycle;\path[fill,line width=0.897pt] (113.8688,168.7413) -- (113.8688,199.2487) --
  (106.3317,199.2487) -- (106.3317,187.8084) -- (93.9044,187.8084) --
  (93.9044,199.2487) -- (86.3225,199.2487) -- (86.3225,168.7413) --
  (93.9044,168.7413) -- (93.9044,180.9891) -- (106.3317,180.9891) --
  (106.3317,168.7413) -- cycle;\path[fill,line width=0.897pt] (140.1309,193.9996) -- (127.2102,193.9996) --
  (125.1913,199.2487) -- (117.1158,199.2487) -- (130.0815,168.5618) --
  (137.3494,168.5618) -- (150.2702,199.2487) -- (142.1947,199.2487) --
  cycle(137.6186,187.5392) -- (133.6706,177.4000) -- (129.7225,187.5392) --
  cycle;\path[fill,line width=0.897pt] (165.8232,190.4105) -- (161.1125,190.4105) --
  (161.1125,199.2487) -- (153.5305,199.2487) -- (153.5305,168.7413) --
  (166.8550,168.7413) .. controls (170.6834,168.7413) and (173.7491,169.7133) ..
  (176.0521,171.6574) .. controls (178.3551,173.6015) and (179.5066,176.2036) ..
  (179.5066,179.4637) .. controls (179.5066,181.5574) and (178.9982,183.4267) ..
  (177.9812,185.0717) .. controls (176.9643,186.7167) and (175.5287,188.0028) ..
  (173.6743,188.9300) -- (181.5255,199.2487) -- (172.4181,199.2487) --
  cycle(161.1125,183.5912) -- (166.7204,183.5912) .. controls
  (168.2458,183.5912) and (169.4870,183.2173) .. (170.4441,182.4696) .. controls
  (171.4311,181.6920) and (171.9246,180.6900) .. (171.9246,179.4638) .. controls
  (171.9246,178.2674) and (171.4311,177.2953) .. (170.4441,176.5476) .. controls
  (169.4571,175.7999) and (168.2159,175.4260) .. (166.7204,175.4260) --
  (161.1125,175.4260) -- cycle;\path[fill,line width=0.897pt] (208.5833,175.5157) -- (198.8478,175.5157) --
  (198.8478,199.2487) -- (191.3107,199.2487) -- (191.3107,175.5157) --
  (181.5753,175.5157) -- (181.5753,168.7413) -- (208.5833,168.7413) -- cycle;\path[fill,line width=0.897pt] (223.4395,168.0684) .. controls
  (226.2808,168.0684) and (228.9577,168.7263) .. (231.4701,170.0423) .. controls
  (232.8459,170.7901) and (233.9825,171.8668) .. (234.8797,173.2725) .. controls
  (235.7770,174.6783) and (236.3154,176.3233) .. (236.4948,178.2076) --
  (228.8680,178.2076) .. controls (228.4193,175.6354) and (226.7743,174.3493) ..
  (223.9330,174.3493) .. controls (222.5871,174.3493) and (221.5253,174.6185) ..
  (220.7476,175.1568) .. controls (219.9700,175.6952) and (219.5812,176.4429) ..
  (219.5812,177.4000) .. controls (219.5812,178.1776) and (219.8504,178.7908) ..
  (220.3887,179.2394) .. controls (220.9570,179.6582) and (221.8543,180.0021) ..
  (223.0806,180.2713) -- (227.3426,181.2583) .. controls (230.4831,181.9761) and
  (232.8010,183.0080) .. (234.2965,184.3539) .. controls (235.7920,185.6699) and
  (236.5397,187.5392) .. (236.5397,189.9619) .. controls (236.5397,193.0126) and
  (235.3882,195.4353) .. (233.0852,197.2298) .. controls (230.7822,198.9944) and
  (227.7165,199.8768) .. (223.8881,199.8768) .. controls (220.1495,199.8768) and
  (217.0090,198.9795) .. (214.4667,197.1849) .. controls (211.8646,195.3306) and
  (210.4140,192.6387) .. (210.1149,189.1095) -- (217.6969,189.1095) .. controls
  (217.9362,190.4554) and (218.6091,191.5171) .. (219.7158,192.2948) .. controls
  (220.8523,193.0425) and (222.3179,193.4164) .. (224.1124,193.4164) .. controls
  (225.5182,193.4164) and (226.5799,193.1771) .. (227.2977,192.6986) .. controls
  (228.0455,192.2200) and (228.4193,191.5620) .. (228.4193,190.7245) .. controls
  (228.4193,189.4684) and (227.3426,188.6010) .. (225.1891,188.1225) --
  (220.1644,187.0009) .. controls (217.1436,186.3129) and (214.9153,185.2512) ..
  (213.4797,183.8155) .. controls (212.0441,182.3799) and (211.3262,180.4507) ..
  (211.3262,178.0281) .. controls (211.3262,175.0073) and (212.4329,172.5996) ..
  (214.6462,170.8050) .. controls (216.8594,168.9806) and (219.7905,168.0683) ..
  (223.4395,168.0683) -- cycle;\end{scope}
\begin{scope}[fill=cffffff,line width=0.624pt]
\path[fill=cffffff,line width=0.624pt] (98.3573,212.4443) -- (98.3573,233.6628)
  -- (93.0839,233.6628) -- (93.0839,212.4443) -- cycle;\path[fill=cffffff,line width=0.624pt] (121.6060,212.4443) --
  (121.6060,233.6628) -- (117.5495,233.6628) -- (107.3771,222.0239) --
  (107.3771,233.6628) -- (102.1037,233.6628) -- (102.1037,212.4443) --
  (106.1290,212.4443) -- (116.3638,224.3954) -- (116.3638,212.4443) -- cycle;\path[fill=cffffff,line width=0.624pt] (125.3236,212.4443) --
  (133.8422,212.4443) .. controls (137.1706,212.4443) and (139.9165,213.4325) ..
  (142.0800,215.4087) .. controls (144.2226,217.3641) and (145.2939,219.8916) ..
  (145.2939,222.9912) .. controls (145.2939,226.1324) and (144.2226,228.7015) ..
  (142.0800,230.6985) .. controls (139.9165,232.6747) and (137.1706,233.6628) ..
  (133.8422,233.6628) -- (125.3236,233.6628) -- cycle(130.5970,228.9199) --
  (133.8422,228.9199) .. controls (135.5688,228.9199) and (137.0146,228.3582) ..
  (138.1795,227.2349) .. controls (139.3445,226.0699) and (139.9269,224.6554) ..
  (139.9269,222.9912) .. controls (139.9269,221.3478) and (139.3445,219.9748) ..
  (138.1795,218.8723) .. controls (137.0354,217.7490) and (135.5896,217.1873) ..
  (133.8422,217.1873) -- (130.5970,217.1873) -- cycle;\path[fill=cffffff,line width=0.624pt] (153.5732,217.0313) --
  (153.5732,220.9006) -- (162.9967,220.9006) -- (162.9967,225.3939) --
  (153.5732,225.3939) -- (153.5732,229.0759) -- (164.3384,229.0759) --
  (164.3384,233.6628) -- (148.2997,233.6628) -- (148.2997,212.4444) --
  (164.3384,212.4444) -- (164.3384,217.0313) -- cycle;\path[fill=cffffff,line width=0.624pt] (180.8881,233.6628) --
  (176.3323,227.1413) -- (171.7454,233.6628) -- (165.6919,233.6628) --
  (172.7751,223.2096) -- (165.7855,212.4444) -- (172.1198,212.4444) --
  (176.3011,218.9971) -- (180.4824,212.4444) -- (186.8167,212.4444) --
  (179.8271,223.2096) -- (186.9104,233.6628) -- cycle;\path[fill=cffffff,line width=0.624pt] (190.7894,223.6153) .. controls
  (190.0197,223.6153) and (189.3540,223.3552) .. (188.7923,222.8352) .. controls
  (188.2515,222.3151) and (187.9810,221.6598) .. (187.9810,220.8693) .. controls
  (187.9810,220.1205) and (188.2515,219.4860) .. (188.7923,218.9659) .. controls
  (189.3540,218.4459) and (190.0197,218.1858) .. (190.7894,218.1858) .. controls
  (191.5799,218.1858) and (192.2559,218.4459) .. (192.8176,218.9659) .. controls
  (193.3793,219.4860) and (193.6601,220.1205) .. (193.6601,220.8693) .. controls
  (193.6601,221.6598) and (193.3793,222.3151) .. (192.8176,222.8352) .. controls
  (192.2767,223.3552) and (191.6007,223.6153) .. (190.7894,223.6153) --
  cycle(192.8176,233.4444) .. controls (192.2767,233.9645) and
  (191.6007,234.2245) .. (190.7894,234.2245) .. controls (189.9989,234.2245) and
  (189.3332,233.9645) .. (188.7923,233.4444) .. controls (188.2307,232.9036) and
  (187.9498,232.2483) .. (187.9498,231.4786) .. controls (187.9706,230.7089) and
  (188.2515,230.0640) .. (188.7923,229.5440) .. controls (189.3540,229.0239) and
  (190.0301,228.7639) .. (190.8206,228.7639) .. controls (191.6111,228.7639) and
  (192.2767,229.0239) .. (192.8176,229.5440) .. controls (193.3793,230.0640) and
  (193.6601,230.7089) .. (193.6601,231.4786) .. controls (193.6601,232.2483) and
  (193.3793,232.9036) .. (192.8176,233.4444) -- cycle;\end{scope}

\end{tikzpicture}
\\
\end{tabularx}
\Large\printcontents[section]{}{1}{}
\vspace*{\fill}
%---------------------------%



%-----Begin Embedded Graphs-----%
\newpage %page style1 doesnt work without this newpage for some reason
\ClearShipoutPicture
\pagestyle{style2}
\setlength\headheight{1.5em}
\renewcommand{\headrulewidth}{0pt}
\include{graph}
%----------------------------------%



%----------Pgfplots templates, uncomment for testing----------%
%\LARGE
%\begin{tikzpicture}% Line graph, use for data over time
%	\begin{axis}
%		[ymin=0,
%		title=\fontsize{25}{25}\selectfont\color{pfgrey}{Warranties Claims - Monthly},
%		tick label style={/pgf/number format/assume math mode},
%		every axis plot/.append style={ultra thick},
%		ymajorgrids,
%		bar width={0.06\textwidth},
%		legend style={
%			at={(0.5,-0.2)},
%			anchor=north,
%			legend columns=-1},
%		ylabel={},
%		symbolic x coords={
%			Jan,
%			Feb,
%			Mar,
%			Apr,
%			May,
%			Jun,
%			Jul,
%			Aug,
%			Sep,
%			Oct,
%			Nov,
%			Dec,
%			},
%		xtick=data,
%		nodes near coords,
%		nodes near coords align={vertical},
%		x tick label style={},]
%	\addplot [draw=pfblue,
%		nodes near coords={\pgfmathfloatifflags{\pgfplotspointmeta}{0}{}{\pgfmathprintnumber{\pgfplotspointmeta}}},
%		nodes near coords align={horizontal},
%		nodes near coords style={font=\Large,/pgf/number format/assume math mode}]
%		coordinates{
%		(Jan,3)
%		(Feb,8)
%		(Mar,0)
%		(Apr,2)
%		(May,0)
%		(Jun,0)
%		(Jul,0)
%		(Aug,0)
%		(Sep,0)
%		(Oct,0)
%		(Nov,0)
%		(Dec,0)
%		};
%	\end{axis}
%\end{tikzpicture}
%
%
%\begin{tikzpicture}% Bar graph, use for current month's values
%	\begin{axis}
%		[ybar,
%		ymin=0,
%		title=\fontsize{25}{25}\selectfont\color{pfgrey}{Warranties Claims - \monthyeardate\today},
%		tick label style={/pgf/number format/assume math mode},
%		every axis plot/.append style={ultra thick},
%		ymajorgrids,
%		bar width={0.06\textwidth},
%		legend style={
%			at={(0.5,-0.2)},
%			anchor=north,
%			legend columns=-1},
%		ylabel={},
%		symbolic x coords={
%			OE ≥ 4000,
%			Thing 2,
%			Thing 3,
%			Thing 4,
%			Thing 5,
%			Thing 6,
%			},
%		xtick=data,
%		nodes near coords,
%		nodes near coords align={vertical},
%		x tick label style={},]
%	\addplot [fill=pfblue,
%		draw=none,
%		nodes near coords={
%			\pgfmathfloatifflags
%			{\pgfplotspointmeta}{0}{}
%			{\pgfmathprintnumber{\pgfplotspointmeta}}},
%		nodes near coords align={south},
%		nodes near coords style={font=\Large,/pgf/number format/assume math mode},
%		every node near coord/.append style={xshift=0pt,yshift=-24pt,anchor=south,font=\color{white}\Large}]
%		coordinates{
%		(OE ≥ 4000,3)
%		(Thing 2,8)
%		(Thing 3,2)
%		(Thing 4,1)
%		(Thing 5,0)
%		(Thing 6,0)
%		};
%	\end{axis}
%% Pareto line
%	\begin{axis}
%		[ymin=0,
%		ymax=100,
%		tick label style={/pgf/number format/assume math mode},
%		every axis plot/.append style={ultra thick},
%		ytick style={draw=none},
%		%bar width={0.06\textwidth},
%		legend style={
%			at={(0.5,-0.2)},
%			anchor=north,
%			legend columns=-1},
%		ylabel={},
%		xtick=data,
%		xticklabels={,,},
%		yticklabel={\pgfmathprintnumber\tick\%},
%		xtick style={draw=none},
%		%nodes near coords,
%		%nodes near coords align={vertical},
%		yticklabel pos=right,
%		y tick label style={},
%		x tick label style={},]
%	\addplot [draw=orange,
%		nodes near coords={}
%%			\pgfmathfloatifflags
%%			{\pgfplotspointmeta}{0}{}
%%			{\pgfmathprintnumber{\pgfplotspointmeta}}},
%%		nodes near coords align={south},
%%		nodes near coords style={font=\Large,/pgf/number format/assume math mode},
%%		every node near coord/.append style={xshift=0pt,yshift=0pt,anchor=south,font=\color{orange}\Large}
%]
%		coordinates{
%		(0,3)
%		(1,30)
%		(2,40)
%		(3,50)
%		(4,60)
%		(5,100)
%		};
%	\end{axis}
%\end{tikzpicture}
\end{document}
